\makeatletter
\newlabel{sup:acknowledgements}{{}{1}{\centering Acknowledgements}{section*.2}{}}
\newlabel{sup:acknowledgements@cref}{{}{[1][1][]1}}
\newlabel{oi:summary}{{}{2}{Contributions at a glance}{section*.4}{}}
\newlabel{oi:summary@cref}{{}{[1][2][]2}}
\newlabel{oi:summary-redesigned}{{}{3}{Contributions at a glance}{section*.6}{}}
\newlabel{oi:summary-redesigned@cref}{{}{[1][3][]3}}
\newlabel{sup:web-applications}{{}{4}{Supplementary web applications}{section*.8}{}}
\newlabel{sup:web-applications@cref}{{}{[1][4][]4}}
\newlabel{sup:fig:volvox-works}{{I}{4}{The portfolio builder landing page}{figure.0.1}{}}
\newlabel{sup:fig:volvox-works@cref}{{[figure][1][0]I}{[1][4][]4}}
\newlabel{sup:fig:volvox-thesis}{{II}{4}{The Volvox page collecting the full thesis content and its visualisers. (Links to individual applets can also be found in the relevant figure captions.)}{figure.0.2}{}}
\newlabel{sup:fig:volvox-thesis@cref}{{[figure][2][0]II}{[1][4][]4}}
\newlabel{sup:fig:markov}{{III}{4}{Markov Side-by-Side, where all of the Markov processes present in this thesis can be seen and interacted with}{figure.0.3}{}}
\newlabel{sup:fig:markov@cref}{{[figure][3][0]III}{[1][4][]4}}
\newlabel{ch:introduction}{{1}{11}{Introduction}{chapter.1}{}}
\newlabel{ch:introduction@cref}{{[chapter][1][]1}{[1][11][]11}}
\newlabel{fig:plague_art}{{1.1}{12}{Plague in the European imagination, in allegory and in reportage. (a)~\textit {The Triumph of Death}, Pieter Bruegel the Elder, c.~1562 (Museo del Prado, Madrid). An army of skeletons lays waste to a landscape in which every human occupation is interrupted at once; in the lower right the card game, the banquet and the lute player have not yet noticed. The painting is an allegory of death in general rather than a record of an epidemic, but it was made in a Europe that had lived with recurrent plague for two centuries. (b)~\textit {Piazza Mercatello a Napoli durante la peste del 1656}, Domenico Gargiulo, called Micco Spadaro (Museo Nazionale di San Martino, Naples). Here there is no allegory: a single square in a real city during a real outbreak, the dead laid out in rows and the living still at work among them. Naples is thought to have lost something close to half its population in that year}{figure.1.1}{}}
\newlabel{fig:plague_art@cref}{{[figure][1][1]1.1}{[1][11][]12}}
\newlabel{fig:yp_micrographs}{{1.2}{14}{The organism. (a) Scanning electron micrograph of \yp {} massed in the foregut of a flea, where the bacteria form the biofilm blockage that makes the vector infectious; the scale bar in the original reads $1.5\,\mu $m. (b) Wayson stain, showing the bipolar ``safety pin'' appearance by which the bacillus is recognised, scattered among host leucocytes. (c) Direct fluorescent antibody stain, in which the antibody binds the F1 capsular antigen --- one of the phase-two factors described below, and absent at flea temperature; the field gives a fair impression of the numbers involved. Sources: (a)~NIAID; (b)~and~(c) CDC Public Health Image Library}{figure.1.2}{}}
\newlabel{fig:yp_micrographs@cref}{{[figure][2][1]1.2}{[1][14][]14}}
\newlabel{fig:phagocytosis}{{1.3}{16}{Phagocytosis, the act on which the whole strategy turns. A neutrophil (yellow) engulfing bacilli (orange); the scale bar is $5\,\mu $m. This is what \yp {} permits in its first phase and refuses in its second, and it is what the injected effectors YopH, YopE, YopT and YpkA exist to prevent. The bacteria shown are \textit {Bacillus anthracis} rather than \yp {}. Image by Volker Brinkmann, cover of \textit {PLoS Pathogens} 1(3), 2005}{figure.1.3}{}}
\newlabel{fig:phagocytosis@cref}{{[figure][3][1]1.3}{[1][16][]16}}
\newlabel{fig:two_columns_six_image}{{1.4}{17}{A schematic representation}{figure.1.4}{}}
\newlabel{fig:two_columns_six_image@cref}{{[figure][4][1]1.4}{[1][16][]17}}
\newlabel{fig:hiv_budding}{{1.5}{20}{Budding release. Scanning electron micrograph of HIV-1 virions (green) assembling and departing from the surface of a cultured lymphocyte, each bump a separate event and the cell intact. Release and death are independent here, which is why the standard model's two constants describe such a cell exactly. A bursting pathogen does the opposite: nothing leaves while the load accumulates, and then everything leaves at once, in the event that ends the cell --- and no pair of constants describes that. \Mextref {p:ch:one-cell-to-population}{Chapter~6} turns on the difference. Image: CDC Public Health Image Library \#10000, photo credit C.~Goldsmith}{figure.1.5}{}}
\newlabel{fig:hiv_budding@cref}{{[figure][5][1]1.5}{[1][19][]20}}
\newlabel{m:ch:math-intro}{{2}{25}{Mathematical introduction}{chapter.2}{}}
\newlabel{m:ch:math-intro@cref}{{[chapter][2][]2}{[1][25][]25}}
\newlabel{m:sec:overview}{{2.1}{25}{Overview}{section.2.1}{}}
\newlabel{m:sec:overview@cref}{{[section][1][2]2.1}{[1][25][]25}}
\newlabel{m:sec:chains}{{2.2}{26}{Markov chains}{section.2.2}{}}
\newlabel{m:sec:chains@cref}{{[section][2][2]2.2}{[1][26][]26}}
\newlabel{m:sec:dtmc}{{2.2.1}{26}{Discrete-time Markov chains}{subsection.2.2.1}{}}
\newlabel{m:sec:dtmc@cref}{{[subsection][1][2,2]2.2.1}{[1][26][]26}}
\newlabel{m:eq:dtmarkov}{{2.1}{26}{Discrete-time Markov chains}{equation.2.2.1}{}}
\newlabel{m:eq:dtmarkov@cref}{{[equation][1][2]2.1}{[1][26][]26}}
\newlabel{m:eq:chapmankolmogorov}{{2.2}{26}{Discrete-time Markov chains}{equation.2.2.2}{}}
\newlabel{m:eq:chapmankolmogorov@cref}{{[equation][2][2]2.2}{[1][26][]26}}
\newlabel{m:sec:randomwalk}{{2.2.1.1}{26}{Simple random walks}{subsubsection.2.2.1.1}{}}
\newlabel{m:sec:randomwalk@cref}{{[subsubsection][1][2,2,1]2.2.1.1}{[1][26][]26}}
\newlabel{m:eq:rwrates}{{2.3}{26}{Simple random walks}{equation.2.2.3}{}}
\newlabel{m:eq:rwrates@cref}{{[equation][3][2]2.3}{[1][26][]26}}
\newlabel{m:eq:rwmoments}{{2.4}{26}{Simple random walks}{equation.2.2.4}{}}
\newlabel{m:eq:rwmoments@cref}{{[equation][4][2]2.4}{[1][26][]26}}
\newlabel{m:eq:rwlaw}{{2.5}{27}{Simple random walks}{equation.2.2.5}{}}
\newlabel{m:eq:rwlaw@cref}{{[equation][5][2]2.5}{[1][26][]27}}
\newlabel{m:fig:rwlaw}{{2.1}{27}{The displacement law \cref {m:eq:rwlaw} after $n=40$ steps, for $q=0.4,0.5,0.6$ and $X_0=0$. Dotted lines mark the means $n(2q-1)$. Changing $q$ mainly translates the distribution over this range, since its standard deviation $\sqrt {4nq(1-q)}$ changes comparatively little}{figure.2.1}{}}
\newlabel{m:fig:rwlaw@cref}{{[figure][1][2]2.1}{[1][27][]27}}
\newlabel{m:eq:embedded}{{2.6}{27}{Simple random walks}{equation.2.2.6}{}}
\newlabel{m:eq:embedded@cref}{{[equation][6][2]2.6}{[1][27][]27}}
\newlabel{m:eq:gamblersruin}{{2.7}{27}{Simple random walks}{equation.2.2.7}{}}
\newlabel{m:eq:gamblersruin@cref}{{[equation][7][2]2.7}{[1][27][]27}}
\newlabel{m:sec:gw}{{2.2.1.2}{27}{Galton--Watson branching processes}{subsubsection.2.2.1.2}{}}
\newlabel{m:sec:gw@cref}{{[subsubsection][2][2,2,1]2.2.1.2}{[1][27][]27}}
\newlabel{m:fig:randomwalk}{{2.2}{28}{Two views of the walk \cref {m:eq:rwrates}. (a) Seven paths from $X_0=0$ with $q=0.45$, together with the mean and one-standard-deviation envelope from \cref {m:eq:rwmoments}. (b) The probability \cref {m:eq:gamblersruin} of reaching $M=20$ before $0$. The symmetric case $\theta =1$ is linear; drift towards $M$ ($\theta <1$) bends the probability upwards, while drift towards $0$ bends it downwards. \explorerlink {https://markov.volvox.site/-/adam74aldridge/simple-random-walk}{Simulate}}{figure.2.2}{}}
\newlabel{m:fig:randomwalk@cref}{{[figure][2][2]2.2}{[1][27][]28}}
\newlabel{m:eq:gwrecursion}{{2.8}{28}{Galton--Watson branching processes}{equation.2.2.8}{}}
\newlabel{m:eq:gwrecursion@cref}{{[equation][8][2]2.8}{[1][28][]28}}
\newlabel{m:eq:gwcomposition}{{2.9}{28}{Galton--Watson branching processes}{equation.2.2.9}{}}
\newlabel{m:eq:gwcomposition@cref}{{[equation][9][2]2.9}{[1][28][]28}}
\newlabel{m:eq:binarylaw}{{2.10}{28}{The binary case}{equation.2.2.10}{}}
\newlabel{m:eq:binarylaw@cref}{{[equation][10][2]2.10}{[1][28][]28}}
\newlabel{m:eq:pgfbinary}{{2.11}{28}{The binary case}{equation.2.2.11}{}}
\newlabel{m:eq:pgfbinary@cref}{{[equation][11][2]2.11}{[1][28][]28}}
\newlabel{m:eq:SnIter}{{2.12}{28}{The binary case}{equation.2.2.12}{}}
\newlabel{m:eq:SnIter@cref}{{[equation][12][2]2.12}{[1][28][]28}}
\newlabel{m:sec:ctmc}{{2.2.2}{29}{Continuous-time Markov chains}{subsection.2.2.2}{}}
\newlabel{m:sec:ctmc@cref}{{[subsection][2][2,2]2.2.2}{[1][28][]29}}
\newlabel{m:eq:markov}{{2.13}{29}{Continuous-time Markov chains}{equation.2.2.13}{}}
\newlabel{m:eq:markov@cref}{{[equation][13][2]2.13}{[1][28][]29}}
\newlabel{m:eq:generator}{{2.14}{29}{Continuous-time Markov chains}{equation.2.2.14}{}}
\newlabel{m:eq:generator@cref}{{[equation][14][2]2.14}{[1][29][]29}}
\newlabel{m:eq:infinitesimal}{{2.15}{29}{Continuous-time Markov chains}{equation.2.2.15}{}}
\newlabel{m:eq:infinitesimal@cref}{{[equation][15][2]2.15}{[1][29][]29}}
\newlabel{m:sec:exp}{{2.2.2.1}{29}{Exponential clocks}{subsubsection.2.2.2.1}{}}
\newlabel{m:sec:exp@cref}{{[subsubsection][1][2,2,2]2.2.2.1}{[1][29][]29}}
\newlabel{m:eq:expdef}{{2.16}{29}{Exponential clocks}{equation.2.2.16}{}}
\newlabel{m:eq:expdef@cref}{{[equation][16][2]2.16}{[1][29][]29}}
\newlabel{m:eq:memoryless}{{2.17}{29}{Exponential clocks}{equation.2.2.17}{}}
\newlabel{m:eq:memoryless@cref}{{[equation][17][2]2.17}{[1][29][]29}}
\newlabel{m:prop:race}{{2.1}{29}{Competing clocks}{theorem.2.1}{}}
\newlabel{m:prop:race@cref}{{[proposition][1][2]2.1}{[1][29][]29}}
\newlabel{m:eq:race}{{2.18}{29}{Competing clocks}{equation.2.2.18}{}}
\newlabel{m:eq:race@cref}{{[equation][18][2]2.18}{[1][29][]29}}
\newlabel{m:rem:gillespie}{{2.2}{29}{Direct simulation}{theorem.2.2}{}}
\newlabel{m:rem:gillespie@cref}{{[remark][2][2]2.2}{[1][29][]29}}
\newlabel{m:sec:poisson}{{2.2.2.2}{30}{The Poisson process}{subsubsection.2.2.2.2}{}}
\newlabel{m:sec:poisson@cref}{{[subsubsection][2][2,2,2]2.2.2.2}{[1][29][]30}}
\newlabel{m:eq:poissonrates}{{2.19}{30}{The Poisson process}{equation.2.2.19}{}}
\newlabel{m:eq:poissonrates@cref}{{[equation][19][2]2.19}{[1][29][]30}}
\newlabel{m:eq:poissonlaw}{{2.20}{30}{The Poisson process}{equation.2.2.20}{}}
\newlabel{m:eq:poissonlaw@cref}{{[equation][20][2]2.20}{[1][30][]30}}
\newlabel{m:fig:poisson}{{2.3}{30}{A Poisson process of rate $\alpha =1.5$. (a) One path, with its exponential holding intervals and mean $\alpha t$. (b) The law \cref {m:eq:poissonlaw} at $t=1,3,6$. Its absolute width grows with time, while its coefficient of variation $1/\sqrt {\alpha t}$ decreases. \explorerlink {https://markov.volvox.site/-/adam74aldridge/poisson-process}{Simulate}}{figure.2.3}{}}
\newlabel{m:fig:poisson@cref}{{[figure][3][2]2.3}{[1][30][]30}}
\newlabel{m:sec:yule}{{2.2.2.3}{30}{The Yule process}{subsubsection.2.2.2.3}{}}
\newlabel{m:sec:yule@cref}{{[subsubsection][3][2,2,2]2.2.2.3}{[1][30][]30}}
\newlabel{m:eq:yulerates}{{2.21}{30}{The Yule process}{equation.2.2.21}{}}
\newlabel{m:eq:yulerates@cref}{{[equation][21][2]2.21}{[1][30][]30}}
\newlabel{m:eq:yulelaw}{{2.22}{30}{The Yule process}{equation.2.2.22}{}}
\newlabel{m:eq:yulelaw@cref}{{[equation][22][2]2.22}{[1][30][]30}}
\newlabel{m:rem:yuleorigin}{{2.3}{31}{Origin and purpose}{theorem.2.3}{}}
\newlabel{m:rem:yuleorigin@cref}{{[remark][3][2]2.3}{[1][30][]31}}
\newlabel{m:sec:bd}{{2.2.2.4}{31}{Birth--death processes}{subsubsection.2.2.2.4}{}}
\newlabel{m:sec:bd@cref}{{[subsubsection][4][2,2,2]2.2.2.4}{[1][31][]31}}
\newlabel{m:eq:bdrates}{{2.23}{31}{Birth--death processes}{equation.2.2.23}{}}
\newlabel{m:eq:bdrates@cref}{{[equation][23][2]2.23}{[1][31][]31}}
\newlabel{m:fig:bdpaths}{{2.4}{31}{Nine paths of the linear birth--death process from $X_0=5$, with $\mu =1$ and three values of $\lambda $; red curves show $5\ee ^{(\lambda -\mu )t}$. At criticality the mean remains five although some paths are already extinct and others are much larger. Supercritical growth also permits early extinction, whose eventual probability here is $(\mu /\lambda )^5$. \explorerlink {https://markov.volvox.site/-/adam74aldridge/linear-birth-death-process}{Simulate}}{figure.2.4}{}}
\newlabel{m:fig:bdpaths@cref}{{[figure][4][2]2.4}{[1][31][]31}}
\newlabel{m:sec:bdc}{{2.2.2.5}{32}{Birth--death--catastrophe processes}{subsubsection.2.2.2.5}{}}
\newlabel{m:sec:bdc@cref}{{[subsubsection][5][2,2,2]2.2.2.5}{[1][31][]32}}
\newlabel{m:def:bdc}{{2.4}{32}{Birth--death--catastrophe process}{theorem.2.4}{}}
\newlabel{m:def:bdc@cref}{{[definition][4][2]2.4}{[1][32][]32}}
\newlabel{m:eq:bdcrates}{{2.24}{32}{Birth--death--catastrophe process}{equation.2.2.24}{}}
\newlabel{m:eq:bdcrates@cref}{{[equation][24][2]2.24}{[1][32][]32}}
\newlabel{m:eq:bdcrelease}{{2.25}{32}{Birth--death--catastrophe processes}{equation.2.2.25}{}}
\newlabel{m:eq:bdcrelease@cref}{{[equation][25][2]2.25}{[1][32][]32}}
\newlabel{m:rem:notmeanfield}{{2.5}{32}{A state-dependent catastrophe is not a mean-field hazard}{theorem.2.5}{}}
\newlabel{m:rem:notmeanfield@cref}{{[remark][5][2]2.5}{[1][32][]32}}
\newlabel{m:eq:norupture}{{2.26}{32}{A state-dependent catastrophe is not a mean-field hazard}{equation.2.2.26}{}}
\newlabel{m:eq:norupture@cref}{{[equation][26][2]2.26}{[1][32][]32}}
\newlabel{m:fig:ratediagram}{{2.5}{33}{Linear birth--death dynamics (top) and the birth--death--catastrophe extension (bottom). Double circles are absorbing. Ordinary extinction enters $0$ only from state $1$; catastrophe instead jumps from any positive count to the distinct state $H$, at total rate $\delta n$}{figure.2.5}{}}
\newlabel{m:fig:ratediagram@cref}{{[figure][5][2]2.5}{[1][32][]33}}
\newlabel{m:sec:inhomogeneous}{{2.2.3}{33}{Time-inhomogeneous processes}{subsection.2.2.3}{}}
\newlabel{m:sec:inhomogeneous@cref}{{[subsection][3][2,2]2.2.3}{[1][33][]33}}
\newlabel{m:eq:inhomogtrans}{{2.27}{33}{Time-inhomogeneous processes}{equation.2.2.27}{}}
\newlabel{m:eq:inhomogtrans@cref}{{[equation][27][2]2.27}{[1][33][]33}}
\newlabel{m:eq:inhomogkolmogorov}{{2.28}{33}{Time-inhomogeneous processes}{equation.2.2.28}{}}
\newlabel{m:eq:inhomogkolmogorov@cref}{{[equation][28][2]2.28}{[1][33][]33}}
\newlabel{m:sec:tibd}{{2.2.3}{33}{Linear birth--death rates}{section*.32}{}}
\newlabel{m:sec:tibd@cref}{{[subsection][3][2,2]2.2.3}{[1][33][]33}}
\newlabel{m:eq:meanodeinhom}{{2.29}{33}{Linear birth--death rates}{equation.2.2.29}{}}
\newlabel{m:eq:meanodeinhom@cref}{{[equation][29][2]2.29}{[1][33][]33}}
\newlabel{m:eq:riccatiext}{{2.30}{34}{Linear birth--death rates}{equation.2.2.30}{}}
\newlabel{m:eq:riccatiext@cref}{{[equation][30][2]2.30}{[1][33][]34}}
\newlabel{m:eq:inhomextforward}{{2.31}{34}{Linear birth--death rates}{equation.2.2.31}{}}
\newlabel{m:eq:inhomextforward@cref}{{[equation][31][2]2.31}{[1][34][]34}}
\newlabel{m:sec:seasonallv}{{2.2.3.1}{34}{The seasonal Lotka--Volterra process}{subsubsection.2.2.3.1}{}}
\newlabel{m:sec:seasonallv@cref}{{[subsubsection][1][2,2,3]2.2.3.1}{[1][34][]34}}
\newlabel{m:eq:lvrate}{{2.32}{34}{The seasonal Lotka--Volterra process}{equation.2.2.32}{}}
\newlabel{m:eq:lvrate@cref}{{[equation][32][2]2.32}{[1][34][]34}}
\newlabel{m:eq:lvkernel}{{2.33}{34}{The seasonal Lotka--Volterra process}{equation.2.2.33}{}}
\newlabel{m:eq:lvkernel@cref}{{[equation][33][2]2.33}{[1][34][]34}}
\newlabel{m:eq:lvmeanfield}{{2.34}{34}{The seasonal Lotka--Volterra process}{equation.2.2.34}{}}
\newlabel{m:eq:lvmeanfield@cref}{{[equation][34][2]2.34}{[1][34][]34}}
\newlabel{m:fig:seasonallv}{{2.6}{35}{The seasonal Lotka--Volterra process with $A=2$, $\omega =2\pi $, $b=2$, $e=0.005$, $d=2$, and $(X_0,Y_0)=(300,100)$. (a) Clipping the forcing creates a closed season occupying one third of every period. (b) Ten exact prey--predator paths. (c) Pointwise medians and 10--90\% bands from 200 paths, with \cref {m:eq:lvmeanfield} dashed. The first cycle is calendar-locked and tightly concentrated; later interaction amplifies phase and amplitude differences. If prey reach zero, the remaining predators form a pure-death process and the path eventually reaches $(0,0)$}{figure.2.6}{}}
\newlabel{m:fig:seasonallv@cref}{{[figure][6][2]2.6}{[1][34][]35}}
\newlabel{m:sec:logisticspeciation}{{2.2.3.2}{35}{The logistic speciation model}{subsubsection.2.2.3.2}{}}
\newlabel{m:sec:logisticspeciation@cref}{{[subsubsection][2][2,2,3]2.2.3.2}{[1][35][]35}}
\newlabel{m:eq:logisticode}{{2.35}{35}{The logistic speciation model}{equation.2.2.35}{}}
\newlabel{m:eq:logisticode@cref}{{[equation][35][2]2.35}{[1][35][]35}}
\newlabel{m:eq:logisticsol}{{2.36}{35}{The logistic speciation model}{equation.2.2.36}{}}
\newlabel{m:eq:logisticsol@cref}{{[equation][36][2]2.36}{[1][35][]35}}
\newlabel{m:eq:speciationrate}{{2.37}{35}{The logistic speciation model}{equation.2.2.37}{}}
\newlabel{m:eq:speciationrate@cref}{{[equation][37][2]2.37}{[1][35][]35}}
\newlabel{m:eq:lambdat}{{2.38}{36}{The logistic speciation model}{equation.2.2.38}{}}
\newlabel{m:eq:lambdat@cref}{{[equation][38][2]2.38}{[1][35][]36}}
\newlabel{m:eq:lambdaint}{{2.39}{36}{The logistic speciation model}{equation.2.2.39}{}}
\newlabel{m:eq:lambdaint@cref}{{[equation][39][2]2.39}{[1][36][]36}}
\newlabel{m:eq:logspecmean}{{2.40}{36}{The logistic speciation model}{equation.2.2.40}{}}
\newlabel{m:eq:logspecmean@cref}{{[equation][40][2]2.40}{[1][36][]36}}
\newlabel{m:fig:logspecmean}{{2.7}{36}{Mean species count \cref {m:eq:logspecmean} for $N_0=1,2,5$, with $C=10$, $r=1$, and $\sigma =0.08$; dashed lines are the limits $(C/N_0)^{\sigma C/r}$. A smaller founding abundance leaves more integrated unused capacity, so the limiting mean is higher. \explorerlink {https://markov.volvox.site/-/adam74aldridge/logistic-environment-type-count-process}{Simulate}}{figure.2.7}{}}
\newlabel{m:fig:logspecmean@cref}{{[figure][7][2]2.7}{[1][36][]36}}
\newlabel{m:fig:lsmean}{{2.8}{37}{Numerical comparison with a logistic-environment birth--death variant. Pale blue curves are species-count paths and the dark blue curve is their sample mean; the lower panel shows the per-species origination rate $\sigma N'(t)$ and a constant extinction rate. Richness turns over after origination falls below extinction. The ordinate $X_t$ in the supplied illustration denotes the species count $S_t$ used in this section; this is not the pure-birth model of \cref {m:eq:speciationrate,m:eq:logspecmean}}{figure.2.8}{}}
\newlabel{m:fig:lsmean@cref}{{[figure][8][2]2.8}{[1][36][]37}}
\newlabel{m:sec:coupled}{{2.2.4}{37}{Coupled ODE--CTMC systems}{subsection.2.2.4}{}}
\newlabel{m:sec:coupled@cref}{{[subsection][4][2,2]2.2.4}{[1][37][]37}}
\newlabel{m:eq:coupledone}{{2.41}{37}{Coupled ODE--CTMC systems}{equation.2.2.41}{}}
\newlabel{m:eq:coupledone@cref}{{[equation][41][2]2.41}{[1][37][]37}}
\newlabel{m:eq:coupledtwo}{{2.42}{37}{Coupled ODE--CTMC systems}{equation.2.2.42}{}}
\newlabel{m:eq:coupledtwo@cref}{{[equation][42][2]2.42}{[1][37][]37}}
\newlabel{m:eq:coupledthree}{{2.43}{37}{Coupled ODE--CTMC systems}{equation.2.2.43}{}}
\newlabel{m:eq:coupledthree@cref}{{[equation][43][2]2.43}{[1][37][]37}}
\newlabel{m:eq:coupledgenerator}{{2.44}{38}{Coupled ODE--CTMC systems}{equation.2.2.44}{}}
\newlabel{m:eq:coupledgenerator@cref}{{[equation][44][2]2.44}{[1][37][]38}}
\newlabel{m:sec:gated}{{2.2.4.1}{38}{Accumulation with gated release}{subsubsection.2.2.4.1}{}}
\newlabel{m:sec:gated@cref}{{[subsubsection][1][2,2,4]2.2.4.1}{[1][38][]38}}
\newlabel{m:eq:gatedflow}{{2.45}{38}{Accumulation with gated release}{equation.2.2.45}{}}
\newlabel{m:eq:gatedflow@cref}{{[equation][45][2]2.45}{[1][38][]38}}
\newlabel{m:eq:gatedjumps}{{2.46}{38}{Accumulation with gated release}{equation.2.2.46}{}}
\newlabel{m:eq:gatedjumps@cref}{{[equation][46][2]2.46}{[1][38][]38}}
\newlabel{m:eq:gatedbalance}{{2.47}{38}{Accumulation with gated release}{equation.2.2.47}{}}
\newlabel{m:eq:gatedbalance@cref}{{[equation][47][2]2.47}{[1][38][]38}}
\newlabel{m:fig:gatedpaths}{{2.9}{39}{Three exact paths of gated release from $(A_0,R_0,G_0)=(5,0,0)$, with $\rho =2$, $c=4$, $\eta =0.015$, and $\gamma =1.45$. While the gate is closed, $A_t$ grows linearly and $R_t$ is constant; an open gate transfers the accumulated amount into $R_t$ and pulls $A_t$ towards $\rho /c$. State-dependent opening produces intermittent releases}{figure.2.9}{}}
\newlabel{m:fig:gatedpaths@cref}{{[figure][9][2]2.9}{[1][38][]39}}
\newlabel{m:sec:surging}{{2.2.4.2}{39}{Accumulation with surging gated release}{subsubsection.2.2.4.2}{}}
\newlabel{m:sec:surging@cref}{{[subsubsection][2][2,2,4]2.2.4.2}{[1][39][]39}}
\newlabel{m:eq:surgingjumps}{{2.48}{39}{Accumulation with surging gated release}{equation.2.2.48}{}}
\newlabel{m:eq:surgingjumps@cref}{{[equation][48][2]2.48}{[1][39][]39}}
\newlabel{m:eq:gatemean}{{2.49}{39}{Accumulation with surging gated release}{equation.2.2.49}{}}
\newlabel{m:eq:gatemean@cref}{{[equation][49][2]2.49}{[1][39][]39}}
\newlabel{m:sec:coupledsim}{{2.2.4.3}{39}{Simulating a coupled system}{subsubsection.2.2.4.3}{}}
\newlabel{m:sec:coupledsim@cref}{{[subsubsection][3][2,2,4]2.2.4.3}{[1][39][]39}}
\newlabel{m:eq:coupledholding}{{2.50}{39}{Simulating a coupled system}{equation.2.2.50}{}}
\newlabel{m:eq:coupledholding@cref}{{[equation][50][2]2.50}{[1][39][]39}}
\newlabel{m:fig:gatedensemble}{{2.10}{40}{Medians and pointwise 10--90\% bands from 2000 paths. Both models use $\rho =2$, $c=4$, $\gamma =1.45$, and $(A_0,R_0,G_0)=(5,0,0)$. Gated release (left) opens at intensity $\eta A_t$ with $\eta =0.015$; surging release (right) opens at the autonomous rate $\nu =0.3$. The parameters have different dimensions, and the two opening mechanisms produce visibly different accumulation levels}{figure.2.10}{}}
\newlabel{m:fig:gatedensemble@cref}{{[figure][10][2]2.10}{[1][39][]40}}
\newlabel{m:sec:methods}{{2.3}{40}{Methods for Markov chains}{section.2.3}{}}
\newlabel{m:sec:methods@cref}{{[section][3][2]2.3}{[1][40][]40}}
\newlabel{m:sec:dtmethods}{{2.3.1}{40}{Discrete-time methods}{subsection.2.3.1}{}}
\newlabel{m:sec:dtmethods@cref}{{[subsection][1][2,3]2.3.1}{[1][40][]40}}
\newlabel{m:sec:firststep}{{2.3.1.1}{40}{First-step analysis}{subsubsection.2.3.1.1}{}}
\newlabel{m:sec:firststep@cref}{{[subsubsection][1][2,3,1]2.3.1.1}{[1][40][]40}}
\newlabel{m:prop:firststep}{{2.6}{40}{First-step principle}{theorem.2.6}{}}
\newlabel{m:prop:firststep@cref}{{[proposition][6][2]2.6}{[1][40][]40}}
\newlabel{m:eq:firststepdiscrete}{{2.51}{41}{First-step principle}{equation.2.3.51}{}}
\newlabel{m:eq:firststepdiscrete@cref}{{[equation][51][2]2.51}{[1][41][]41}}
\newlabel{m:eq:firststep}{{2.52}{41}{First-step principle}{equation.2.3.52}{}}
\newlabel{m:eq:firststep@cref}{{[equation][52][2]2.52}{[1][41][]41}}
\newlabel{m:eq:extinctionfixedpoint}{{2.53}{41}{First-step analysis}{equation.2.3.53}{}}
\newlabel{m:eq:extinctionfixedpoint@cref}{{[equation][53][2]2.53}{[1][41][]41}}
\newlabel{m:sec:pgfmethod}{{2.3.1.2}{41}{Generating functions}{subsubsection.2.3.1.2}{}}
\newlabel{m:sec:pgfmethod@cref}{{[subsubsection][2][2,3,1]2.3.1.2}{[1][41][]41}}
\newlabel{m:eq:pgfdef}{{2.54}{41}{Generating functions}{equation.2.3.54}{}}
\newlabel{m:eq:pgfdef@cref}{{[equation][54][2]2.54}{[1][41][]41}}
\newlabel{m:eq:pgfmoments}{{2.55}{41}{Generating functions}{equation.2.3.55}{}}
\newlabel{m:eq:pgfmoments@cref}{{[equation][55][2]2.55}{[1][41][]41}}
\newlabel{m:eq:pgfbivariate}{{2.56}{41}{Generating functions}{equation.2.3.56}{}}
\newlabel{m:eq:pgfbivariate@cref}{{[equation][56][2]2.56}{[1][41][]41}}
\newlabel{m:sec:iteration}{{2.3.1.3}{42}{Functional iteration}{subsubsection.2.3.1.3}{}}
\newlabel{m:sec:iteration@cref}{{[subsubsection][3][2,3,1]2.3.1.3}{[1][41][]42}}
\newlabel{m:sec:ctmethods}{{2.3.2}{42}{Continuous-time methods}{subsection.2.3.2}{}}
\newlabel{m:sec:ctmethods@cref}{{[subsection][2][2,3]2.3.2}{[1][42][]42}}
\newlabel{m:sec:backward}{{2.3.2.1}{42}{Forward and backward equations}{subsubsection.2.3.2.1}{}}
\newlabel{m:sec:backward@cref}{{[subsubsection][1][2,3,2]2.3.2.1}{[1][42][]42}}
\newlabel{m:eq:forward}{{2.57}{42}{Forward and backward equations}{equation.2.3.57}{}}
\newlabel{m:eq:forward@cref}{{[equation][57][2]2.57}{[1][42][]42}}
\newlabel{m:eq:backwardeq}{{2.58}{42}{Forward and backward equations}{equation.2.3.58}{}}
\newlabel{m:eq:backwardeq@cref}{{[equation][58][2]2.58}{[1][42][]42}}
\newlabel{m:eq:master}{{2.59}{42}{Forward and backward equations}{equation.2.3.59}{}}
\newlabel{m:eq:master@cref}{{[equation][59][2]2.59}{[1][42][]42}}
\newlabel{m:eq:bdmaster}{{2.60}{42}{Forward and backward equations}{equation.2.3.60}{}}
\newlabel{m:eq:bdmaster@cref}{{[equation][60][2]2.60}{[1][42][]42}}
\newlabel{m:eq:bdpde}{{2.61}{42}{Forward and backward equations}{equation.2.3.61}{}}
\newlabel{m:eq:bdpde@cref}{{[equation][61][2]2.61}{[1][42][]42}}
\newlabel{m:eq:bdsolution}{{2.62}{43}{Forward and backward equations}{equation.2.3.62}{}}
\newlabel{m:eq:bdsolution@cref}{{[equation][62][2]2.62}{[1][42][]43}}
\newlabel{m:sec:mean}{{2.3.2.2}{43}{Mean population size}{subsubsection.2.3.2.2}{}}
\newlabel{m:sec:mean@cref}{{[subsubsection][2][2,3,2]2.3.2.2}{[1][43][]43}}
\newlabel{m:eq:meanode}{{2.63}{43}{Mean population size}{equation.2.3.63}{}}
\newlabel{m:eq:meanode@cref}{{[equation][63][2]2.63}{[1][43][]43}}
\newlabel{m:eq:bdmean}{{2.64}{43}{Mean population size}{equation.2.3.64}{}}
\newlabel{m:eq:bdmean@cref}{{[equation][64][2]2.64}{[1][43][]43}}
\newlabel{m:sec:variance}{{2.3.2.3}{43}{Variance and higher moments}{subsubsection.2.3.2.3}{}}
\newlabel{m:sec:variance@cref}{{[subsubsection][3][2,3,2]2.3.2.3}{[1][43][]43}}
\newlabel{m:eq:secondmomentode}{{2.65}{43}{Variance and higher moments}{equation.2.3.65}{}}
\newlabel{m:eq:secondmomentode@cref}{{[equation][65][2]2.65}{[1][43][]43}}
\newlabel{m:eq:bdvariance}{{2.66}{43}{Variance and higher moments}{equation.2.3.66}{}}
\newlabel{m:eq:bdvariance@cref}{{[equation][66][2]2.66}{[1][43][]43}}
\newlabel{m:sec:hitting}{{2.3.2.4}{43}{Hitting probabilities}{subsubsection.2.3.2.4}{}}
\newlabel{m:sec:hitting@cref}{{[subsubsection][4][2,3,2]2.3.2.4}{[1][43][]43}}
\newlabel{m:eq:bdjumpratios}{{2.67}{43}{Hitting probabilities}{equation.2.3.67}{}}
\newlabel{m:eq:bdjumpratios@cref}{{[equation][67][2]2.67}{[1][43][]43}}
\newlabel{m:eq:bdruin}{{2.68}{43}{Hitting probabilities}{equation.2.3.68}{}}
\newlabel{m:eq:bdruin@cref}{{[equation][68][2]2.68}{[1][43][]43}}
\newlabel{m:sec:extinctionprob}{{2.3.2.5}{44}{Extinction probabilities}{subsubsection.2.3.2.5}{}}
\newlabel{m:sec:extinctionprob@cref}{{[subsubsection][5][2,3,2]2.3.2.5}{[1][44][]44}}
\newlabel{m:eq:p0t}{{2.69}{44}{Extinction probabilities}{equation.2.3.69}{}}
\newlabel{m:eq:p0t@cref}{{[equation][69][2]2.69}{[1][44][]44}}
\newlabel{m:fig:extinction}{{2.11}{44}{(a) The extinction probability \cref {m:eq:p0t} against time, with $\mu =1$ and $\lambda =0.6,\,1,\,1.4$; solid curves start from $N=1$, dashed from $N=5$. Below and at criticality every curve climbs to one, however large the founding cohort; above it the curves level off at $(\mu /\lambda )^N$, marked dotted for $N=1$, and a larger cohort buys a permanently smaller extinction probability rather than merely a delay. (b) The critical case in detail, from one founder: the law of $X_t$ on the positive states at $t=1,3,10$, on a logarithmic scale. Absorbing mass accumulates at zero --- $p_0$ rises from $0.50$ to $0.91$ --- while the surviving tail flattens and stretches. The mean is exactly $1$ at all three times}{figure.2.11}{}}
\newlabel{m:fig:extinction@cref}{{[figure][11][2]2.11}{[1][44][]44}}
\newlabel{m:eq:ctsS}{{2.70}{44}{Extinction probabilities}{equation.2.3.70}{}}
\newlabel{m:eq:ctsS@cref}{{[equation][70][2]2.70}{[1][44][]44}}
\newlabel{m:sec:condmeans}{{2.3.2.6}{45}{Conditional means}{subsubsection.2.3.2.6}{}}
\newlabel{m:sec:condmeans@cref}{{[subsubsection][6][2,3,2]2.3.2.6}{[1][44][]45}}
\newlabel{m:def:qsd}{{2.7}{45}{Quasi-stationary and Yaglom laws}{theorem.2.7}{}}
\newlabel{m:def:qsd@cref}{{[definition][7][2]2.7}{[1][45][]45}}
\newlabel{m:eq:condit}{{2.71}{45}{The continuous-time constant}{equation.2.3.71}{}}
\newlabel{m:eq:condit@cref}{{[equation][71][2]2.71}{[1][45][]45}}
\newlabel{m:eq:ctscondmean}{{2.72}{45}{The continuous-time constant}{equation.2.3.72}{}}
\newlabel{m:eq:ctscondmean@cref}{{[equation][72][2]2.72}{[1][45][]45}}
\newlabel{m:eq:ctsqsmean}{{2.73}{45}{The continuous-time constant}{equation.2.3.73}{}}
\newlabel{m:eq:ctsqsmean@cref}{{[equation][73][2]2.73}{[1][45][]45}}
\newlabel{m:eq:Acts}{{2.74}{45}{The continuous-time constant}{equation.2.3.74}{}}
\newlabel{m:eq:Acts@cref}{{[equation][74][2]2.74}{[1][45][]45}}
\newlabel{m:eq:Apdef}{{2.75}{45}{The discrete-time constant}{equation.2.3.75}{}}
\newlabel{m:eq:Apdef@cref}{{[equation][75][2]2.75}{[1][45][]45}}
\newlabel{m:eq:qsmean}{{2.76}{46}{The discrete-time constant}{equation.2.3.76}{}}
\newlabel{m:eq:qsmean@cref}{{[equation][76][2]2.76}{[1][45][]46}}
\newlabel{m:eq:twoconstants}{{2.77}{46}{The two constants compared, and the handover}{equation.2.3.77}{}}
\newlabel{m:eq:twoconstants@cref}{{[equation][77][2]2.77}{[1][46][]46}}
\newlabel{m:fig:dtct}{{2.12}{46}{The continuous-time constant $\Ac (p)=(1-2p)/(1-p)$ (green) against high-precision numerical values of the discrete-time constant $A(p)$ (red points). Both vanish at $p=\tfrac 12$, but at $p=0$ the continuous curve starts at $1$ and the discrete one at $\tfrac 12$, and the two are not related by a constant factor in between. The black line is a formula obtained by symbolic regression; it tracks the points closely but is not exact, and \ChA {} explains why no elementary formula could be}{figure.2.12}{}}
\newlabel{m:fig:dtct@cref}{{[figure][12][2]2.12}{[1][46][]46}}
\newlabel{m:sec:moc}{{2.3.3}{47}{The method of characteristics}{subsection.2.3.3}{}}
\newlabel{m:sec:moc@cref}{{[subsection][3][2,3]2.3.3}{[1][46][]47}}
\newlabel{m:sec:mocidea}{{2.3.3.1}{47}{The idea}{subsubsection.2.3.3.1}{}}
\newlabel{m:sec:mocidea@cref}{{[subsubsection][1][2,3,3]2.3.3.1}{[1][47][]47}}
\newlabel{m:eq:genericfirstorder}{{2.78}{47}{The idea}{equation.2.3.78}{}}
\newlabel{m:eq:genericfirstorder@cref}{{[equation][78][2]2.78}{[1][47][]47}}
\newlabel{m:eq:chainrule}{{2.79}{47}{The idea}{equation.2.3.79}{}}
\newlabel{m:eq:chainrule@cref}{{[equation][79][2]2.79}{[1][47][]47}}
\newlabel{m:eq:genericcharacteristics}{{2.80}{47}{The idea}{equation.2.3.80}{}}
\newlabel{m:eq:genericcharacteristics@cref}{{[equation][80][2]2.80}{[1][47][]47}}
\newlabel{m:eq:ICpar}{{2.81}{47}{The idea}{equation.2.3.81}{}}
\newlabel{m:eq:ICpar@cref}{{[equation][81][2]2.81}{[1][47][]47}}
\newlabel{m:sec:mocworked}{{2.3.3.2}{47}{A worked example: absorption--death}{subsubsection.2.3.3.2}{}}
\newlabel{m:sec:mocworked@cref}{{[subsubsection][2][2,3,3]2.3.3.2}{[1][47][]47}}
\newlabel{m:fig:characteristics}{{2.13}{48}{Characteristics of the $x$-equation \cref {m:eq:ch13}, $\dd x/\dd \tau =\mu (1-x)$ with $\tau =-t$, drawn for $\mu =1$. Each grey curve is one characteristic, $x(t)=1-(1-\sigma )\ee ^{\mu t}$, labelled by the point $\sigma $ at which it meets the initial surface $t=0$, drawn as the heavy vertical line; together they foliate the domain, and exactly one passes through any given $(x,t)$. Since $G$ is constant along characteristics, its value at a point is its value at the foot of that point's curve --- so solving the equation amounts to running the curve backwards from $(x,t)$ to $t=0$ and reading off $\sigma $, here $\sigma =1-(1-x)\ee ^{-\mu t}$. The full problem is three-variate; this is its projection onto the $(x,t)$ plane, which is where the cascade starts, because \cref {m:eq:ch13} does not involve $y$}{figure.2.13}{}}
\newlabel{m:fig:characteristics@cref}{{[figure][13][2]2.13}{[1][47][]48}}
\newlabel{m:fig:compartment}{{2.14}{48}{The single-compartment absorption models. Particles begin in the exterior medium and are taken up one at a time, each independently at rate $\alpha $, so the absorption channel carries total rate $Y_t\alpha $; once inside, a particle dies at rate $\mu $ and, in the fullest version of the model, divides at rate $\lambda $. Only the three transitions drawn here occur, and each is proportional to the count in the compartment it acts on. Note that nothing returns from the compartment to the medium, which is why the exterior count is unaffected by the interior dynamics}{figure.2.14}{}}
\newlabel{m:fig:compartment@cref}{{[figure][14][2]2.14}{[1][47][]48}}
\newlabel{m:eq:FK2}{{2.82}{48}{A worked example: absorption--death}{equation.2.3.82}{}}
\newlabel{m:eq:FK2@cref}{{[equation][82][2]2.82}{[1][47][]48}}
\newlabel{m:eq:PDE2}{{2.83}{48}{A worked example: absorption--death}{equation.2.3.83}{}}
\newlabel{m:eq:PDE2@cref}{{[equation][83][2]2.83}{[1][48][]48}}
\newlabel{m:eq:absdeathstates}{{2.84}{49}{A worked example: absorption--death}{equation.2.3.84}{}}
\newlabel{m:eq:absdeathstates@cref}{{[equation][84][2]2.84}{[1][49][]49}}
\newlabel{m:eq:Gabsdeath}{{2.85}{49}{A worked example: absorption--death}{equation.2.3.85}{}}
\newlabel{m:eq:Gabsdeath@cref}{{[equation][85][2]2.85}{[1][49][]49}}
\newlabel{m:fig:abs2}{{2.15}{49}{State probabilities for the single-particle absorption--death model, \cref {m:eq:absdeathstates}. The particle starts outside ($p_{0,1}$, blue), is absorbed into the compartment ($p_{1,0}$, red, which rises and then decays as death takes over), and finally dies ($p_{0,0}$, green, absorbing). \explorerlink {https://markov.volvox.site/-/adam74aldridge/absorption-death-model}{Simulate}}{figure.2.15}{}}
\newlabel{m:fig:abs2@cref}{{[figure][15][2]2.15}{[1][49][]49}}
\newlabel{m:eq:ch11}{{2.86}{49}{A worked example: absorption--death}{equation.2.3.86}{}}
\newlabel{m:eq:ch11@cref}{{[equation][86][2]2.86}{[1][49][]49}}
\newlabel{m:eq:ch12}{{2.87}{49}{A worked example: absorption--death}{equation.2.3.87}{}}
\newlabel{m:eq:ch12@cref}{{[equation][87][2]2.87}{[1][49][]49}}
\newlabel{m:eq:ch13}{{2.88}{49}{A worked example: absorption--death}{equation.2.3.88}{}}
\newlabel{m:eq:ch13@cref}{{[equation][88][2]2.88}{[1][49][]49}}
\newlabel{m:eq:ch14}{{2.89}{49}{A worked example: absorption--death}{equation.2.3.89}{}}
\newlabel{m:eq:ch14@cref}{{[equation][89][2]2.89}{[1][49][]49}}
\newlabel{m:eq:GetaN}{{2.90}{49}{The equation for $G$, \eqref {m:eq:ch11}}{equation.2.3.90}{}}
\newlabel{m:eq:GetaN@cref}{{[equation][90][2]2.90}{[1][49][]49}}
\newlabel{m:eq:minustau}{{2.91}{49}{The equation for $t$, \eqref {m:eq:ch12}}{equation.2.3.91}{}}
\newlabel{m:eq:minustau@cref}{{[equation][91][2]2.91}{[1][49][]49}}
\newlabel{m:eq:sandxwithtau}{{2.92}{49}{The equation for $x$, \eqref {m:eq:ch13}}{equation.2.3.92}{}}
\newlabel{m:eq:sandxwithtau@cref}{{[equation][92][2]2.92}{[1][49][]49}}
\newlabel{m:eq:ywitheta}{{2.93}{50}{The equation for $y$, \eqref {m:eq:ch14}}{equation.2.3.93}{}}
\newlabel{m:eq:ywitheta@cref}{{[equation][93][2]2.93}{[1][50][]50}}
\newlabel{m:sec:mocwhat}{{2.3.3.3}{50}{What the example shows}{subsubsection.2.3.3.3}{}}
\newlabel{m:sec:mocwhat@cref}{{[subsubsection][3][2,3,3]2.3.3.3}{[1][50][]50}}
\newlabel{m:app:smallpops}{{2.A}{50}{Survival of small populations}{section.2.appendix.A}{}}
\newlabel{m:app:smallpops@cref}{{[section][1][2]2.A}{[1][50][]50}}
\newlabel{m:fig:earlySn}{{2.16}{51}{Survival probability over the first four generations from a single progenitor. Even in the supercritical case ($p=0.55$) the early decay is steep, and the three curves have barely separated by $n=4$}{figure.2.16}{}}
\newlabel{m:fig:earlySn@cref}{{[figure][16][2]2.16}{[1][51][]51}}
\newlabel{m:eq:Snk}{{2.94}{51}{An initial cohort of size \texorpdfstring {$k$}{k}}{equation.2.appendix.A.94}{}}
\newlabel{m:eq:Snk@cref}{{[equation][94][2]2.94}{[1][51][]51}}
\newlabel{m:app:absorption}{{2.B}{51}{The absorption models in full}{section.2.appendix.B}{}}
\newlabel{m:app:absorption@cref}{{[section][2][2]2.B}{[1][51][]51}}
\newlabel{m:fig:kvals}{{2.17}{52}{Survival probability $S^{(k)}_n$ from \cref {m:eq:Snk} out to $500$ generations at exact criticality, $p=0.5$, for $k=1,10,100,1000$ (blue through to red). A larger founding cohort buys longevity, but at criticality every curve is still headed for zero}{figure.2.17}{}}
\newlabel{m:fig:kvals@cref}{{[figure][17][2]2.17}{[1][51][]52}}
\newlabel{m:fig:kvals2}{{2.18}{52}{As \cref {m:fig:kvals}, but with $p$ displaced slightly from criticality: $p=0.505$ on the left (supercritical), $p=0.495$ on the right (subcritical). A displacement of half a per cent separates curves that were nearly coincident in \cref {m:fig:kvals}}{figure.2.18}{}}
\newlabel{m:fig:kvals2@cref}{{[figure][18][2]2.18}{[1][51][]52}}
\newlabel{m:app:absonly}{{2.B.1}{52}{Absorption only}{subsection.2.appendix.B.1}{}}
\newlabel{m:app:absonly@cref}{{[subsection][1][2,2]2.B.1}{[1][52][]52}}
\newlabel{m:eq:FK1}{{2.95}{53}{Absorption only}{equation.2.appendix.B.95}{}}
\newlabel{m:eq:FK1@cref}{{[equation][95][2]2.95}{[1][52][]53}}
\newlabel{m:eq:PDE1}{{2.96}{53}{Absorption only}{equation.2.appendix.B.96}{}}
\newlabel{m:eq:PDE1@cref}{{[equation][96][2]2.96}{[1][53][]53}}
\newlabel{m:eq:indRelPGF}{{2.97}{53}{Absorption only}{equation.2.appendix.B.97}{}}
\newlabel{m:eq:indRelPGF@cref}{{[equation][97][2]2.97}{[1][53][]53}}
\newlabel{m:eq:Gabsonly}{{2.98}{53}{Absorption only}{equation.2.appendix.B.98}{}}
\newlabel{m:eq:Gabsonly@cref}{{[equation][98][2]2.98}{[1][53][]53}}
\newlabel{m:fig:abs1}{{2.19}{53}{State probabilities for the single-particle absorption-only model. The particle begins outside the compartment ($p_{0,1}$, decaying exponentially at rate $\alpha $) and ends inside it ($p_{1,0}$). There is nowhere else for it to go, so the two curves sum to one at every time. \explorerlink {https://markov.volvox.site/-/adam74aldridge/absorption-only-model}{Simulate}}{figure.2.19}{}}
\newlabel{m:fig:abs1@cref}{{[figure][19][2]2.19}{[1][53][]53}}
\newlabel{m:app:absbirthdeath}{{2.B.2}{53}{Absorption--birth--death}{subsection.2.appendix.B.2}{}}
\newlabel{m:app:absbirthdeath@cref}{{[subsection][2][2,2]2.B.2}{[1][53][]53}}
\newlabel{m:eq:FK3}{{2.99}{54}{Absorption--birth--death}{equation.2.appendix.B.99}{}}
\newlabel{m:eq:FK3@cref}{{[equation][99][2]2.99}{[1][53][]54}}
\newlabel{m:eq:PDE3}{{2.100}{54}{Absorption--birth--death}{equation.2.appendix.B.100}{}}
\newlabel{m:eq:PDE3@cref}{{[equation][100][2]2.100}{[1][54][]54}}
\newlabel{m:eq:ch21}{{2.101}{54}{Absorption--birth--death}{equation.2.appendix.B.101}{}}
\newlabel{m:eq:ch21@cref}{{[equation][101][2]2.101}{[1][54][]54}}
\newlabel{m:eq:ch22}{{2.102}{54}{Absorption--birth--death}{equation.2.appendix.B.102}{}}
\newlabel{m:eq:ch22@cref}{{[equation][102][2]2.102}{[1][54][]54}}
\newlabel{m:eq:ch23}{{2.103}{54}{Absorption--birth--death}{equation.2.appendix.B.103}{}}
\newlabel{m:eq:ch23@cref}{{[equation][103][2]2.103}{[1][54][]54}}
\newlabel{m:eq:ch24}{{2.104}{54}{Absorption--birth--death}{equation.2.appendix.B.104}{}}
\newlabel{m:eq:ch24@cref}{{[equation][104][2]2.104}{[1][54][]54}}
\newlabel{m:eq:GetaN2}{{2.105}{54}{The equations for $G$ and $t$}{equation.2.appendix.B.105}{}}
\newlabel{m:eq:GetaN2@cref}{{[equation][105][2]2.105}{[1][54][]54}}
\newlabel{m:eq:xpm}{{2.106}{54}{The equation for $x$, \eqref {m:eq:ch23}}{equation.2.appendix.B.106}{}}
\newlabel{m:eq:xpm@cref}{{[equation][106][2]2.106}{[1][54][]54}}
\newlabel{m:eq:b1}{{2.107}{54}{The equation for $x$, \eqref {m:eq:ch23}}{equation.2.appendix.B.107}{}}
\newlabel{m:eq:b1@cref}{{[equation][107][2]2.107}{[1][54][]54}}
\newlabel{m:eq:xratio}{{2.108}{54}{The equation for $x$, \eqref {m:eq:ch23}}{equation.2.appendix.B.108}{}}
\newlabel{m:eq:xratio@cref}{{[equation][108][2]2.108}{[1][54][]54}}
\newlabel{m:eq:xabd}{{2.109}{54}{The equation for $x$, \eqref {m:eq:ch23}}{equation.2.appendix.B.109}{}}
\newlabel{m:eq:xabd@cref}{{[equation][109][2]2.109}{[1][54][]54}}
\newlabel{m:eq:yint}{{2.110}{55}{The equation for $y$, \eqref {m:eq:ch24}}{equation.2.appendix.B.110}{}}
\newlabel{m:eq:yint@cref}{{[equation][110][2]2.110}{[1][54][]55}}
\newlabel{m:eq:xgeom}{{2.111}{55}{The equation for $y$, \eqref {m:eq:ch24}}{equation.2.appendix.B.111}{}}
\newlabel{m:eq:xgeom@cref}{{[equation][111][2]2.111}{[1][54][]55}}
\newlabel{m:eq:Psidef}{{2.112}{55}{The equation for $y$, \eqref {m:eq:ch24}}{equation.2.appendix.B.112}{}}
\newlabel{m:eq:Psidef@cref}{{[equation][112][2]2.112}{[1][55][]55}}
\newlabel{m:eq:etasolved}{{2.113}{55}{The equation for $y$, \eqref {m:eq:ch24}}{equation.2.appendix.B.113}{}}
\newlabel{m:eq:etasolved@cref}{{[equation][113][2]2.113}{[1][55][]55}}
\newlabel{m:prop:Gabd}{{2.8}{55}{Generating function for absorption--birth--death}{theorem.2.8}{}}
\newlabel{m:prop:Gabd@cref}{{[proposition][8][2]2.8}{[1][55][]55}}
\newlabel{m:eq:Gabd}{{2.114}{55}{Generating function for absorption--birth--death}{equation.2.appendix.B.114}{}}
\newlabel{m:eq:Gabd@cref}{{[equation][114][2]2.114}{[1][55][]55}}
\newlabel{m:app:abdregular}{{2.B.3}{55}{A parameter-regular representation}{subsection.2.appendix.B.3}{}}
\newlabel{m:app:abdregular@cref}{{[subsection][3][2,2]2.B.3}{[1][55][]55}}
\newlabel{m:eq:interiorbackward}{{2.115}{56}{A parameter-regular representation}{equation.2.appendix.B.115}{}}
\newlabel{m:eq:interiorbackward@cref}{{[equation][115][2]2.115}{[1][55][]56}}
\newlabel{m:eq:interiorratio}{{2.116}{56}{A parameter-regular representation}{equation.2.appendix.B.116}{}}
\newlabel{m:eq:interiorratio@cref}{{[equation][116][2]2.116}{[1][55][]56}}
\newlabel{m:eq:interiorcritical}{{2.117}{56}{A parameter-regular representation}{equation.2.appendix.B.117}{}}
\newlabel{m:eq:interiorcritical@cref}{{[equation][117][2]2.117}{[1][56][]56}}
\newlabel{m:eq:abdconvolution}{{2.118}{56}{A parameter-regular representation}{equation.2.appendix.B.118}{}}
\newlabel{m:eq:abdconvolution@cref}{{[equation][118][2]2.118}{[1][56][]56}}
\newlabel{m:app:abdhyp}{{2.B.4}{56}{Hypergeometric form, domain and resonance}{subsection.2.appendix.B.4}{}}
\newlabel{m:app:abdhyp@cref}{{[subsection][4][2,2]2.B.4}{[1][56][]56}}
\newlabel{m:eq:2F1def}{{2.119}{56}{Hypergeometric form, domain and resonance}{equation.2.appendix.B.119}{}}
\newlabel{m:eq:2F1def@cref}{{[equation][119][2]2.119}{[1][56][]56}}
\newlabel{m:eq:pochcancel}{{2.120}{56}{Hypergeometric form, domain and resonance}{equation.2.appendix.B.120}{}}
\newlabel{m:eq:pochcancel@cref}{{[equation][120][2]2.120}{[1][56][]56}}
\newlabel{m:eq:Psihyp}{{2.121}{56}{Hypergeometric form, domain and resonance}{equation.2.appendix.B.121}{}}
\newlabel{m:eq:Psihyp@cref}{{[equation][121][2]2.121}{[1][56][]56}}
\newlabel{m:rem:domain}{{2.9}{56}{Domain of validity}{theorem.2.9}{}}
\newlabel{m:rem:domain@cref}{{[remark][9][2]2.9}{[1][56][]56}}
\newlabel{m:rem:resonance}{{2.10}{57}{Resonance}{theorem.2.10}{}}
\newlabel{m:rem:resonance@cref}{{[remark][10][2]2.10}{[1][57][]57}}
\newlabel{m:eq:resonance}{{2.122}{57}{Resonance}{equation.2.appendix.B.122}{}}
\newlabel{m:eq:resonance@cref}{{[equation][122][2]2.122}{[1][57][]57}}
\newlabel{m:app:coefrecovery}{{2.B.5}{57}{Recovering state probabilities}{subsection.2.appendix.B.5}{}}
\newlabel{m:app:coefrecovery@cref}{{[subsection][5][2,2]2.B.5}{[1][57][]57}}
\newlabel{m:eq:bivariatecoef}{{2.123}{57}{Recovering state probabilities}{equation.2.appendix.B.123}{}}
\newlabel{m:eq:bivariatecoef@cref}{{[equation][123][2]2.123}{[1][57][]57}}
\newlabel{m:eq:hdef}{{2.124}{57}{Recovering state probabilities}{equation.2.appendix.B.124}{}}
\newlabel{m:eq:hdef@cref}{{[equation][124][2]2.124}{[1][57][]57}}
\newlabel{m:eq:abdcoef}{{2.125}{57}{Recovering state probabilities}{equation.2.appendix.B.125}{}}
\newlabel{m:eq:abdcoef@cref}{{[equation][125][2]2.125}{[1][57][]57}}
\newlabel{m:app:hyp}{{2.C}{57}{An integral identity for the hypergeometric function}{section.2.appendix.C}{}}
\newlabel{m:app:hyp@cref}{{[section][3][2]2.C}{[1][57][]57}}
\newlabel{m:eq:Iform}{{2.126}{57}{An integral identity for the hypergeometric function}{equation.2.appendix.C.126}{}}
\newlabel{m:eq:Iform@cref}{{[equation][126][2]2.126}{[1][57][]57}}
\newlabel{m:prop:hypIden}{{2.11}{58}{}{theorem.2.11}{}}
\newlabel{m:prop:hypIden@cref}{{[proposition][11][2]2.11}{[1][58][]58}}
\newlabel{m:eq:hypIden}{{2.127}{58}{}{equation.2.appendix.C.127}{}}
\newlabel{m:eq:hypIden@cref}{{[equation][127][2]2.127}{[1][58][]58}}
\newlabel{m:eq:idenHyp}{{2.128}{58}{An integral identity for the hypergeometric function}{equation.2.appendix.C.128}{}}
\newlabel{m:eq:idenHyp@cref}{{[equation][128][2]2.128}{[1][58][]58}}
\newlabel{m:eq:backat}{{2.129}{58}{An integral identity for the hypergeometric function}{equation.2.appendix.C.129}{}}
\newlabel{m:eq:backat@cref}{{[equation][129][2]2.129}{[1][58][]58}}
\newlabel{m:app:extract}{{2.D}{59}{Extracting state probabilities from a generating function}{section.2.appendix.D}{}}
\newlabel{m:app:extract@cref}{{[section][4][2]2.D}{[1][58][]59}}
\newlabel{m:eq:examplePgf}{{2.130}{59}{Extracting state probabilities from a generating function}{equation.2.appendix.D.130}{}}
\newlabel{m:eq:examplePgf@cref}{{[equation][130][2]2.130}{[1][58][]59}}
\newlabel{m:eq:series1}{{2.131}{59}{Extracting state probabilities from a generating function}{equation.2.appendix.D.131}{}}
\newlabel{m:eq:series1@cref}{{[equation][131][2]2.131}{[1][58][]59}}
\newlabel{m:eq:stateProb}{{2.132}{59}{Extracting state probabilities from a generating function}{equation.2.appendix.D.132}{}}
\newlabel{m:eq:stateProb@cref}{{[equation][132][2]2.132}{[1][59][]59}}
\newlabel{m:eq:bivariateCauchy}{{2.133}{59}{Extracting state probabilities from a generating function}{equation.2.appendix.D.133}{}}
\newlabel{m:eq:bivariateCauchy@cref}{{[equation][133][2]2.133}{[1][59][]59}}
\newlabel{m:eq:multilinear}{{2.134}{59}{The multilinear case}{equation.2.appendix.D.134}{}}
\newlabel{m:eq:multilinear@cref}{{[equation][134][2]2.134}{[1][59][]59}}
\newlabel{m:eq:derivs}{{2.135}{59}{The multilinear case}{equation.2.appendix.D.135}{}}
\newlabel{m:eq:derivs@cref}{{[equation][135][2]2.135}{[1][59][]59}}
\newlabel{m:eq:stateProbSpec}{{2.136}{59}{The multilinear case}{equation.2.appendix.D.136}{}}
\newlabel{m:eq:stateProbSpec@cref}{{[equation][136][2]2.136}{[1][59][]59}}
\newlabel{gw:ch:galton-watson}{{3}{61}{Galton--Watson branching and the constant $A(p)$}{chapter.3}{}}
\newlabel{gw:ch:galton-watson@cref}{{[chapter][3][]3}{[1][61][]61}}
\newlabel{fig:example}{{\caption@xref {fig:example}{ on input line 76}}{62}{Introduction}{section.3.1}{}}
\newlabel{fig:example@cref}{{[section][1][3]3.1}{[1][62][]62}}
\newlabel{eq:unrec}{{3.1}{64}{Survival probability and extinction}{equation.3.2.1}{}}
\newlabel{eq:unrec@cref}{{[equation][1][3]3.1}{[1][64][]64}}
\newlabel{eq:SnIter}{{3.2}{64}{Survival probability and extinction}{equation.3.2.2}{}}
\newlabel{eq:SnIter@cref}{{[equation][2][3]3.2}{[1][64][]64}}
\newlabel{fig:phase_transition}{{3.1}{65}{Fixed points of the survival recursion $S\mapsto 2pS-pS^2$. Each panel shows the map (blue) against the diagonal (dashed), with fixed points marked in red. Below and at criticality the only fixed point in $[0,1]$ is the origin; for $p>\tfrac 12$ a stable positive root $S_\infty =(2p-1)/p$ separates from it (at $p=0.74$, this is approximately $0.649$)}{figure.3.1}{}}
\newlabel{fig:phase_transition@cref}{{[figure][1][3]3.1}{[1][64][]65}}
\newlabel{sec:limcondmean}{{3.3}{65}{Limiting Conditional Mean \texorpdfstring {$\mathbb {E}(Z_n|Z_n>0) = 1/A(p)$}{E(Zn|Zn>0) = 1/A(p)}}{section.3.3}{}}
\newlabel{sec:limcondmean@cref}{{[section][3][3]3.3}{[1][65][]65}}
\newlabel{eq:Ap-def-main}{{3.3}{66}{Limiting Conditional Mean \texorpdfstring {$\mathbb {E}(Z_n|Z_n>0) = 1/A(p)$}{E(Zn|Zn>0) = 1/A(p)}}{equation.3.3.3}{}}
\newlabel{eq:Ap-def-main@cref}{{[equation][3][3]3.3}{[1][66][]66}}
\newlabel{eq:Sn-asymp}{{3.4}{66}{Limiting Conditional Mean \texorpdfstring {$\mathbb {E}(Z_n|Z_n>0) = 1/A(p)$}{E(Zn|Zn>0) = 1/A(p)}}{equation.3.3.4}{}}
\newlabel{eq:Sn-asymp@cref}{{[equation][4][3]3.4}{[1][66][]66}}
\newlabel{fig:condmean}{{3.2}{67}{The conditional mean $(2p)^n/S_n$ against generation number $n$, with $S_n$ obtained by iterating \eqref {eq:SnIter}. From bottom to top, $p=0.3,\,0.4,\,0.45,\,0.48,\,0.49$. Each curve approaches the constant $1/A(p)$; the limiting values are $2.825$, $4.208$, $6.893$, $14.703$ and $27.476$. The lowest three have essentially converged by $n=100$, but the $p=0.49$ curve has reached only $24.1$ of its limiting $27.5$}{figure.3.2}{}}
\newlabel{fig:condmean@cref}{{[figure][2][3]3.2}{[1][66][]67}}
\newlabel{sec:condvarmain}{{3.3.1}{67}{The limiting conditional variance}{subsection.3.3.1}{}}
\newlabel{sec:condvarmain@cref}{{[subsection][1][3,3]3.3.1}{[1][66][]67}}
\newlabel{eq:condvar}{{3.5}{67}{The limiting conditional variance}{equation.3.3.5}{}}
\newlabel{eq:condvar@cref}{{[equation][5][3]3.5}{[1][66][]67}}
\newlabel{fig:yalimit}{{3.3}{68}{Ensemble of binary Galton--Watson realisations at $p=0.49732333$, followed for $670$ generations. Grey paths are individual population sizes $Z_n$; the blue curve is the mean among lineages still alive; the red dashed line is the limiting conditional mean $1/A(p)$. \explorerlink {https://files.volvox.site/users/DmOAcD4aqPPnvhOsvqYYqFgxEZ53/postFiles/1769320791182-cond.html}{Yaglom visualiser}}{figure.3.3}{}}
\newlabel{fig:yalimit@cref}{{[figure][3][3]3.3}{[1][67][]68}}
\newlabel{sec:product}{{3.4}{68}{An infinite-product representation of \texorpdfstring {$A(p)$}{A(p)}}{section.3.4}{}}
\newlabel{sec:product@cref}{{[section][4][3]3.4}{[1][68][]68}}
\newlabel{eq:prodFormula}{{3.6}{69}{An infinite-product representation of \texorpdfstring {$A(p)$}{A(p)}}{equation.3.4.6}{}}
\newlabel{eq:prodFormula@cref}{{[equation][6][3]3.6}{[1][69][]69}}
\newlabel{sec:koenigs}{{3.5}{70}{The logistic map and Koenigs representation}{section.3.5}{}}
\newlabel{sec:koenigs@cref}{{[section][5][3]3.5}{[1][69][]70}}
\newlabel{eq:logistic}{{3.7}{70}{The logistic map and Koenigs representation}{equation.3.5.7}{}}
\newlabel{eq:logistic@cref}{{[equation][7][3]3.7}{[1][69][]70}}
\newlabel{fig:period_double}{{3.4}{70}{\emph {Aside, not used in the argument.} Phase diagram of the logistic map that shows the period doubling of stable states, ultimately resulting in chaotic behaviour. Blue sample line on the left at $r=3.5255$, right plot shows a trajectory of the map with the same value of $r$. The subcritical branching window is $r=2p\in (0,1)$, at the far left of this diagram, where the only attractor is the fixed point at the origin. \explorerlink {https://files.volvox.site/users/DmOAcD4aqPPnvhOsvqYYqFgxEZ53/postFiles/1771487127831-logistic_spiderweb.html}{Bifurcation visualiser}}{figure.3.4}{}}
\newlabel{fig:period_double@cref}{{[figure][4][3]3.4}{[1][70][]70}}
\newlabel{def1}{{3.1}{70}{}{theorem.3.1}{}}
\newlabel{def1@cref}{{[definition][1][3]3.1}{[1][70][]70}}
\newlabel{eq:second}{{3.8}{71}{The identity \texorpdfstring {$A(p)=2\psi _{2p}(1/2)$}{A(p) = 2 psi\_2p(1/2)}}{equation.3.5.8}{}}
\newlabel{eq:second@cref}{{[equation][8][3]3.8}{[1][71][]71}}
\newlabel{eq:fourth}{{3.9}{71}{The identity \texorpdfstring {$A(p)=2\psi _{2p}(1/2)$}{A(p) = 2 psi\_2p(1/2)}}{equation.3.5.9}{}}
\newlabel{eq:fourth@cref}{{[equation][9][3]3.9}{[1][71][]71}}
\newlabel{eq:induction}{{3.10}{71}{The identity \texorpdfstring {$A(p)=2\psi _{2p}(1/2)$}{A(p) = 2 psi\_2p(1/2)}}{equation.3.5.10}{}}
\newlabel{eq:induction@cref}{{[equation][10][3]3.10}{[1][71][]71}}
\newlabel{sec:closedform}{{3.5.2}{71}{The closed-form question}{subsection.3.5.2}{}}
\newlabel{sec:closedform@cref}{{[subsection][2][3,5]3.5.2}{[1][71][]71}}
\newlabel{fig:power_law}{{3.5}{73}{Power-law behaviour in a critical Galton--Watson branching process ($p = 0.5$).\\ (a) Survival probability $S_n = \mathbb {P}(Z_n > 0)$, showing the late generation asymptotic scaling $S_n \sim 2/n$.\\ (b) Probability mass function of the extinction time $T_{\mathrm {extinct}} = \inf \{n : Z_n = 0\}$, exhibiting a heavy tail $\mathbb {P}(T_{\mathrm {extinct}} = n) \sim 2/n^2$.\\ (c) Probability mass function of the total progeny $\mathcal {T} = \sum _{k \ge 0} Z_k$, displaying the classical critical branching asymptotic $\mathbb {P}(\mathcal {T} = n) \sim n^{-3/2}$. ($10^6$ realisations)}{figure.3.5}{}}
\newlabel{fig:power_law@cref}{{[figure][5][3]3.5}{[1][72][]73}}
\newlabel{eq:critrec}{{3.11}{75}{Expected lifetime}{equation.3.6.11}{}}
\newlabel{eq:critrec@cref}{{[equation][11][3]3.11}{[1][74][]75}}
\newlabel{eq:recincrement}{{3.12}{75}{Expected lifetime}{equation.3.6.12}{}}
\newlabel{eq:recincrement@cref}{{[equation][12][3]3.12}{[1][75][]75}}
\newlabel{twoovrn_sn}{{3.13}{75}{Expected lifetime}{equation.3.6.13}{}}
\newlabel{twoovrn_sn@cref}{{[equation][13][3]3.13}{[1][75][]75}}
\newlabel{sec:nearcrit}{{3.7}{75}{Near-critical behaviour of \texorpdfstring {$A(p)$}{A(p)}}{section.3.7}{}}
\newlabel{sec:nearcrit@cref}{{[section][7][3]3.7}{[1][75][]75}}
\newlabel{eq:nearcrit-profile}{{3.17}{76}{Continuum ODE}{equation.3.7.17}{}}
\newlabel{eq:nearcrit-profile@cref}{{[equation][17][3]3.17}{[1][76][]76}}
\newlabel{eq:ratiomonitor}{{3.18}{76}{Direct numerical approximation}{equation.3.8.18}{}}
\newlabel{eq:ratiomonitor@cref}{{[equation][18][3]3.18}{[1][76][]76}}
\newlabel{eq:trunc-bound}{{3.19}{76}{Direct numerical approximation}{equation.3.8.19}{}}
\newlabel{eq:trunc-bound@cref}{{[equation][19][3]3.19}{[1][76][]76}}
\newlabel{eq:stoprule}{{3.20}{77}{Direct numerical approximation}{equation.3.8.20}{}}
\newlabel{eq:stoprule@cref}{{[equation][20][3]3.20}{[1][76][]77}}
\newlabel{eq:cost}{{3.21}{77}{Direct numerical approximation}{equation.3.8.21}{}}
\newlabel{eq:cost@cref}{{[equation][21][3]3.21}{[1][77][]77}}
\newlabel{eq:crossover}{{3.22}{77}{Product/asymptotic hybrid near criticality}{equation.3.8.22}{}}
\newlabel{eq:crossover@cref}{{[equation][22][3]3.22}{[1][77][]77}}
\newlabel{eq:hybrid}{{3.23}{78}{Product/asymptotic hybrid near criticality}{equation.3.8.23}{}}
\newlabel{eq:hybrid@cref}{{[equation][23][3]3.23}{[1][77][]78}}
\newlabel{sec:koenigsvisual}{{3.9}{78}{The linearising coordinate in the complex plane}{section.3.9}{}}
\newlabel{sec:koenigsvisual@cref}{{[section][9][3]3.9}{[1][78][]78}}
\newlabel{fig:koenigs-orbit}{{3.7}{78}{Orbit fates for $f_r(z)=rz(1-z)$ at $r=0.64$ ($p=0.32$). The seeds $z=-0.229139+0.514714i$, $z=-0.812055+0.625998i$ and $z=1.043377+1.900730i$ in (a)--(c) give, respectively, rapid capture, capture after a long transient and escape. The first two therefore lie in the same attracting basin}{figure.3.7}{}}
\newlabel{fig:koenigs-orbit@cref}{{[figure][7][3]3.7}{[1][78][]78}}
\newlabel{fig:koenigs-surface}{{3.8}{80}{The Koenigs coordinate at $r=0.64$ ($p=0.32$). In (a), the marked seed $z=\tfrac 12$ has product-computed height $A(p)=0.335521$ on $|2\psi _r(z)|$; the dashed line matches it to the cyan $A$-against-$p$ curve. Panels (b) and (c) lower and remove the surface. The funnels mark $\psi _r(0)=\psi _r(1)=0$, and the reciprocal height gives the quasi-stationary mean $1/A(p)=2.980439$. \explorerlink {https://files.volvox.site/users/69b423f5faa2bfbb5b696c61/pages/69b423fcfaa2bfbb5b696c77/files/1787185484035-q1arbz-koenigs_explorer.html}{Koenigs visualiser}}{figure.3.8a}{}}
\newlabel{fig:koenigs-surface@cref}{{[figure][8][3]3.8}{[1][80][]80}}
\newlabel{fig:koenigs-seq}{{3.9}{81}{Continuation from the real branching slice at $\operatorname {Re}r=0.64$. Panel (a) is the branching case $p=0.32$, with $A(p)=0.335521$ and $1/A(p)=2.980439$. In (b)--(e), $|r|$ is $0.646$, $0.720$, $0.912$ and $0.986$: the Koenigs coordinate persists, but $r/2$ is no longer a probability. In (f), $|r|\approx 1.009$ and the attracting chart no longer applies. The marked point is $z=\tfrac 12$ throughout}{figure.3.9}{}}
\newlabel{fig:koenigs-seq@cref}{{[figure][9][3]3.9}{[1][80][]81}}
\newlabel{fig:koenigs-gallery}{{3.10}{82}{Five deep zooms outside the attracting window, with $|r|$ equal to $1.147$, $1.582$, $1.582$, $1.568$ and $1.015$ in panel order. Since all exceed one, none shows the attracting Koenigs chart; the images instead place the branching segment $r=2p\in (0,1)$ within the wider complex quadratic family. \explorerlink {https://files.volvox.site/users/69b423f5faa2bfbb5b696c61/pages/69b423fcfaa2bfbb5b696c77/files/1787185484035-q1arbz-koenigs_explorer.html}{Koenigs visualiser}}{figure.3.10}{}}
\newlabel{fig:koenigs-gallery@cref}{{[figure][10][3]3.10}{[1][81][]82}}
\newlabel{app:hyper}{{3.A}{83}{No elementary formula for the Koenigs function of the logistic map when\texorpdfstring {\protect \\}{ }$0<r<1$}{section.3.appendix.A}{}}
\newlabel{app:hyper@cref}{{[section][1][3]3.A}{[1][83][]83}}
\newlabel{eq:koenigs}{{3.25}{84}{The question}{equation.3.appendix.A.25}{}}
\newlabel{eq:koenigs@cref}{{[equation][25][3]3.25}{[1][84][]84}}
\newlabel{eq:ade}{{3.26}{84}{}{equation.3.appendix.A.26}{}}
\newlabel{eq:ade@cref}{{[equation][26][3]3.26}{[1][84][]84}}
\newlabel{thm:bb}{{3.3}{85}{Becker--Bergweiler}{theorem.3.3}{}}
\newlabel{thm:bb@cref}{{[theorem][3][3]3.3}{[1][84][]85}}
\newlabel{eq:source-schroeder}{{3.27}{85}{Becker--Bergweiler}{equation.3.appendix.A.27}{}}
\newlabel{eq:source-schroeder@cref}{{[equation][27][3]3.27}{[1][85][]85}}
\newlabel{lem:inverse}{{3.4}{85}{}{theorem.3.4}{}}
\newlabel{lem:inverse@cref}{{[lemma][4][3]3.4}{[1][85][]85}}
\newlabel{gw:thm:main}{{3.5}{85}{}{theorem.3.5}{}}
\newlabel{gw:thm:main@cref}{{[theorem][5][3]3.5}{[1][85][]85}}
\newlabel{eq:limit}{{3.28}{86}{What the result means}{equation.3.appendix.A.28}{}}
\newlabel{eq:limit@cref}{{[equation][28][3]3.28}{[1][86][]86}}
\newlabel{sec:search}{{3.B}{86}{Search for an elementary formula}{section.3.appendix.B}{}}
\newlabel{sec:search@cref}{{[section][2][3]3.B}{[1][86][]86}}
\newlabel{fig:gasearch}{{3.11}{87}{A genetic algorithm evolving polynomials to fit numerically generated $A(p)$ data. Each genome encodes the terms and coefficients of a polynomial. Grey curves are candidates from a population of $50$ (mutation rate $30\%$); blue points are the data; the red curve is the best fit after $90$ generations (MSE $1.35\times 10^{-4}$). The first visualiser is a vibe-coded interface to a Python implementation I wrote by hand; v2 is a later vibe-coded attempt at a better solver, made before the Koenigs connection. \explorerlink {https://files.volvox.site/users/69b423f5faa2bfbb5b696c61/pages/69b423fcfaa2bfbb5b696c77/files/1787182834062-65kw9b-Genetic_algorithm_for_polynomial_searching.html}{Polynomial-fit visualiser} \explorerlink {https://files.volvox.site/users/69b423f5faa2bfbb5b696c61/pages/69b423fcfaa2bfbb5b696c77/files/1787182532533-vkazf0-Evolutionary_function_finder_2.html}{Polynomial-fit visualiser (v2)}}{figure.3.11}{}}
\newlabel{fig:gasearch@cref}{{[figure][11][3]3.11}{[1][86][]87}}
\newlabel{app:condvar}{{3.C}{89}{The limiting conditional variance}{section.3.appendix.C}{}}
\newlabel{app:condvar@cref}{{[section][3][3]3.C}{[1][88][]89}}
\newlabel{eq:phinprime}{{3.29}{89}{The limiting conditional variance}{equation.3.appendix.C.29}{}}
\newlabel{eq:phinprime@cref}{{[equation][29][3]3.29}{[1][89][]89}}
\newlabel{eq:phinsecond}{{3.30}{89}{The limiting conditional variance}{equation.3.appendix.C.30}{}}
\newlabel{eq:phinsecond@cref}{{[equation][30][3]3.30}{[1][89][]89}}
\newlabel{eq:vnrecursion}{{3.31}{89}{The limiting conditional variance}{equation.3.appendix.C.31}{}}
\newlabel{eq:vnrecursion@cref}{{[equation][31][3]3.31}{[1][89][]89}}
\newlabel{eq:factorialmoment}{{3.32}{89}{The limiting conditional variance}{equation.3.appendix.C.32}{}}
\newlabel{eq:factorialmoment@cref}{{[equation][32][3]3.32}{[1][89][]89}}
\newlabel{eq:condsecondmoment}{{3.33}{89}{The limiting conditional variance}{equation.3.appendix.C.33}{}}
\newlabel{eq:condsecondmoment@cref}{{[equation][33][3]3.33}{[1][89][]89}}
\newlabel{eq:condvar-app}{{3.34}{89}{The limiting conditional variance}{equation.3.appendix.C.34}{}}
\newlabel{eq:condvar-app@cref}{{[equation][34][3]3.34}{[1][89][]89}}
\newlabel{app:nearcritmoments}{{3.D}{89}{Near-critical growth of mean and variance}{section.3.appendix.D}{}}
\newlabel{app:nearcritmoments@cref}{{[section][4][3]3.D}{[1][89][]89}}
\newlabel{eq:A-nearcrit}{{3.35}{89}{Near-critical growth of mean and variance}{equation.3.appendix.D.35}{}}
\newlabel{eq:A-nearcrit@cref}{{[equation][35][3]3.35}{[1][89][]89}}
\newlabel{eq:mean-nearcrit}{{3.36}{90}{Near-critical growth of mean and variance}{equation.3.appendix.D.36}{}}
\newlabel{eq:mean-nearcrit@cref}{{[equation][36][3]3.36}{[1][89][]90}}
\newlabel{eq:var-nearcrit}{{3.37}{90}{Near-critical growth of mean and variance}{equation.3.appendix.D.37}{}}
\newlabel{eq:var-nearcrit@cref}{{[equation][37][3]3.37}{[1][90][]90}}
\newlabel{bdc:ch:core}{{4}{91}{The birth--death--catastrophe process}{chapter.4}{}}
\newlabel{bdc:ch:core@cref}{{[chapter][4][]4}{[1][91][]91}}
\newlabel{bdc:sec:intro}{{4.1}{91}{Introduction}{section.4.1}{}}
\newlabel{bdc:sec:intro@cref}{{[section][1][4]4.1}{[1][91][]91}}
\newlabel{bdc:rem:novelty}{{4.1}{92}{What is classical here, and what is developed here}{theorem.4.1}{}}
\newlabel{bdc:rem:novelty@cref}{{[remark][1][4]4.1}{[1][92][]92}}
\newlabel{bdc:sec:process}{{4.2}{93}{The process}{section.4.2}{}}
\newlabel{bdc:sec:process@cref}{{[section][2][4]4.2}{[1][92][]93}}
\newlabel{bdc:def:bdc}{{4.2}{93}{Birth--death--catastrophe process}{theorem.4.2}{}}
\newlabel{bdc:def:bdc@cref}{{[definition][2][4]4.2}{[1][93][]93}}
\newlabel{bdc:eq:rates}{{4.1}{93}{Birth--death--catastrophe process}{equation.4.2.1}{}}
\newlabel{bdc:eq:rates@cref}{{[equation][1][4]4.1}{[1][93][]93}}
\newlabel{bdc:fig:ratediagram}{{4.1}{93}{Linear birth--death dynamics (top) and the birth--death--catastrophe process used throughout this chapter (bottom). Double circles are absorbing. Ordinary extinction enters $0$ only from state $1$; catastrophe jumps from any positive count to the distinct state $H$, at total rate $\delta n$}{figure.4.1}{}}
\newlabel{bdc:fig:ratediagram@cref}{{[figure][1][4]4.1}{[1][93][]93}}
\newlabel{bdc:sec:paths}{{4.2.1}{93}{Sample paths}{subsection.4.2.1}{}}
\newlabel{bdc:sec:paths@cref}{{[subsection][1][4,2]4.2.1}{[1][93][]93}}
\newlabel{bdc:fig:gillespie}{{4.2}{94}{One hundred exact realisations of the load $X_t$ (top) and the released count $W_t$ (bottom), generated by the Gillespie algorithm \cite {Gillespie1977}. The step of $X_t$ to the plotting baseline marks entry into $H$, not the count state $0$; with $\mu =0$ no realisation reaches $0$ before rupturing. Each $W_t$ path remains at zero until rupture, then jumps to the load present at that instant. $\lambda =1$, $\mu =0$, $\delta =0.1$, $X_0=1$}{figure.4.2}{}}
\newlabel{bdc:fig:gillespie@cref}{{[figure][2][4]4.2}{[1][93][]94}}
\newlabel{bdc:rem:Jmixes}{{4.3}{94}{The mean load describes no typical path}{theorem.4.3}{}}
\newlabel{bdc:rem:Jmixes@cref}{{[remark][3][4]4.3}{[1][94][]94}}
\newlabel{bdc:fig:fatepaths}{{4.3}{95}{Twelve exact realisations of $X_t$, coloured by terminal fate. Grey paths ending in crosses reach internal extinction; vermillion paths ending in filled circles enter $H$. The bold curve is the exact mean $J(t)$. $\lambda =1$, $\mu =0.2$, $\delta =0.05$, $X_0=1$; seed 154. \explorerlink {https://markov.volvox.site/-/adam74aldridge/birth-death-catastrophe-process}{Simulate}}{figure.4.3}{}}
\newlabel{bdc:fig:fatepaths@cref}{{[figure][3][4]4.3}{[1][93][]95}}
\newlabel{bdc:sec:released}{{4.2.2}{95}{The released count and the joint process}{subsection.4.2.2}{}}
\newlabel{bdc:sec:released@cref}{{[subsection][2][4,2]4.2.2}{[1][95][]95}}
\newlabel{bdc:eq:release}{{4.2}{95}{The released count and the joint process}{equation.4.2.2}{}}
\newlabel{bdc:eq:release@cref}{{[equation][2][4]4.2}{[1][95][]95}}
\newlabel{bdc:eq:joint}{{4.3}{95}{The released count and the joint process}{equation.4.2.3}{}}
\newlabel{bdc:eq:joint@cref}{{[equation][3][4]4.3}{[1][95][]95}}
\newlabel{bdc:sec:conventions}{{4.2.3}{95}{Two conventions for rupture}{subsection.4.2.3}{}}
\newlabel{bdc:sec:conventions@cref}{{[subsection][3][4,2]4.2.3}{[1][95][]95}}
\newlabel{bdc:fig:killing}{{4.4}{96}{Two conventions for rupture. Under killing (left), rupture from state $n$ ends at $0$ and is indistinguishable from internal extinction. Under catastrophe (right), used here, the same transition ends at the separate state $H$. Birth and death rates remain $\lambda n$ and $\mu n$}{figure.4.4}{}}
\newlabel{bdc:fig:killing@cref}{{[figure][4][4]4.4}{[1][95][]96}}
\newlabel{bdc:sec:notation}{{4.2.4}{96}{Notation}{subsection.4.2.4}{}}
\newlabel{bdc:sec:notation@cref}{{[subsection][4][4,2]4.2.4}{[1][96][]96}}
\newlabel{bdc:sec:riccati}{{4.3}{96}{The Riccati equation and its two roots}{section.4.3}{}}
\newlabel{bdc:sec:riccati@cref}{{[section][3][4]4.3}{[1][96][]96}}
\newlabel{bdc:sec:firststep}{{4.3.1}{96}{First-step decomposition}{subsection.4.3.1}{}}
\newlabel{bdc:sec:firststep@cref}{{[subsection][1][4,3]4.3.1}{[1][96][]96}}
\newlabel{bdc:eq:riccati}{{4.4}{97}{First-step decomposition}{equation.4.3.4}{}}
\newlabel{bdc:eq:riccati@cref}{{[equation][4][4]4.4}{[1][97][]97}}
\newlabel{bdc:sec:roots}{{4.3.2}{97}{The two roots}{subsection.4.3.2}{}}
\newlabel{bdc:sec:roots@cref}{{[subsection][2][4,3]4.3.2}{[1][97][]97}}
\newlabel{bdc:lem:roots}{{4.4}{97}{The characteristic quadratic has one root on either side of unity}{theorem.4.4}{}}
\newlabel{bdc:lem:roots@cref}{{[lemma][4][4]4.4}{[1][97][]97}}
\newlabel{bdc:rem:assumptions}{{4.5}{97}{Standing assumptions}{theorem.4.5}{}}
\newlabel{bdc:rem:assumptions@cref}{{[remark][5][4]4.5}{[1][97][]97}}
\newlabel{bdc:eq:riccatifact}{{4.5}{97}{The two roots}{equation.4.3.5}{}}
\newlabel{bdc:eq:riccatifact@cref}{{[equation][5][4]4.5}{[1][97][]97}}
\newlabel{bdc:fig:roots}{{4.5}{98}{The characteristic quadratic and its flow for $\lambda =1$, $\mu =0.2$ and $\delta =0.05$. The roots $b=0.188$ and $a=1.062$ lie on either side of unity because $q(1)=-\delta $. A trajectory beginning at $I(0)=1$ moves towards the stable root $b$. The distances $A=a-1$ and $B=1-b$ recur in the closed forms}{figure.4.5}{}}
\newlabel{bdc:fig:roots@cref}{{[figure][5][4]4.5}{[1][97][]98}}
\newlabel{bdc:eq:identities}{{4.6}{98}{The two roots}{equation.4.3.6}{}}
\newlabel{bdc:eq:identities@cref}{{[equation][6][4]4.6}{[1][98][]98}}
\newlabel{bdc:sec:solveI}{{4.3.3}{98}{Solving for \texorpdfstring {$I(t)$}{I(t)}}{subsection.4.3.3}{}}
\newlabel{bdc:sec:solveI@cref}{{[subsection][3][4,3]4.3.3}{[1][98][]98}}
\newlabel{bdc:eq:ratio}{{4.7}{98}{Solving for \texorpdfstring {$I(t)$}{I(t)}}{equation.4.3.7}{}}
\newlabel{bdc:eq:ratio@cref}{{[equation][7][4]4.7}{[1][98][]98}}
\newlabel{bdc:thm:I}{{4.6}{99}{The non-catastrophe probability}{theorem.4.6}{}}
\newlabel{bdc:thm:I@cref}{{[theorem][6][4]4.6}{[1][98][]99}}
\newlabel{bdc:eq:I}{{4.8}{99}{The non-catastrophe probability}{equation.4.3.8}{}}
\newlabel{bdc:eq:I@cref}{{[equation][8][4]4.8}{[1][98][]99}}
\newlabel{bdc:sec:D}{{4.3.4}{99}{The same equation again: internal extinction}{subsection.4.3.4}{}}
\newlabel{bdc:sec:D@cref}{{[subsection][4][4,3]4.3.4}{[1][99][]99}}
\newlabel{bdc:eq:riccatiD}{{4.9}{99}{The same equation again: internal extinction}{equation.4.3.9}{}}
\newlabel{bdc:eq:riccatiD@cref}{{[equation][9][4]4.9}{[1][99][]99}}
\newlabel{bdc:thm:D}{{4.7}{99}{The internal-extinction probability}{theorem.4.7}{}}
\newlabel{bdc:thm:D@cref}{{[theorem][7][4]4.7}{[1][99][]99}}
\newlabel{bdc:eq:D}{{4.10}{99}{The internal-extinction probability}{equation.4.3.10}{}}
\newlabel{bdc:eq:D@cref}{{[equation][10][4]4.10}{[1][99][]99}}
\newlabel{bdc:sec:Ifix}{{4.3.5}{99}{The non-fixation probability}{subsection.4.3.5}{}}
\newlabel{bdc:sec:Ifix@cref}{{[subsection][5][4,3]4.3.5}{[1][99][]99}}
\newlabel{bdc:eq:ddIfix}{{4.11}{100}{The non-fixation probability}{equation.4.3.11}{}}
\newlabel{bdc:eq:ddIfix@cref}{{[equation][11][4]4.11}{[1][99][]100}}
\newlabel{bdc:thm:Ifix}{{4.8}{100}{The non-fixation probability}{theorem.4.8}{}}
\newlabel{bdc:thm:Ifix@cref}{{[theorem][8][4]4.8}{[1][100][]100}}
\newlabel{bdc:eq:Ifix}{{4.12}{100}{The non-fixation probability}{equation.4.3.12}{}}
\newlabel{bdc:eq:Ifix@cref}{{[equation][12][4]4.12}{[1][100][]100}}
\newlabel{bdc:fig:fixation}{{4.6}{100}{Non-catastrophe, internal-extinction, and productive probabilities for $\lambda =1$, $\mu =0.2$ and $\delta =0.05$. The non-catastrophe probability $I(t)$ falls to $b$ while the internal-extinction probability $D(t)$ rises to it. Their difference $\Ifix (t)$ is the probability of continued productivity and tends to zero}{figure.4.6}{}}
\newlabel{bdc:fig:fixation@cref}{{[figure][6][4]4.6}{[1][100][]100}}
\newlabel{bdc:sec:multifounder}{{4.3.6}{100}{Several founders}{subsection.4.3.6}{}}
\newlabel{bdc:sec:multifounder@cref}{{[subsection][6][4,3]4.3.6}{[1][100][]100}}
\newlabel{bdc:eq:Ifixk}{{4.13}{101}{Several founders}{equation.4.3.13}{}}
\newlabel{bdc:eq:Ifixk@cref}{{[equation][13][4]4.13}{[1][100][]101}}
\newlabel{bdc:rem:multifounder}{{4.9}{101}{Non-fixation does not factorise}{theorem.4.9}{}}
\newlabel{bdc:rem:multifounder@cref}{{[remark][9][4]4.9}{[1][100][]101}}
\newlabel{bdc:fig:multifounder}{{4.7}{101}{Non-fixation from several founders. (a) Exact probabilities $\Ifix [,{k}](t)=I(t)^k-D(t)^k$ for $k=1,2,3,5$. (b) For two founders, $[I(t)-D(t)]^2$ lies below the exact result because it omits mixed configurations with one extinct and one productive lineage. $\lambda =1$, $\mu =0.2$ and $\delta =0.05$}{figure.4.7}{}}
\newlabel{bdc:fig:multifounder@cref}{{[figure][7][4]4.7}{[1][101][]101}}
\newlabel{bdc:sec:moments}{{4.4}{101}{Moments}{section.4.4}{}}
\newlabel{bdc:sec:moments@cref}{{[section][4][4]4.4}{[1][101][]101}}
\newlabel{bdc:rem:closure}{{4.10}{101}{Why closed forms exist}{theorem.4.10}{}}
\newlabel{bdc:rem:closure@cref}{{[remark][10][4]4.10}{[1][101][]101}}
\newlabel{bdc:sec:J}{{4.4.1}{102}{The mean load}{subsection.4.4.1}{}}
\newlabel{bdc:sec:J@cref}{{[subsection][1][4,4]4.4.1}{[1][102][]102}}
\newlabel{bdc:eq:IJ}{{4.14}{102}{The mean load}{equation.4.4.14}{}}
\newlabel{bdc:eq:IJ@cref}{{[equation][14][4]4.14}{[1][102][]102}}
\newlabel{bdc:thm:J}{{4.11}{102}{The mean load}{theorem.4.11}{}}
\newlabel{bdc:thm:J@cref}{{[theorem][11][4]4.11}{[1][102][]102}}
\newlabel{bdc:eq:J}{{4.15}{102}{The mean load}{equation.4.4.15}{}}
\newlabel{bdc:eq:J@cref}{{[equation][15][4]4.15}{[1][102][]102}}
\newlabel{bdc:eq:ddJ}{{4.16}{102}{The mean load}{equation.4.4.16}{}}
\newlabel{bdc:eq:ddJ@cref}{{[equation][16][4]4.16}{[1][102][]102}}
\newlabel{bdc:sec:K}{{4.4.2}{102}{The second moment}{subsection.4.4.2}{}}
\newlabel{bdc:sec:K@cref}{{[subsection][2][4,4]4.4.2}{[1][102][]102}}
\newlabel{bdc:eq:Kder1}{{4.17}{102}{The second moment}{equation.4.4.17}{}}
\newlabel{bdc:eq:Kder1@cref}{{[equation][17][4]4.17}{[1][102][]102}}
\newlabel{bdc:eq:Kder2}{{4.18}{103}{The second moment}{equation.4.4.18}{}}
\newlabel{bdc:eq:Kder2@cref}{{[equation][18][4]4.18}{[1][102][]103}}
\newlabel{bdc:thm:K}{{4.12}{103}{The second moment of the load}{theorem.4.12}{}}
\newlabel{bdc:thm:K@cref}{{[theorem][12][4]4.12}{[1][103][]103}}
\newlabel{bdc:eq:K}{{4.19}{103}{The second moment of the load}{equation.4.4.19}{}}
\newlabel{bdc:eq:K@cref}{{[equation][19][4]4.19}{[1][103][]103}}
\newlabel{bdc:eq:KI}{{4.20}{103}{The second moment of the load}{equation.4.4.20}{}}
\newlabel{bdc:eq:KI@cref}{{[equation][20][4]4.20}{[1][103][]103}}
\newlabel{bdc:sec:V}{{4.4.3}{103}{The mean release}{subsection.4.4.3}{}}
\newlabel{bdc:sec:V@cref}{{[subsection][3][4,4]4.4.3}{[1][103][]103}}
\newlabel{bdc:eq:ddV}{{4.21}{103}{The mean release}{equation.4.4.21}{}}
\newlabel{bdc:eq:ddV@cref}{{[equation][21][4]4.21}{[1][103][]103}}
\newlabel{bdc:thm:V}{{4.13}{103}{The mean release}{theorem.4.13}{}}
\newlabel{bdc:thm:V@cref}{{[theorem][13][4]4.13}{[1][103][]103}}
\newlabel{bdc:eq:V}{{4.22}{103}{The mean release}{equation.4.4.22}{}}
\newlabel{bdc:eq:V@cref}{{[equation][22][4]4.22}{[1][103][]103}}
\newlabel{bdc:eq:Vinf}{{4.23}{103}{The mean release}{equation.4.4.23}{}}
\newlabel{bdc:eq:Vinf@cref}{{[equation][23][4]4.23}{[1][103][]103}}
\newlabel{bdc:sec:JV}{{4.4.4}{104}{Load and release together}{subsection.4.4.4}{}}
\newlabel{bdc:sec:JV@cref}{{[subsection][4][4,4]4.4.4}{[1][103][]104}}
\newlabel{bdc:fig:jv}{{4.8}{104}{Mean intracellular load and release for $\lambda =1$, $\mu =0.2$ and $\delta =0.05$. (a) The time courses $J(t)$ and $V(t)$: the load first grows and then disappears, while release increases to $V_\infty =13.99$. (b) The same solution as the parametric path $(J(t),V(t))$, with selected times marked. \explorerlink {https://markov.volvox.site/-/adam74aldridge/birth-death-catastrophe-closed-moment-system}{Simulate}}{figure.4.8}{}}
\newlabel{bdc:fig:jv@cref}{{[figure][8][4]4.8}{[1][104][]104}}
\newlabel{bdc:fig:vt}{{4.9}{104}{The same solution read against the no-catastrophe probability, for $\lambda =1$, $\mu =0.2$ and $\delta =0.05$: the mean release $V(t)=\ex {W_t}$ climbs onto its plateau $V_\infty =aB/A=13.99$ (dashed), fed by the hump of the mean load $J(t)$; below, $I(t)$ falls from $1$ to its floor $b=0.188$ (dashed)}{figure.4.9}{}}
\newlabel{bdc:fig:vt@cref}{{[figure][9][4]4.9}{[1][104][]104}}
\newlabel{bdc:thm:Wsq}{{4.14}{105}{The second moment of the release}{theorem.4.14}{}}
\newlabel{bdc:thm:Wsq@cref}{{[theorem][14][4]4.14}{[1][104][]105}}
\newlabel{bdc:eq:Wsq}{{4.24}{105}{The second moment of the release}{equation.4.4.24}{}}
\newlabel{bdc:eq:Wsq@cref}{{[equation][24][4]4.24}{[1][105][]105}}
\newlabel{bdc:sec:variances}{{4.4.5}{105}{Variances}{subsection.4.4.5}{}}
\newlabel{bdc:sec:variances@cref}{{[subsection][5][4,4]4.4.5}{[1][105][]105}}
\newlabel{bdc:eq:VarX}{{4.25}{105}{Variances}{equation.4.4.25}{}}
\newlabel{bdc:eq:VarX@cref}{{[equation][25][4]4.25}{[1][105][]105}}
\newlabel{bdc:eq:VarW}{{4.26}{105}{Variances}{equation.4.4.26}{}}
\newlabel{bdc:eq:VarW@cref}{{[equation][26][4]4.26}{[1][105][]105}}
\newlabel{bdc:fig:meansd}{{4.10}{105}{Means and one-standard-deviation envelopes for (a) the intracellular load $X_t$ and (b) the released count $W_t$, with $\lambda =1$, $\mu =0.2$ and $\delta =0.05$. Lower envelopes are clipped at the support boundary zero}{figure.4.10}{}}
\newlabel{bdc:fig:meansd@cref}{{[figure][10][4]4.10}{[1][105][]105}}
\newlabel{bdc:sec:pair}{{4.4.6}{105}{The pair \texorpdfstring {$(X_t,W_t)$}{(Xt,Wt)}}{subsection.4.4.6}{}}
\newlabel{bdc:sec:pair@cref}{{[subsection][6][4,4]4.4.6}{[1][105][]105}}
\newlabel{bdc:cor:exclusion}{{4.15}{105}{Exclusion, and the covariance}{theorem.4.15}{}}
\newlabel{bdc:cor:exclusion@cref}{{[corollary][15][4]4.15}{[1][105][]105}}
\newlabel{bdc:eq:cov}{{4.27}{105}{Exclusion, and the covariance}{equation.4.4.27}{}}
\newlabel{bdc:eq:cov@cref}{{[equation][27][4]4.27}{[1][105][]105}}
\newlabel{bdc:cor:yield}{{4.16}{106}{Current internal and released yield}{theorem.4.16}{}}
\newlabel{bdc:cor:yield@cref}{{[corollary][16][4]4.16}{[1][105][]106}}
\newlabel{bdc:eq:yield}{{4.28}{106}{Current internal and released yield}{equation.4.4.28}{}}
\newlabel{bdc:eq:yield@cref}{{[equation][28][4]4.28}{[1][105][]106}}
\newlabel{bdc:sec:economy}{{4.5}{106}{The \texorpdfstring {$I$}{I}-economy}{section.4.5}{}}
\newlabel{bdc:sec:economy@cref}{{[section][5][4]4.5}{[1][106][]106}}
\newlabel{bdc:prop:polynomial}{{4.17}{106}{Every load moment is a polynomial in $I$}{theorem.4.17}{}}
\newlabel{bdc:prop:polynomial@cref}{{[proposition][17][4]4.17}{[1][106][]106}}
\newlabel{bdc:eq:VarXI}{{4.29}{107}{The \texorpdfstring {$I$}{I}-economy}{equation.4.5.29}{}}
\newlabel{bdc:eq:VarXI@cref}{{[equation][29][4]4.29}{[1][106][]107}}
\newlabel{bdc:fig:economy}{{4.11}{107}{The mean load $J$, mean release $V$ and second load moment $K$ as functions of $I$, for $\lambda =1$, $\mu =0.2$ and $\delta =0.05$. The trajectory runs from $I=1$ towards $b$; the symmetric-logarithmic ordinate keeps the three scales visible without changing their values}{figure.4.11}{}}
\newlabel{bdc:fig:economy@cref}{{[figure][11][4]4.11}{[1][107][]107}}
\newlabel{bdc:sec:burst}{{4.6}{107}{Burst size and burst time}{section.4.6}{}}
\newlabel{bdc:sec:burst@cref}{{[section][6][4]4.6}{[1][107][]107}}
\newlabel{bdc:sec:conditioning}{{4.6.1}{108}{Conditioning on rupture}{subsection.4.6.1}{}}
\newlabel{bdc:sec:conditioning@cref}{{[subsection][1][4,6]4.6.1}{[1][107][]108}}
\newlabel{bdc:cor:burstmean}{{4.18}{108}{The mean burst size}{theorem.4.18}{}}
\newlabel{bdc:cor:burstmean@cref}{{[corollary][18][4]4.18}{[1][108][]108}}
\newlabel{bdc:eq:burstmean}{{4.30}{108}{The mean burst size}{equation.4.6.30}{}}
\newlabel{bdc:eq:burstmean@cref}{{[equation][30][4]4.30}{[1][108][]108}}
\newlabel{bdc:cor:burstvar}{{4.19}{108}{The first two moments of the burst law}{theorem.4.19}{}}
\newlabel{bdc:cor:burstvar@cref}{{[corollary][19][4]4.19}{[1][108][]108}}
\newlabel{bdc:eq:burstvar}{{4.31}{108}{The first two moments of the burst law}{equation.4.6.31}{}}
\newlabel{bdc:eq:burstvar@cref}{{[equation][31][4]4.31}{[1][108][]108}}
\newlabel{bdc:sec:bursttime}{{4.6.2}{108}{The burst time}{subsection.4.6.2}{}}
\newlabel{bdc:sec:bursttime@cref}{{[subsection][2][4,6]4.6.2}{[1][108][]108}}
\newlabel{bdc:fig:burstsize}{{4.12}{109}{The burst-size law in the no-death case: $\pr {\Kb =k}=(\delta /\lambda )a^{-k}$, a geometric law on $\{1,2,\ldots \}$ with ratio $1/a$. Here $\lambda =1$, $\mu =0$, $\delta =0.1$ and $a=1.1$}{figure.4.12}{}}
\newlabel{bdc:fig:burstsize@cref}{{[figure][12][4]4.12}{[1][108][]109}}
\newlabel{bdc:eq:bursttime}{{4.32}{109}{The burst time}{equation.4.6.32}{}}
\newlabel{bdc:eq:bursttime@cref}{{[equation][32][4]4.32}{[1][108][]109}}
\newlabel{bdc:sec:plateau}{{4.6.3}{109}{The conditional mean and its plateau}{subsection.4.6.3}{}}
\newlabel{bdc:sec:plateau@cref}{{[subsection][3][4,6]4.6.3}{[1][109][]109}}
\newlabel{bdc:fig:condmean}{{4.13}{110}{Conditional mean loads. (a) For $\mu =0$, $\lambda =1$ and $\delta =0.1$, $J(t)/I(t)$ and $V(t)$ tend to $1+\lambda /\delta =11$. (b) For $\lambda =1$, $\mu =0.2$ and $\delta =0.05$, $J(t)/\Ifix (t)$ approaches $a/(a-1)$, whereas $J(t)/I(t)$ tends to zero because non-catastrophe also includes cells that died internally}{figure.4.13}{}}
\newlabel{bdc:fig:condmean@cref}{{[figure][13][4]4.13}{[1][109][]110}}
\newlabel{bdc:sec:summary}{{4.7}{110}{Summary of results}{section.4.7}{}}
\newlabel{bdc:sec:summary@cref}{{[section][7][4]4.7}{[1][110][]110}}
\newlabel{bdc:fig:landscape}{{4.14}{110}{Rate landscapes at $\lambda =1$, for $\mu \in [0,1]$ and $\delta \in [0.01,1]$. (a) The probability $b$ of internal extinction before rupture. (b) The mean release $V_\infty $. The star marks $(\mu ,\delta )=(0.2,0.05)$, used in most figures; the diamond marks the no-death point $(0,0.1)$. The catastrophe-rate axis is logarithmic}{figure.4.14}{}}
\newlabel{bdc:fig:landscape@cref}{{[figure][14][4]4.14}{[1][110][]110}}
\newlabel{bdc:tab:notation}{{4.1}{111}{Notation for the birth--death--catastrophe process}{table.4.1}{}}
\newlabel{bdc:tab:notation@cref}{{[table][1][4]4.1}{[1][110][]111}}
\newlabel{bdc:tab:closedforms}{{4.2}{111}{Principal closed forms, from one founder}{table.4.2}{}}
\newlabel{bdc:tab:closedforms@cref}{{[table][2][4]4.2}{[1][110][]111}}
\newlabel{bdc:tab:odes}{{4.3}{112}{The four differential relations underlying the closed forms}{table.4.3}{}}
\newlabel{bdc:tab:odes@cref}{{[table][3][4]4.3}{[1][110][]112}}
\newlabel{bdc:sec:next}{{4.8}{112}{Conclusion}{section.4.8}{}}
\newlabel{bdc:sec:next@cref}{{[section][8][4]4.8}{[1][112][]112}}
\newlabel{bdc:sec:next-singlecell}{{4.8.1}{112}{Completing the single-cell theory}{subsection.4.8.1}{}}
\newlabel{bdc:sec:next-singlecell@cref}{{[subsection][1][4,8]4.8.1}{[1][112][]112}}
\newlabel{bdc:sec:next-population}{{4.8.2}{112}{Embedding the cell in a population}{subsection.4.8.2}{}}
\newlabel{bdc:sec:next-population@cref}{{[subsection][2][4,8]4.8.2}{[1][112][]112}}
\newlabel{bdc:app:limits}{{4.A}{113}{Checks and limiting cases}{section.4.appendix.A}{}}
\newlabel{bdc:app:limits@cref}{{[section][1][4]4.A}{[1][113][]113}}
\newlabel{bdc:app:initial}{{4.A.1}{113}{Initial values}{subsection.4.appendix.A.1}{}}
\newlabel{bdc:app:initial@cref}{{[subsection][1][4,1]4.A.1}{[1][113][]113}}
\newlabel{bdc:app:late}{{4.A.2}{113}{Late-time values}{subsection.4.appendix.A.2}{}}
\newlabel{bdc:app:late@cref}{{[subsection][2][4,1]4.A.2}{[1][113][]113}}
\newlabel{bdc:app:nodeath}{{4.A.3}{113}{No death}{subsection.4.appendix.A.3}{}}
\newlabel{bdc:app:nodeath@cref}{{[subsection][3][4,1]4.A.3}{[1][113][]113}}
\newlabel{bdc:app:nocatastrophe}{{4.A.4}{114}{No catastrophe}{subsection.4.appendix.A.4}{}}
\newlabel{bdc:app:nocatastrophe@cref}{{[subsection][4][4,1]4.A.4}{[1][113][]114}}
\newlabel{bdc:app:critical}{{4.A.5}{114}{Critical births and deaths}{subsection.4.appendix.A.5}{}}
\newlabel{bdc:app:critical@cref}{{[subsection][5][4,1]4.A.5}{[1][114][]114}}
\newlabel{bdc:app:nobirths}{{4.A.6}{114}{No births}{subsection.4.appendix.A.6}{}}
\newlabel{bdc:app:nobirths@cref}{{[subsection][6][4,1]4.A.6}{[1][114][]114}}
\newlabel{bdc:app:Wsq}{{4.A.7}{114}{The second moment of the release}{subsection.4.appendix.A.7}{}}
\newlabel{bdc:app:Wsq@cref}{{[subsection][7][4,1]4.A.7}{[1][114][]114}}
\newlabel{bdc:eq:ddWsq}{{4.33}{114}{The second moment of the release}{equation.4.appendix.A.33}{}}
\newlabel{bdc:eq:ddWsq@cref}{{[equation][33][4]4.33}{[1][114][]114}}
\newlabel{bdc:tab:xsqchanges}{{4.4}{115}{Changes in $X_t^2$ under the three event types}{table.4.4}{}}
\newlabel{bdc:tab:xsqchanges@cref}{{[table][4][4]4.4}{[1][114][]115}}
\newlabel{bdc:eq:ddK}{{4.34}{115}{The second moment of the release}{equation.4.appendix.A.34}{}}
\newlabel{bdc:eq:ddK@cref}{{[equation][34][4]4.34}{[1][114][]115}}
\newlabel{bdc:eq:dWsq}{{4.35}{115}{The second moment of the release}{equation.4.appendix.A.35}{}}
\newlabel{bdc:eq:dWsq@cref}{{[equation][35][4]4.35}{[1][114][]115}}
\newlabel{dist:ch:distribution-theory}{{5}{116}{The birth--death--catastrophe process: distribution theory, quasi-stationarity, and burst statistics}{chapter.5}{}}
\newlabel{dist:ch:distribution-theory@cref}{{[chapter][5][]5}{[1][116][]116}}
\newlabel{dist:sec:intro}{{5.1}{116}{Introduction}{section.5.1}{}}
\newlabel{dist:sec:intro@cref}{{[section][1][5]5.1}{[1][116][]116}}
\newlabel{dist:sec:assumptions}{{5.1.1}{117}{Standing assumptions}{subsection.5.1.1}{}}
\newlabel{dist:sec:assumptions@cref}{{[subsection][1][5,1]5.1.1}{[1][117][]117}}
\newlabel{dist:sec:contributions}{{5.1.2}{117}{Contributions}{subsection.5.1.2}{}}
\newlabel{dist:sec:contributions@cref}{{[subsection][2][5,1]5.1.2}{[1][117][]117}}
\newlabel{dist:sec:exports}{{5.1.3}{118}{What this chapter exports}{subsection.5.1.3}{}}
\newlabel{dist:sec:exports@cref}{{[subsection][3][5,1]5.1.3}{[1][118][]118}}
\newlabel{dist:sec:recap}{{5.2}{118}{The process, its roots, and its closed forms}{section.5.2}{}}
\newlabel{dist:sec:recap@cref}{{[section][2][5]5.2}{[1][118][]118}}
\newlabel{dist:sec:notation-tab}{{5.2.1}{118}{Canonical notation}{subsection.5.2.1}{}}
\newlabel{dist:sec:notation-tab@cref}{{[subsection][1][5,2]5.2.1}{[1][118][]118}}
\newlabel{dist:sec:process}{{5.2.2}{118}{The process and its release counter}{subsection.5.2.2}{}}
\newlabel{dist:sec:process@cref}{{[subsection][2][5,2]5.2.2}{[1][118][]118}}
\newlabel{dist:def:bdc}{{5.1}{118}{Birth--death--catastrophe process}{theorem.5.1}{}}
\newlabel{dist:def:bdc@cref}{{[definition][1][5]5.1}{[1][118][]118}}
\newlabel{dist:tab:notation}{{5.1}{119}{Canonical notation for this chapter}{table.5.1}{}}
\newlabel{dist:tab:notation@cref}{{[table][1][5]5.1}{[1][118][]119}}
\newlabel{dist:eq:wdef}{{5.1}{119}{The process and its release counter}{equation.5.2.1}{}}
\newlabel{dist:eq:wdef@cref}{{[equation][1][5]5.1}{[1][119][]119}}
\newlabel{dist:sec:roots}{{5.2.3}{119}{Roots and the workhorse identity}{subsection.5.2.3}{}}
\newlabel{dist:sec:roots@cref}{{[subsection][3][5,2]5.2.3}{[1][119][]119}}
\newlabel{dist:eq:roots}{{5.2}{119}{Roots and the workhorse identity}{equation.5.2.2}{}}
\newlabel{dist:eq:roots@cref}{{[equation][2][5]5.2}{[1][119][]119}}
\newlabel{dist:fig:joint}{{5.1}{120}{A representative path of the joint process $(\Xt ,\Wt )$. The intracellular load follows a birth--death path until rupture at time $\tau $, when $\Xt $ enters the absorbing state $H$. At the same instant the released count jumps from zero to the pre-rupture load $X_{\tau ^-}$ and remains there. The two axes share a time origin: nothing leaves the cell before $\tau $, and everything leaves at $\tau $}{figure.5.1}{}}
\newlabel{dist:fig:joint@cref}{{[figure][1][5]5.1}{[1][119][]120}}
\newlabel{dist:eq:abtheta}{{5.3}{120}{Roots and the workhorse identity}{equation.5.2.3}{}}
\newlabel{dist:eq:abtheta@cref}{{[equation][3][5]5.3}{[1][119][]120}}
\newlabel{dist:eq:ab}{{5.4}{120}{Roots and the workhorse identity}{equation.5.2.4}{}}
\newlabel{dist:eq:ab@cref}{{[equation][4][5]5.4}{[1][120][]120}}
\newlabel{dist:sec:moments-recap}{{5.2.4}{120}{The closed forms}{subsection.5.2.4}{}}
\newlabel{dist:sec:moments-recap@cref}{{[subsection][4][5,2]5.2.4}{[1][120][]120}}
\newlabel{dist:def:fixation}{{5.2}{120}{Fixation functions}{theorem.5.2}{}}
\newlabel{dist:def:fixation@cref}{{[definition][2][5]5.2}{[1][120][]120}}
\newlabel{dist:eq:idefs}{{5.5}{120}{Fixation functions}{equation.5.2.5}{}}
\newlabel{dist:eq:idefs@cref}{{[equation][5][5]5.5}{[1][120][]120}}
\newlabel{dist:eq:ifixdef}{{5.6}{120}{Fixation functions}{equation.5.2.6}{}}
\newlabel{dist:eq:ifixdef@cref}{{[equation][6][5]5.6}{[1][120][]120}}
\newlabel{dist:eq:riccati}{{5.7}{120}{The closed forms}{equation.5.2.7}{}}
\newlabel{dist:eq:riccati@cref}{{[equation][7][5]5.7}{[1][120][]120}}
\newlabel{dist:eq:idihat}{{5.8}{121}{The closed forms}{equation.5.2.8}{}}
\newlabel{dist:eq:idihat@cref}{{[equation][8][5]5.8}{[1][120][]121}}
\newlabel{dist:eq:j}{{5.9}{121}{The closed forms}{equation.5.2.9}{}}
\newlabel{dist:eq:j@cref}{{[equation][9][5]5.9}{[1][121][]121}}
\newlabel{dist:eq:k}{{5.10}{121}{The closed forms}{equation.5.2.10}{}}
\newlabel{dist:eq:k@cref}{{[equation][10][5]5.10}{[1][121][]121}}
\newlabel{dist:eq:v}{{5.11}{121}{The closed forms}{equation.5.2.11}{}}
\newlabel{dist:eq:v@cref}{{[equation][11][5]5.11}{[1][121][]121}}
\newlabel{dist:eq:vinf}{{5.12}{121}{The closed forms}{equation.5.2.12}{}}
\newlabel{dist:eq:vinf@cref}{{[equation][12][5]5.12}{[1][121][]121}}
\newlabel{dist:eq:ew2}{{5.13}{121}{The closed forms}{equation.5.2.13}{}}
\newlabel{dist:eq:ew2@cref}{{[equation][13][5]5.13}{[1][121][]121}}
\newlabel{dist:sec:pgf}{{5.3}{121}{The generating function}{section.5.3}{}}
\newlabel{dist:sec:pgf@cref}{{[section][3][5]5.3}{[1][121][]121}}
\newlabel{dist:eq:gdef}{{5.14}{121}{The generating function}{equation.5.3.14}{}}
\newlabel{dist:eq:gdef@cref}{{[equation][14][5]5.14}{[1][121][]121}}
\newlabel{dist:sec:masterpde}{{5.3.1}{122}{From the master equations to a PDE}{subsection.5.3.1}{}}
\newlabel{dist:sec:masterpde@cref}{{[subsection][1][5,3]5.3.1}{[1][121][]122}}
\newlabel{dist:eq:kolmogorov}{{5.15}{122}{From the master equations to a PDE}{equation.5.3.15}{}}
\newlabel{dist:eq:kolmogorov@cref}{{[equation][15][5]5.15}{[1][121][]122}}
\newlabel{dist:eq:kolmogorovcat}{{5.16}{122}{From the master equations to a PDE}{equation.5.3.16}{}}
\newlabel{dist:eq:kolmogorovcat@cref}{{[equation][16][5]5.16}{[1][122][]122}}
\newlabel{dist:eq:killpde}{{5.17}{122}{From the master equations to a PDE}{equation.5.3.17}{}}
\newlabel{dist:eq:killpde@cref}{{[equation][17][5]5.17}{[1][122][]122}}
\newlabel{dist:eq:catpde}{{5.18}{122}{From the master equations to a PDE}{equation.5.3.18}{}}
\newlabel{dist:eq:catpde@cref}{{[equation][18][5]5.18}{[1][122][]122}}
\newlabel{dist:sec:characteristics}{{5.3.2}{122}{Solving by characteristics}{subsection.5.3.2}{}}
\newlabel{dist:sec:characteristics@cref}{{[subsection][2][5,3]5.3.2}{[1][122][]122}}
\newlabel{dist:eq:zchar}{{5.19}{122}{Solving by characteristics}{equation.5.3.19}{}}
\newlabel{dist:eq:zchar@cref}{{[equation][19][5]5.19}{[1][122][]122}}
\newlabel{dist:eq:zsolve}{{5.20}{122}{Solving by characteristics}{equation.5.3.20}{}}
\newlabel{dist:eq:zsolve@cref}{{[equation][20][5]5.20}{[1][122][]122}}
\newlabel{dist:eq:sinv}{{5.21}{122}{Solving by characteristics}{equation.5.3.21}{}}
\newlabel{dist:eq:sinv@cref}{{[equation][21][5]5.21}{[1][122][]122}}
\newlabel{dist:eq:gcat}{{5.22}{122}{Solving by characteristics}{equation.5.3.22}{}}
\newlabel{dist:eq:gcat@cref}{{[equation][22][5]5.22}{[1][122][]122}}
\newlabel{dist:sec:pgf-k-founders}{{5.3.3}{123}{Several founders}{subsection.5.3.3}{}}
\newlabel{dist:sec:pgf-k-founders@cref}{{[subsection][3][5,3]5.3.3}{[1][122][]123}}
\newlabel{dist:eq:gk}{{5.23}{123}{Several founders}{equation.5.3.23}{}}
\newlabel{dist:eq:gk@cref}{{[equation][23][5]5.23}{[1][123][]123}}
\newlabel{dist:sec:conventions}{{5.3.4}{123}{The two conventions differ by a function of time}{subsection.5.3.4}{}}
\newlabel{dist:sec:conventions@cref}{{[subsection][4][5,3]5.3.4}{[1][123][]123}}
\newlabel{dist:prop:conventions}{{5.3}{123}{The two rupture conventions}{theorem.5.3}{}}
\newlabel{dist:prop:conventions@cref}{{[proposition][3][5]5.3}{[1][123][]123}}
\newlabel{dist:eq:conventions}{{5.24}{123}{The two rupture conventions}{equation.5.3.24}{}}
\newlabel{dist:eq:conventions@cref}{{[equation][24][5]5.24}{[1][123][]123}}
\newlabel{dist:eq:gkill}{{5.25}{123}{The two conventions differ by a function of time}{equation.5.3.25}{}}
\newlabel{dist:eq:gkill@cref}{{[equation][25][5]5.25}{[1][123][]123}}
\newlabel{dist:sec:geometric}{{5.4}{123}{The load is geometric at every time}{section.5.4}{}}
\newlabel{dist:sec:geometric@cref}{{[section][4][5]5.4}{[1][123][]123}}
\newlabel{dist:eq:dnG}{{5.26}{123}{The load is geometric at every time}{equation.5.4.26}{}}
\newlabel{dist:eq:dnG@cref}{{[equation][26][5]5.26}{[1][123][]123}}
\newlabel{dist:eq:stateprob}{{5.27}{124}{The load is geometric at every time}{equation.5.4.27}{}}
\newlabel{dist:eq:stateprob@cref}{{[equation][27][5]5.27}{[1][123][]124}}
\newlabel{dist:sec:sliding}{{5.4.1}{124}{The sliding ratio}{subsection.5.4.1}{}}
\newlabel{dist:sec:sliding@cref}{{[subsection][1][5,4]5.4.1}{[1][124][]124}}
\newlabel{dist:eq:pdef}{{5.28}{124}{The sliding ratio}{equation.5.4.28}{}}
\newlabel{dist:eq:pdef@cref}{{[equation][28][5]5.28}{[1][124][]124}}
\newlabel{dist:thm:geometric}{{5.4}{124}{The load is geometric at every time}{theorem.5.4}{}}
\newlabel{dist:thm:geometric@cref}{{[theorem][4][5]5.4}{[1][124][]124}}
\newlabel{dist:eq:geometric}{{5.29}{124}{The load is geometric at every time}{equation.5.4.29}{}}
\newlabel{dist:eq:geometric@cref}{{[equation][29][5]5.29}{[1][124][]124}}
\newlabel{dist:eq:plimits}{{5.30}{124}{The load is geometric at every time}{equation.5.4.30}{}}
\newlabel{dist:eq:plimits@cref}{{[equation][30][5]5.30}{[1][124][]124}}
\newlabel{dist:sec:sum-ihat}{{5.4.2}{124}{Normalisation: the sum is \texorpdfstring {$\Ifix $}{Ifix}}{subsection.5.4.2}{}}
\newlabel{dist:sec:sum-ihat@cref}{{[subsection][2][5,4]5.4.2}{[1][124][]124}}
\newlabel{dist:cor:sum}{{5.5}{124}{The positive states carry the non-fixation probability}{theorem.5.5}{}}
\newlabel{dist:cor:sum@cref}{{[corollary][5][5]5.5}{[1][124][]124}}
\newlabel{dist:eq:pnsum}{{5.31}{124}{The positive states carry the non-fixation probability}{equation.5.4.31}{}}
\newlabel{dist:eq:pnsum@cref}{{[equation][31][5]5.31}{[1][124][]124}}
\newlabel{dist:fig:readings}{{5.2}{125}{One endpoint, three readings. The ratio $P(t)$ of \cref {dist:thm:geometric} slides from $0$ at founding to $1/a$, and each of the next three sections applies a different operation to that endpoint: \cref {dist:sec:qsd} takes it as the quasi-stationary ratio, \cref {dist:sec:burst-size} integrates the rupture hazard against it, and \cref {dist:sec:burst-time} conditions the rupture time on the size it produces. The first two return the same law, differing only by the eventual-burst probability $1-b$; that coincidence is \cref {dist:thm:identity}}{figure.5.2}{}}
\newlabel{dist:fig:readings@cref}{{[figure][2][5]5.2}{[1][124][]125}}
\newlabel{dist:fig:slide}{{5.3}{125}{State probabilities $p_n(t)$ at successive times for $(\lambda ,\mu ,\delta )=(1,0.2,0.05)$. On the logarithmic vertical scale each geometric law is a straight line whose ratio $P(t)$ increases towards $1/a$; the inset shows that approach directly. The curves fall in $n$ and rise in $t$: the load spreads as $P$ slides}{figure.5.3}{}}
\newlabel{dist:fig:slide@cref}{{[figure][3][5]5.3}{[1][124][]125}}
\newlabel{dist:sec:qsd}{{5.5}{126}{The quasi-stationary distribution}{section.5.5}{}}
\newlabel{dist:sec:qsd@cref}{{[section][5][5]5.5}{[1][126][]126}}
\newlabel{dist:sec:conditioning}{{5.5.1}{126}{Conditioning on survival}{subsection.5.5.1}{}}
\newlabel{dist:sec:conditioning@cref}{{[subsection][1][5,5]5.5.1}{[1][126][]126}}
\newlabel{dist:eq:fixprob}{{5.32}{126}{Conditioning on survival}{equation.5.5.32}{}}
\newlabel{dist:eq:fixprob@cref}{{[equation][32][5]5.32}{[1][126][]126}}
\newlabel{dist:eq:condp}{{5.33}{126}{Conditioning on survival}{equation.5.5.33}{}}
\newlabel{dist:eq:condp@cref}{{[equation][33][5]5.33}{[1][126][]126}}
\newlabel{dist:eq:qsddef}{{5.34}{126}{Conditioning on survival}{equation.5.5.34}{}}
\newlabel{dist:eq:qsddef@cref}{{[equation][34][5]5.34}{[1][126][]126}}
\newlabel{dist:sec:qsd-geometric}{{5.5.2}{126}{The limit is geometric}{subsection.5.5.2}{}}
\newlabel{dist:sec:qsd-geometric@cref}{{[subsection][2][5,5]5.5.2}{[1][126][]126}}
\newlabel{dist:thm:qsd}{{5.6}{126}{Quasi-stationary distribution of the BDC process}{theorem.5.6}{}}
\newlabel{dist:thm:qsd@cref}{{[theorem][6][5]5.6}{[1][126][]126}}
\newlabel{dist:eq:qsd}{{5.35}{126}{Quasi-stationary distribution of the BDC process}{equation.5.5.35}{}}
\newlabel{dist:eq:qsd@cref}{{[equation][35][5]5.35}{[1][126][]126}}
\newlabel{dist:fig:qsmeans}{{5.4}{127}{Quasi-stationary and release means across internal death rates at $\lambda =1$ and $\delta =0.05$. The quasi-stationary mean $\langle X\rangle _{\rm QS}=a/(a-1)$ is identically the conditional burst mean $V_\infty /(1-b)$. At $\mu =0$, $b=0$ and both also equal $V_\infty $; once $\mu >0$, internally cleared cells contribute zero release and place $V_\infty $ strictly below the other two; the arrow marks that gap, $b\langle X\rangle _{\rm QS}=3.25$, at the working point. In the $\mu =0$ case the conditional mean and the cumulative release share a limit, $\lim _t J(t)/I(t)=\lim _t V(t)=1+\lambda /\delta $, which diverges as $\delta $ shrinks}{figure.5.4}{}}
\newlabel{dist:fig:qsmeans@cref}{{[figure][4][5]5.4}{[1][127][]127}}
\newlabel{dist:eq:condgeom}{{5.36}{127}{The limit is geometric}{equation.5.5.36}{}}
\newlabel{dist:eq:condgeom@cref}{{[equation][36][5]5.36}{[1][126][]127}}
\newlabel{dist:eq:qsmoments}{{5.37}{127}{The limit is geometric}{equation.5.5.37}{}}
\newlabel{dist:eq:qsmoments@cref}{{[equation][37][5]5.37}{[1][127][]127}}
\newlabel{dist:sec:decay-rate}{{5.5.3}{127}{The decay rate, and what it validates}{subsection.5.5.3}{}}
\newlabel{dist:sec:decay-rate@cref}{{[subsection][3][5,5]5.5.3}{[1][127][]127}}
\newlabel{dist:eq:survivalloss}{{5.38}{127}{The decay rate, and what it validates}{equation.5.5.38}{}}
\newlabel{dist:eq:survivalloss@cref}{{[equation][38][5]5.38}{[1][127][]127}}
\newlabel{dist:prop:decay}{{5.7}{128}{Quasi-stationary decay rate}{theorem.5.7}{}}
\newlabel{dist:prop:decay@cref}{{[proposition][7][5]5.7}{[1][128][]128}}
\newlabel{dist:eq:decayrate}{{5.39}{128}{Quasi-stationary decay rate}{equation.5.5.39}{}}
\newlabel{dist:eq:decayrate@cref}{{[equation][39][5]5.39}{[1][128][]128}}
\newlabel{dist:eq:decaycheck}{{5.40}{128}{The decay rate, and what it validates}{equation.5.5.40}{}}
\newlabel{dist:eq:decaycheck@cref}{{[equation][40][5]5.40}{[1][128][]128}}
\newlabel{dist:sec:lifetime}{{5.5.4}{128}{Mean productive lifetime}{subsection.5.5.4}{}}
\newlabel{dist:sec:lifetime@cref}{{[subsection][4][5,5]5.5.4}{[1][128][]128}}
\newlabel{dist:eq:tprodint}{{5.41}{128}{Mean productive lifetime}{equation.5.5.41}{}}
\newlabel{dist:eq:tprodint@cref}{{[equation][41][5]5.41}{[1][128][]128}}
\newlabel{dist:prop:lifetime}{{5.8}{128}{Mean productive lifetime}{theorem.5.8}{}}
\newlabel{dist:prop:lifetime@cref}{{[proposition][8][5]5.8}{[1][128][]128}}
\newlabel{dist:eq:tprod}{{5.42}{128}{Mean productive lifetime}{equation.5.5.42}{}}
\newlabel{dist:eq:tprod@cref}{{[equation][42][5]5.42}{[1][128][]128}}
\newlabel{dist:fig:lifetime}{{5.5}{129}{Mean productive lifetime $\E {T_{\rm prod}}=\lambda ^{-1}\log \bigl (a/(a-1)\bigr )$. (a) Its dependence on $\mu \in [0,1]$ and $\delta \in [0.005,1]$ at $\lambda =1$; contours give lifetime values and the star marks the working point. (b) The productive-survival curve $\Ifix (t)$ at $(\lambda ,\mu ,\delta )=(1,0.2,0.05)$, with its integral $\E {T_{\rm prod}}=2.847$ shaded. The matched classical death rate for \ChPop {} is $d_{\Icell }=1/\E {T_{\rm prod}}=0.351$}{figure.5.5}{}}
\newlabel{dist:fig:lifetime@cref}{{[figure][5][5]5.5}{[1][129][]129}}
\newlabel{dist:sec:burst-size}{{5.6}{129}{Burst size, and why it is the quasi-stationary law}{section.5.6}{}}
\newlabel{dist:sec:burst-size@cref}{{[section][6][5]5.6}{[1][129][]129}}
\newlabel{dist:sec:burst-law}{{5.6.1}{129}{The burst-size law}{subsection.5.6.1}{}}
\newlabel{dist:sec:burst-law@cref}{{[subsection][1][5,6]5.6.1}{[1][129][]129}}
\newlabel{dist:eq:burstintegral}{{5.43}{129}{The burst-size law}{equation.5.6.43}{}}
\newlabel{dist:eq:burstintegral@cref}{{[equation][43][5]5.43}{[1][129][]129}}
\newlabel{dist:eq:antideriv}{{5.44}{130}{The burst-size law}{equation.5.6.44}{}}
\newlabel{dist:eq:antideriv@cref}{{[equation][44][5]5.44}{[1][130][]130}}
\newlabel{dist:eq:pkintegral}{{5.45}{130}{The burst-size law}{equation.5.6.45}{}}
\newlabel{dist:eq:pkintegral@cref}{{[equation][45][5]5.45}{[1][130][]130}}
\newlabel{dist:thm:burst}{{5.9}{130}{Burst-size distribution}{theorem.5.9}{}}
\newlabel{dist:thm:burst@cref}{{[theorem][9][5]5.9}{[1][130][]130}}
\newlabel{dist:eq:burstlaw}{{5.46}{130}{Burst-size distribution}{equation.5.6.46}{}}
\newlabel{dist:eq:burstlaw@cref}{{[equation][46][5]5.46}{[1][130][]130}}
\newlabel{dist:eq:burstlawcond}{{5.47}{130}{Burst-size distribution}{equation.5.6.47}{}}
\newlabel{dist:eq:burstlawcond@cref}{{[equation][47][5]5.47}{[1][130][]130}}
\newlabel{dist:sec:identity}{{5.6.2}{130}{Why the two laws agree}{subsection.5.6.2}{}}
\newlabel{dist:sec:identity@cref}{{[subsection][2][5,6]5.6.2}{[1][130][]130}}
\newlabel{dist:thm:identity}{{5.10}{130}{The burst size is the quasi-stationary law}{theorem.5.10}{}}
\newlabel{dist:thm:identity@cref}{{[theorem][10][5]5.10}{[1][130][]130}}
\newlabel{dist:eq:ponelaw}{{5.48}{131}{Why the two laws agree}{equation.5.6.48}{}}
\newlabel{dist:eq:ponelaw@cref}{{[equation][48][5]5.48}{[1][130][]131}}
\newlabel{dist:eq:clockchange}{{5.49}{131}{Why the two laws agree}{equation.5.6.49}{}}
\newlabel{dist:eq:clockchange@cref}{{[equation][49][5]5.49}{[1][131][]131}}
\newlabel{dist:eq:telescope}{{5.50}{131}{Why the two laws agree}{equation.5.6.50}{}}
\newlabel{dist:eq:telescope@cref}{{[equation][50][5]5.50}{[1][131][]131}}
\newlabel{dist:rem:criterion}{{5.11}{131}{What the argument needs}{theorem.5.11}{}}
\newlabel{dist:rem:criterion@cref}{{[remark][11][5]5.11}{[1][131][]131}}
\newlabel{dist:sec:size-bias}{{5.6.3}{131}{Size-biasing and the size of a late burst}{subsection.5.6.3}{}}
\newlabel{dist:sec:size-bias@cref}{{[subsection][3][5,6]5.6.3}{[1][131][]131}}
\newlabel{dist:eq:sizebias}{{5.51}{131}{Size-biasing and the size of a late burst}{equation.5.6.51}{}}
\newlabel{dist:eq:sizebias@cref}{{[equation][51][5]5.51}{[1][131][]131}}
\newlabel{dist:eq:latelimit}{{5.52}{131}{Size-biasing and the size of a late burst}{equation.5.6.52}{}}
\newlabel{dist:eq:latelimit@cref}{{[equation][52][5]5.52}{[1][131][]131}}
\newlabel{dist:fig:two-mechanisms}{{5.6}{132}{Two mechanisms, one law. (a) The load conditioned on being neither ruptured nor extinct, $p_n(t)/\Ifix (t)$, at $t=1,5,20$ from \cref {dist:eq:stateprob,dist:eq:idihat}, with the quasi-stationary law $(a-1)a^{-n}$ dashed; on the logarithmic scale each transient law is a straight line whose slope rotates onto the limiting slope as $P(t)$ slides up to $1/a$. (b) The burst-size law $(\delta /\lambda )a^{-k}$ of \cref {dist:eq:burstlaw} (markers) and the same law conditioned on the burst occurring (dashed) --- the identical line, reached without taking a limit at all; the two series differ by the constant factor $1-b=0.812$. The first panel arrives by conditioning on survival along a single trajectory, the second by weighting each state by its rupture hazard and integrating across trajectories. (c) The two routes on one axis, each computed on its own terms: the conditional law of (a) evaluated at $t=60$, against the occupation integral of (b) evaluated by numerical quadrature rather than in closed form. They agree to $10^{-8}$, which is the accuracy of the quadrature grid and not of the identity --- \cref {dist:thm:identity} makes the two exactly equal. Parameters $(\lambda ,\mu ,\delta )=(1,0.2,0.05)$, so $a=1.0616$ and $b=0.1884$; closed forms throughout except where panel (c) integrates numerically, and no simulation anywhere}{figure.5.6}{}}
\newlabel{dist:fig:two-mechanisms@cref}{{[figure][6][5]5.6}{[1][131][]132}}
\newlabel{dist:fig:latemean}{{5.7}{132}{The size of a late burst at $(\lambda ,\mu ,\delta )=(1,0.2,0.05)$. The conditional mean $\E {\Kb \mid \tau =t}=K(t)/J(t)$ of \cref {dist:eq:sizebias} rises from $1$ at founding to $(a+1)/(a-1)=33.46$; the quasi-stationary mean $a/(a-1)=17.23$ is marked for contrast, their ratio being $(a+1)/a=1.94$}{figure.5.7}{}}
\newlabel{dist:fig:latemean@cref}{{[figure][7][5]5.7}{[1][131][]132}}
\newlabel{dist:fig:jointtauk}{{5.8}{133}{Joint density, with respect to time and counting measure, of burst time $\tau $ and burst size $\Kb $ for $(\lambda ,\mu ,\delta )=(1,0,0.1)$. The overlay is the conditional mean $\E {\Kb \mid \tau =t}=K(t)/J(t)$, rising from $1$ towards $1+2\lambda /\delta =21$; the upward tilt of the mass is size-biasing made visible}{figure.5.8}{}}
\newlabel{dist:fig:jointtauk@cref}{{[figure][8][5]5.8}{[1][131][]133}}
\newlabel{dist:sec:burst-time}{{5.7}{133}{Burst time}{section.5.7}{}}
\newlabel{dist:sec:burst-time@cref}{{[section][7][5]5.7}{[1][131][]133}}
\newlabel{dist:sec:burst-time-marginal}{{5.7.1}{133}{The marginal law, and its defect}{subsection.5.7.1}{}}
\newlabel{dist:sec:burst-time-marginal@cref}{{[subsection][1][5,7]5.7.1}{[1][133][]133}}
\newlabel{dist:eq:phi}{{5.53}{133}{The marginal law, and its defect}{equation.5.7.53}{}}
\newlabel{dist:eq:phi@cref}{{[equation][53][5]5.53}{[1][133][]133}}
\newlabel{dist:eq:phidefect}{{5.54}{133}{The marginal law, and its defect}{equation.5.7.54}{}}
\newlabel{dist:eq:phidefect@cref}{{[equation][54][5]5.54}{[1][133][]133}}
\newlabel{dist:sec:tau-given-k}{{5.7.2}{133}{Burst time given burst size, without death}{subsection.5.7.2}{}}
\newlabel{dist:sec:tau-given-k@cref}{{[subsection][2][5,7]5.7.2}{[1][133][]133}}
\newlabel{dist:eq:tn}{{5.55}{133}{Burst time given burst size, without death}{equation.5.7.55}{}}
\newlabel{dist:eq:tn@cref}{{[equation][55][5]5.55}{[1][133][]133}}
\newlabel{dist:fig:bursttime}{{5.9}{134}{Defective burst-time density $\varphi (t)=\delta J(t)$ and its conditional counterpart $\varphi (t)/(1-b)$. For $(\lambda ,\mu ,\delta )=(1,0.2,0.05)$ the finite-time density has total mass $1-b$ and the remaining mass $b=0.1884$ lies at $\tau =\infty $; for $(1,0,0.1)$, $b=0$ and the two densities coincide. The gap between the curves is the whole content of the defect}{figure.5.9}{}}
\newlabel{dist:fig:bursttime@cref}{{[figure][9][5]5.9}{[1][133][]134}}
\newlabel{dist:prop:taugivenk}{{5.12}{134}{Conditional rupture time, $\mu =0$}{theorem.5.12}{}}
\newlabel{dist:prop:taugivenk@cref}{{[proposition][12][5]5.12}{[1][134][]134}}
\newlabel{dist:eq:taugivenk}{{5.56}{134}{Conditional rupture time, $\mu =0$}{equation.5.7.56}{}}
\newlabel{dist:eq:taugivenk@cref}{{[equation][56][5]5.56}{[1][134][]134}}
\newlabel{dist:eq:taugivenkasymp}{{5.57}{134}{Burst time given burst size, without death}{equation.5.7.57}{}}
\newlabel{dist:eq:taugivenkasymp@cref}{{[equation][57][5]5.57}{[1][134][]134}}
\newlabel{dist:fig:taugivenk}{{5.10}{135}{Conditional mean rupture time given burst size $n$ when $\mu =0$, with $\lambda =1$, $\delta =0.1$ and hence $\theta =1.1$. Exact harmonic means $\theta ^{-1}H_n$ from \cref {dist:eq:taugivenk} are shown with the large-$n$ approximation $\theta ^{-1}(\log n+\gamma )$ and the linear budding comparator $(n+1)/(d_{\Icell }+p)$ at $d_{\Icell }+p=8$. The two grow at different orders: logarithmically for bursting, linearly for budding}{figure.5.10}{}}
\newlabel{dist:fig:taugivenk@cref}{{[figure][10][5]5.10}{[1][134][]135}}
\newlabel{dist:eq:condcdf}{{5.58}{135}{Burst time given burst size, without death}{equation.5.7.58}{}}
\newlabel{dist:eq:condcdf@cref}{{[equation][58][5]5.58}{[1][134][]135}}
\newlabel{dist:eq:gumbel}{{5.59}{135}{Burst time given burst size, without death}{equation.5.7.59}{}}
\newlabel{dist:eq:gumbel@cref}{{[equation][59][5]5.59}{[1][135][]135}}
\newlabel{dist:sec:tau-burst-mu-pos}{{5.7.3}{135}{The mean burst time when death is possible}{subsection.5.7.3}{}}
\newlabel{dist:sec:tau-burst-mu-pos@cref}{{[subsection][3][5,7]5.7.3}{[1][135][]135}}
\newlabel{dist:prop:tauburst}{{5.13}{135}{Mean burst time conditioned on bursting}{theorem.5.13}{}}
\newlabel{dist:prop:tauburst@cref}{{[proposition][13][5]5.13}{[1][135][]135}}
\newlabel{dist:eq:tauburst}{{5.60}{135}{Mean burst time conditioned on bursting}{equation.5.7.60}{}}
\newlabel{dist:eq:tauburst@cref}{{[equation][60][5]5.60}{[1][135][]135}}
\newlabel{dist:sec:cond-means}{{5.8}{136}{Conditional burst means}{section.5.8}{}}
\newlabel{dist:sec:cond-means@cref}{{[section][8][5]5.8}{[1][136][]136}}
\newlabel{dist:sec:overall-mean}{{5.8.1}{136}{The overall expected burst size}{subsection.5.8.1}{}}
\newlabel{dist:sec:overall-mean@cref}{{[subsection][1][5,8]5.8.1}{[1][136][]136}}
\newlabel{dist:eq:voveri}{{5.61}{136}{The overall expected burst size}{equation.5.8.61}{}}
\newlabel{dist:eq:voveri@cref}{{[equation][61][5]5.61}{[1][136][]136}}
\newlabel{dist:eq:meank}{{5.62}{136}{The overall expected burst size}{equation.5.8.62}{}}
\newlabel{dist:eq:meank@cref}{{[equation][62][5]5.62}{[1][136][]136}}
\newlabel{dist:sec:circularity}{{5.8.2}{136}{The apparent circularity}{subsection.5.8.2}{}}
\newlabel{dist:sec:circularity@cref}{{[subsection][2][5,8]5.8.2}{[1][136][]136}}
\newlabel{dist:eq:kdecomp}{{5.63}{136}{The apparent circularity}{equation.5.8.63}{}}
\newlabel{dist:eq:kdecomp@cref}{{[equation][63][5]5.63}{[1][136][]136}}
\newlabel{dist:eq:kcond}{{5.64}{136}{The apparent circularity}{equation.5.8.64}{}}
\newlabel{dist:eq:kcond@cref}{{[equation][64][5]5.64}{[1][136][]136}}
\newlabel{dist:tab:threemeans}{{5.2}{137}{Three means, three conditionings. The first two differ by the factor $1-b$, the probability of eventual rupture; the second and third by the size-biasing factor $(a+1)/a$ of \cref {dist:sec:size-bias}}{table.5.2}{}}
\newlabel{dist:tab:threemeans@cref}{{[table][2][5]5.2}{[1][137][]137}}
\newlabel{dist:eq:consistent}{{5.65}{137}{The apparent circularity}{equation.5.8.65}{}}
\newlabel{dist:eq:consistent@cref}{{[equation][65][5]5.65}{[1][136][]137}}
\newlabel{dist:eq:patient}{{5.66}{137}{The apparent circularity}{equation.5.8.66}{}}
\newlabel{dist:eq:patient@cref}{{[equation][66][5]5.66}{[1][137][]137}}
\newlabel{dist:eq:latemean}{{5.67}{137}{The apparent circularity}{equation.5.8.67}{}}
\newlabel{dist:eq:latemean@cref}{{[equation][67][5]5.67}{[1][137][]137}}
\newlabel{dist:eq:kjP}{{5.68}{137}{The apparent circularity}{equation.5.8.68}{}}
\newlabel{dist:eq:kjP@cref}{{[equation][68][5]5.68}{[1][137][]137}}
\newlabel{dist:rem:circularity}{{5.14}{137}{Where the circularity was}{theorem.5.14}{}}
\newlabel{dist:rem:circularity@cref}{{[remark][14][5]5.14}{[1][137][]137}}
\newlabel{dist:sec:moi}{{5.9}{137}{Several founders}{section.5.9}{}}
\newlabel{dist:sec:moi@cref}{{[section][9][5]5.9}{[1][137][]137}}
\newlabel{dist:fig:threemeans}{{5.11}{138}{Three distinct averages at $(\lambda ,\mu ,\delta )=(1,0.2,0.05)$, drawn as trajectories. The unconditional mean release $V(t)$ approaches $V_\infty =13.99$ and includes cells that never burst; the burst mean $V_\infty /(1-b)=a/(a-1)=17.23$ is the horizontal line; conditioning on a very late eventual burst approaches the size-biased limit $(a+1)/(a-1)=33.46$. Keeping the conditioning explicit removes the apparent circularity}{figure.5.11}{}}
\newlabel{dist:fig:threemeans@cref}{{[figure][11][5]5.11}{[1][137][]138}}
\newlabel{dist:sec:moi-fixation}{{5.9.1}{138}{Fixation functions for \texorpdfstring {$k$}{k} founders}{subsection.5.9.1}{}}
\newlabel{dist:sec:moi-fixation@cref}{{[subsection][1][5,9]5.9.1}{[1][138][]138}}
\newlabel{dist:eq:ihatk}{{5.69}{138}{Fixation functions for \texorpdfstring {$k$}{k} founders}{equation.5.9.69}{}}
\newlabel{dist:eq:ihatk@cref}{{[equation][69][5]5.69}{[1][138][]138}}
\newlabel{dist:rem:ihatk}{{5.15}{138}{A guess that fails}{theorem.5.15}{}}
\newlabel{dist:rem:ihatk@cref}{{[remark][15][5]5.15}{[1][138][]138}}
\newlabel{dist:sec:moi-moments}{{5.9.2}{138}{Moments and the release kernel}{subsection.5.9.2}{}}
\newlabel{dist:sec:moi-moments@cref}{{[subsection][2][5,9]5.9.2}{[1][138][]138}}
\newlabel{dist:eq:jk}{{5.70}{138}{Moments and the release kernel}{equation.5.9.70}{}}
\newlabel{dist:eq:jk@cref}{{[equation][70][5]5.70}{[1][138][]138}}
\newlabel{dist:eq:kk}{{5.71}{139}{Moments and the release kernel}{equation.5.9.71}{}}
\newlabel{dist:eq:kk@cref}{{[equation][71][5]5.71}{[1][138][]139}}
\newlabel{dist:eq:gkflux}{{5.72}{139}{Moments and the release kernel}{equation.5.9.72}{}}
\newlabel{dist:eq:gkflux@cref}{{[equation][72][5]5.72}{[1][139][]139}}
\newlabel{dist:rem:gk}{{5.16}{139}{The free-sum trap}{theorem.5.16}{}}
\newlabel{dist:rem:gk@cref}{{[remark][16][5]5.16}{[1][139][]139}}
\newlabel{dist:sec:moi-yield}{{5.9.3}{139}{Lifetime yield}{subsection.5.9.3}{}}
\newlabel{dist:sec:moi-yield@cref}{{[subsection][3][5,9]5.9.3}{[1][139][]139}}
\newlabel{dist:eq:vk1}{{5.73}{139}{Lifetime yield}{equation.5.9.73}{}}
\newlabel{dist:eq:vk1@cref}{{[equation][73][5]5.73}{[1][139][]139}}
\newlabel{dist:eq:vk}{{5.74}{139}{Lifetime yield}{equation.5.9.74}{}}
\newlabel{dist:eq:vk@cref}{{[equation][74][5]5.74}{[1][139][]139}}
\newlabel{dist:eq:vkmuzero}{{5.75}{139}{Lifetime yield}{equation.5.9.75}{}}
\newlabel{dist:eq:vkmuzero@cref}{{[equation][75][5]5.75}{[1][139][]139}}
\newlabel{dist:fig:subextensive}{{5.12}{140}{Unconditional mean yield $V_\infty ^{(k)}$ against founding multiplicity $k$, compared with the naive independent-cell line $kV_\infty $, for $(\lambda ,\mu ,\delta )=(1,0.2,0.05)$. The true yield is sub-extensive because all founders share one cell and hence one catastrophe; for $k=2$, $V_\infty ^{(2)}=17.43$ rather than $2V_\infty =27.97$}{figure.5.12}{}}
\newlabel{dist:fig:subextensive@cref}{{[figure][12][5]5.12}{[1][140][]140}}
\newlabel{dist:fig:moi}{{5.13}{140}{Multiplicity of infection at $(\lambda ,\mu ,\delta )=(1,0.2,0.05)$. (a) Release fluxes $g_k(t)=\delta K_k(t)$ from \cref {dist:eq:gkflux} for $k=1,\ldots ,4$: early ages are ordered superlinearly in founding load, while late fluxes reverse as multi-founder lineages fix faster, so the enhancement is age-dependent rather than a pointwise ordering valid for all $t$. (b) Conditional mean burst $V_\infty ^{(k)}/(1-b^k)$ increases with $k$; the dashed line is the $k=1$ value $a/(a-1)=17.23$}{figure.5.13}{}}
\newlabel{dist:fig:moi@cref}{{[figure][13][5]5.13}{[1][140][]140}}
\newlabel{dist:sec:chained}{{5.10}{140}{Chained immediate transfer}{section.5.10}{}}
\newlabel{dist:sec:chained@cref}{{[section][10][5]5.10}{[1][140][]140}}
\newlabel{dist:eq:srho}{{5.76}{141}{Chained immediate transfer}{equation.5.10.76}{}}
\newlabel{dist:eq:srho@cref}{{[equation][76][5]5.76}{[1][141][]141}}
\newlabel{dist:sec:key-observation}{{5.10.1}{141}{The key observation: event types ignore the load}{subsection.5.10.1}{}}
\newlabel{dist:sec:key-observation@cref}{{[subsection][1][5,10]5.10.1}{[1][141][]141}}
\newlabel{dist:eq:geomevents}{{5.77}{141}{The key observation: event types ignore the load}{equation.5.10.77}{}}
\newlabel{dist:eq:geomevents@cref}{{[equation][77][5]5.77}{[1][141][]141}}
\newlabel{dist:eq:rkdecomp}{{5.78}{141}{The key observation: event types ignore the load}{equation.5.10.78}{}}
\newlabel{dist:eq:rkdecomp@cref}{{[equation][78][5]5.78}{[1][141][]141}}
\newlabel{dist:sec:rupture-sizes}{{5.10.2}{141}{Rupture-size laws}{subsection.5.10.2}{}}
\newlabel{dist:sec:rupture-sizes@cref}{{[subsection][2][5,10]5.10.2}{[1][141][]141}}
\newlabel{dist:eq:rklaw}{{5.79}{141}{Rupture-size laws}{equation.5.10.79}{}}
\newlabel{dist:eq:rklaw@cref}{{[equation][79][5]5.79}{[1][141][]141}}
\newlabel{dist:eq:r1}{{5.80}{141}{Rupture-size laws}{equation.5.10.80}{}}
\newlabel{dist:eq:r1@cref}{{[equation][80][5]5.80}{[1][141][]141}}
\newlabel{dist:eq:r234}{{5.81}{141}{Rupture-size laws}{equation.5.10.81}{}}
\newlabel{dist:eq:r234@cref}{{[equation][81][5]5.81}{[1][141][]141}}
\newlabel{dist:eq:arec}{{5.82}{141}{Rupture-size laws}{equation.5.10.82}{}}
\newlabel{dist:eq:arec@cref}{{[equation][82][5]5.82}{[1][141][]141}}
\newlabel{dist:rem:coefficient}{{5.17}{142}{The coefficient}{theorem.5.17}{}}
\newlabel{dist:rem:coefficient@cref}{{[remark][17][5]5.17}{[1][141][]142}}
\newlabel{dist:eq:rkmoments}{{5.83}{142}{Rupture-size laws}{equation.5.10.83}{}}
\newlabel{dist:eq:rkmoments@cref}{{[equation][83][5]5.83}{[1][142][]142}}
\newlabel{dist:eq:rkgf}{{5.84}{142}{Rupture-size laws}{equation.5.10.84}{}}
\newlabel{dist:eq:rkgf@cref}{{[equation][84][5]5.84}{[1][142][]142}}
\newlabel{dist:sec:rupture-times}{{5.10.3}{142}{Rupture times}{subsection.5.10.3}{}}
\newlabel{dist:sec:rupture-times@cref}{{[subsection][3][5,10]5.10.3}{[1][142][]142}}
\newlabel{dist:eq:t1m}{{5.85}{142}{Rupture times}{equation.5.10.85}{}}
\newlabel{dist:eq:t1m@cref}{{[equation][85][5]5.85}{[1][142][]142}}
\newlabel{dist:eq:genfix}{{5.86}{142}{Rupture times}{equation.5.10.86}{}}
\newlabel{dist:eq:genfix@cref}{{[equation][86][5]5.86}{[1][142][]142}}
\newlabel{dist:sec:chained-scope}{{5.10.4}{142}{Where the structure breaks}{subsection.5.10.4}{}}
\newlabel{dist:sec:chained-scope@cref}{{[subsection][4][5,10]5.10.4}{[1][142][]142}}
\newlabel{dist:sec:budding}{{5.11}{142}{Bursting against budding}{section.5.11}{}}
\newlabel{dist:sec:budding@cref}{{[section][11][5]5.11}{[1][142][]142}}
\newlabel{dist:fig:chained}{{5.14}{143}{Chained immediate transfer with $(\lambda ,\delta )=(1,1)$ and $s=\rho =1/2$. (a) Simulated rupture-size probabilities (bars) from $10^6$ chain realisations with seed 42, with the exact negative-binomial laws \cref {dist:eq:rklaw} (curves) for $k=1,\ldots ,4$. (b) The analytic mean interval $\E {T(k)}$ of \cref {dist:eq:meant} at fixed founding load $k$, and the simulated mean interval at chain position $k$; both decrease over $k=1,\ldots ,6$, the latter faster because its founding loads are random and growing}{figure.5.14}{}}
\newlabel{dist:fig:chained@cref}{{[figure][14][5]5.14}{[1][142][]143}}
\newlabel{dist:eq:bud}{{5.87}{143}{Bursting against budding}{equation.5.11.87}{}}
\newlabel{dist:eq:bud@cref}{{[equation][87][5]5.87}{[1][143][]143}}
\newlabel{dist:eq:budgeom}{{5.88}{143}{Bursting against budding}{equation.5.11.88}{}}
\newlabel{dist:eq:budgeom@cref}{{[equation][88][5]5.88}{[1][143][]143}}
\newlabel{dist:tab:budburst}{{5.3}{144}{Single-cell budding and bursting compared, $\mu =0$ for bursting. Budding parameters: release rate $p$ and infected-cell death rate $d_{\Icell }$. Bursting parameters: division rate $\lambda $ and rupture rate $\delta $ per intracellular particle; $a=(\lambda +\delta )/\lambda $ here}{table.5.3}{}}
\newlabel{dist:tab:budburst@cref}{{[table][3][5]5.3}{[1][143][]144}}
\newlabel{dist:fig:budding}{{5.15}{144}{Release-count laws for the BDC process at $(\lambda ,\mu ,\delta )=(1,0.2,0.05)$ and a budding comparator whose mean $p/d_{\Icell }$ is matched to the BDC conditional mean $V_\infty /(1-b)=17.23$. Panel~(a): the BDC law conditioned on bursting is geometric on $\{1,2,\ldots \}$, whereas the budding law is geometric on $\{0,1,2,\ldots \}$ and retains probability $0.055$ of zero release. Panel~(b): survival functions $\pr {\Kb \ge k}$ and $\pr {\Kb _{\rm bud}\ge k}$. Matching the means matches neither the support nor the tails}{figure.5.15}{}}
\newlabel{dist:fig:budding@cref}{{[figure][15][5]5.15}{[1][144][]144}}
\newlabel{dist:sec:discussion}{{5.12}{144}{Discussion}{section.5.12}{}}
\newlabel{dist:sec:discussion@cref}{{[section][12][5]5.12}{[1][144][]144}}
\newlabel{dist:sec:discussion-assay}{{5.12.1}{145}{Single-cell assay predictions for lytic systems}{subsection.5.12.1}{}}
\newlabel{dist:sec:discussion-assay@cref}{{[subsection][1][5,12]5.12.1}{[1][145][]145}}
\newlabel{dist:sec:two-extremes}{{5.12.2}{145}{The same theory at two extremes}{subsection.5.12.2}{}}
\newlabel{dist:sec:two-extremes@cref}{{[subsection][2][5,12]5.12.2}{[1][145][]145}}
\newlabel{dist:tab:twoextremes}{{5.4}{146}{The same three-rate theory at two published parameter sets, from \cite {carruthers2020stochastic} and \cite {williams2021anthrax}. Every entry is a closed form of this chapter evaluated at those rates: \cref {dist:eq:roots,dist:eq:vinf,dist:eq:qsd,dist:eq:tprod,dist:eq:tauburst}. The digits shown are those of the evaluation, not of the estimate: each entry is exact given the rates in the first row, and inherits whatever uncertainty those rates carry. The two columns come from different methods applied to different data, so the comparison between them is qualitative even where the arithmetic within a column is not}{table.5.4}{}}
\newlabel{dist:tab:twoextremes@cref}{{[table][4][5]5.4}{[1][145][]146}}
\newlabel{dist:sec:probes}{{5.12.3}{146}{What the two probes showed}{subsection.5.12.3}{}}
\newlabel{dist:sec:probes@cref}{{[subsection][3][5,12]5.12.3}{[1][146][]146}}
\newlabel{dist:sec:open-problems}{{5.12.4}{146}{Open problems}{subsection.5.12.4}{}}
\newlabel{dist:sec:open-problems@cref}{{[subsection][4][5,12]5.12.4}{[1][146][]146}}
\newlabel{dist:sec:next}{{5.12.5}{147}{What comes next}{subsection.5.12.5}{}}
\newlabel{dist:sec:next@cref}{{[subsection][5][5,12]5.12.5}{[1][147][]147}}
\newlabel{dist:app:formulae}{{5.A}{147}{Master formula table: single cell}{section.5.appendix.A}{}}
\newlabel{dist:app:formulae@cref}{{[section][1][5]5.A}{[1][147][]147}}
\newlabel{dist:tab:formulae}{{5.5}{147}{Every closed form the chapter uses, with its provenance}{table.5.5}{}}
\newlabel{dist:tab:formulae@cref}{{[table][5][5]5.5}{[1][147][]147}}
\newlabel{dist:app:verification}{{5.B}{148}{Verification record: chained immediate transfer}{section.5.appendix.B}{}}
\newlabel{dist:app:verification@cref}{{[section][2][5]5.B}{[1][148][]148}}
\newlabel{dist:app:technical}{{5.C}{149}{Technical derivations}{section.5.appendix.C}{}}
\newlabel{dist:app:technical@cref}{{[section][3][5]5.C}{[1][149][]149}}
\newlabel{dist:app:vk}{{5.C.1}{149}{Lifetime yield from \texorpdfstring {$k$}{k} founders}{subsection.5.appendix.C.1}{}}
\newlabel{dist:app:vk@cref}{{[subsection][1][5,3]5.C.1}{[1][149][]149}}
\newlabel{dist:app:extract}{{5.C.2}{150}{Coefficient extraction for general initial conditions}{subsection.5.appendix.C.2}{}}
\newlabel{dist:app:extract@cref}{{[subsection][2][5,3]5.C.2}{[1][150][]150}}
\newlabel{dist:eq:gkgamma}{{5.89}{150}{Coefficient extraction for general initial conditions}{equation.5.appendix.C.89}{}}
\newlabel{dist:eq:gkgamma@cref}{{[equation][89][5]5.89}{[1][150][]150}}
\newlabel{dist:eq:zHGk}{{5.90}{150}{Coefficient extraction for general initial conditions}{equation.5.appendix.C.90}{}}
\newlabel{dist:eq:zHGk@cref}{{[equation][90][5]5.90}{[1][150][]150}}
\newlabel{dist:eq:cexpand}{{5.91}{150}{Coefficient extraction for general initial conditions}{equation.5.appendix.C.91}{}}
\newlabel{dist:eq:cexpand@cref}{{[equation][91][5]5.91}{[1][150][]150}}
\newlabel{dist:eq:hypantideriv}{{5.92}{150}{Coefficient extraction for general initial conditions}{equation.5.appendix.C.92}{}}
\newlabel{dist:eq:hypantideriv@cref}{{[equation][92][5]5.92}{[1][150][]150}}
\newlabel{dist:app:chained-transforms}{{5.C.3}{150}{Inter-rupture intervals and the second rupture time}{subsection.5.appendix.C.3}{}}
\newlabel{dist:app:chained-transforms@cref}{{[subsection][3][5,3]5.C.3}{[1][150][]150}}
\newlabel{dist:eq:tkdecomp}{{5.93}{150}{Inter-rupture intervals and the second rupture time}{equation.5.appendix.C.93}{}}
\newlabel{dist:eq:tkdecomp@cref}{{[equation][93][5]5.93}{[1][150][]150}}
\newlabel{dist:eq:ltk}{{5.94}{150}{Inter-rupture intervals and the second rupture time}{equation.5.appendix.C.94}{}}
\newlabel{dist:eq:ltk@cref}{{[equation][94][5]5.94}{[1][150][]150}}
\newlabel{dist:eq:ltkhyp}{{5.95}{151}{Inter-rupture intervals and the second rupture time}{equation.5.appendix.C.95}{}}
\newlabel{dist:eq:ltkhyp@cref}{{[equation][95][5]5.95}{[1][150][]151}}
\newlabel{dist:eq:meant}{{5.96}{151}{Inter-rupture intervals and the second rupture time}{equation.5.appendix.C.96}{}}
\newlabel{dist:eq:meant@cref}{{[equation][96][5]5.96}{[1][150][]151}}
\newlabel{dist:eq:t2cond}{{5.97}{151}{Inter-rupture intervals and the second rupture time}{equation.5.appendix.C.97}{}}
\newlabel{dist:eq:t2cond@cref}{{[equation][97][5]5.97}{[1][151][]151}}
\newlabel{dist:eq:lt2}{{5.98}{151}{Inter-rupture intervals and the second rupture time}{equation.5.appendix.C.98}{}}
\newlabel{dist:eq:lt2@cref}{{[equation][98][5]5.98}{[1][151][]151}}
\newlabel{p:ch:one-cell-to-population}{{6}{152}{Burst-aware within-host dynamics}{chapter.6}{}}
\newlabel{p:ch:one-cell-to-population@cref}{{[chapter][6][]6}{[1][152][]152}}
\newlabel{p:sec:introduction}{{6.1}{152}{Introduction}{section.6.1}{}}
\newlabel{p:sec:introduction@cref}{{[section][1][6]6.1}{[1][152][]152}}
\newlabel{p:sec:assumptions}{{6.1.1}{153}{Standing assumptions}{subsection.6.1.1}{}}
\newlabel{p:sec:assumptions@cref}{{[subsection][1][6,1]6.1.1}{[1][153][]153}}
\newlabel{p:tab:notation}{{6.1}{154}{Population notation. Single-cell symbols are those of \cref {p:app:quoted}}{table.6.1}{}}
\newlabel{p:tab:notation@cref}{{[table][1][6]6.1}{[1][153][]154}}
\newlabel{p:sec:comparator}{{6.2}{154}{The standard model and the cell it describes exactly}{section.6.2}{}}
\newlabel{p:sec:comparator@cref}{{[section][2][6]6.2}{[1][154][]154}}
\newlabel{p:sec:classical}{{6.2.1}{154}{The standard model}{subsection.6.2.1}{}}
\newlabel{p:sec:classical@cref}{{[subsection][1][6,2]6.2.1}{[1][154][]154}}
\newlabel{p:eq:bmvr-classical}{{6.1}{154}{The standard model}{equation.6.2.1}{}}
\newlabel{p:eq:bmvr-classical@cref}{{[equation][1][6]6.1}{[1][154][]154}}
\newlabel{p:eq:R0classical}{{6.2}{154}{The standard model}{equation.6.2.2}{}}
\newlabel{p:eq:R0classical@cref}{{[equation][2][6]6.2}{[1][154][]154}}
\newlabel{p:sec:budding}{{6.2.2}{154}{The budding cell is described exactly}{subsection.6.2.2}{}}
\newlabel{p:sec:budding@cref}{{[subsection][2][6,2]6.2.2}{[1][154][]154}}
\newlabel{p:eq:bud-cell}{{6.3}{154}{The budding cell is described exactly}{equation.6.2.3}{}}
\newlabel{p:eq:bud-cell@cref}{{[equation][3][6]6.3}{[1][154][]154}}
\newlabel{p:fig:constant-fails}{{6.1}{155}{A synchronous cohort releases on a schedule, not at a rate. Both schedules have lifetime yield $V_\infty =11$ and mean productive lifetime $\mean {T_{\rm prod}}=2.398$, with $d_{\Icell }=1/\mean {T_{\rm prod}}$, $p=V_\infty d_{\Icell }$, $(\lambda ,\mu ,\delta )=(1,0,0.1)$ and $c=1$: as fairly matched as two constants allow. (a) At matched yield the release kernel $g(\alpha )=\delta K(\alpha )$ still lags the budding cell's $p\ee ^{-d_{\Icell }\alpha }$, because a bursting cell must first accumulate load. (b) Linear clearance preserves that lag, carrying it into the free-particle response}{figure.6.1}{}}
\newlabel{p:fig:constant-fails@cref}{{[figure][1][6]6.1}{[1][155][]155}}
\newlabel{p:eq:bud-geom}{{6.4}{155}{The budding cell is described exactly}{equation.6.2.4}{}}
\newlabel{p:eq:bud-geom@cref}{{[equation][4][6]6.4}{[1][154][]155}}
\newlabel{p:sec:obstruction}{{6.2.3}{155}{Bursting, and why no constant repairs it}{subsection.6.2.3}{}}
\newlabel{p:sec:obstruction@cref}{{[subsection][3][6,2]6.2.3}{[1][155][]155}}
\newlabel{p:sec:kernels}{{6.3}{156}{The intracellular process and its two kernels}{section.6.3}{}}
\newlabel{p:sec:kernels@cref}{{[section][3][6]6.3}{[1][155][]156}}
\newlabel{p:eq:roots}{{6.5}{156}{The process}{equation.6.3.5}{}}
\newlabel{p:eq:roots@cref}{{[equation][5][6]6.5}{[1][156][]156}}
\newlabel{p:eq:Vinf-Tprod}{{6.6}{156}{The process}{equation.6.3.6}{}}
\newlabel{p:eq:Vinf-Tprod@cref}{{[equation][6][6]6.6}{[1][156][]156}}
\newlabel{p:eq:Skernel}{{6.7}{156}{The two kernels}{equation.6.3.7}{}}
\newlabel{p:eq:Skernel@cref}{{[equation][7][6]6.7}{[1][156][]156}}
\newlabel{p:eq:gkernel}{{6.8}{156}{The two kernels}{equation.6.3.8}{}}
\newlabel{p:eq:gkernel@cref}{{[equation][8][6]6.8}{[1][156][]156}}
\newlabel{p:sec:trio}{{6.3.3}{156}{Three regimes, used throughout}{subsection.6.3.3}{}}
\newlabel{p:sec:trio@cref}{{[subsection][3][6,3]6.3.3}{[1][156][]156}}
\newlabel{p:tab:trio}{{6.2}{156}{The three regimes used throughout. The first has $L>1$, the second sits on the boundary $L=1$, the third has $L<1$. The last two columns are the $r=0$ endpoints of the effective-parameter map of \cref {p:sec:projection}}{table.6.2}{}}
\newlabel{p:tab:trio@cref}{{[table][2][6]6.2}{[1][156][]156}}
\newlabel{p:fig:kernels}{{6.2}{157}{Survival $\Ifix $ (a)--(c) and release $g=\delta K$ (d)--(f) across the flooding boundary of \cref {p:tab:trio}. Release rises then falls --- a cell must accumulate before it bursts --- while survival decays; shaded areas check $\int _0^\infty \Ifix =\mean {T_{\rm prod}}$ and $\int _0^\infty g=V_\infty $. Parameters, left to right: $(\lambda ,\mu ,\delta )=(1,0,0.1)$, $(1,0.5,1/3)$, $(1,0.9,0.1)$, with $L=1.10$, $1.00$, $0.42$ and $V_\infty =11.00$, $2.00$, $1.32$. Independent $g$-scales keep the timing shift visible}{figure.6.2}{}}
\newlabel{p:fig:kernels@cref}{{[figure][2][6]6.2}{[1][156][]157}}
\newlabel{p:sec:renewal}{{6.4}{157}{Burst-aware renewal dynamics}{section.6.4}{}}
\newlabel{p:sec:renewal@cref}{{[section][4][6]6.4}{[1][157][]157}}
\newlabel{p:def:renewal-system}{{6.1}{157}{The burst-aware renewal system}{theorem.6.1}{}}
\newlabel{p:def:renewal-system@cref}{{[definition][1][6]6.1}{[1][157][]157}}
\newlabel{p:eq:Icell}{{6.9}{157}{The burst-aware renewal system}{equation.6.4.9}{}}
\newlabel{p:eq:Icell@cref}{{[equation][9][6]6.9}{[1][157][]157}}
\newlabel{p:eq:Vfree}{{6.10}{157}{The burst-aware renewal system}{equation.6.4.10}{}}
\newlabel{p:eq:Vfree@cref}{{[equation][10][6]6.10}{[1][157][]157}}
\newlabel{p:fig:renewal-schematic}{{6.3}{158}{Bookkeeping of \cref {p:def:renewal-system}: infection events start cohorts; the kernels $\Ifix $ and $g=\delta K$ weight them by age; the feedback $i=\gamma T\Vfree $ closes the loop}{figure.6.3}{}}
\newlabel{p:fig:renewal-schematic@cref}{{[figure][3][6]6.3}{[1][157][]158}}
\newlabel{p:tab:correspondence}{{6.3}{158}{What the renewal system replaces. Both constants go, not only the release rate}{table.6.3}{}}
\newlabel{p:tab:correspondence@cref}{{[table][3][6]6.3}{[1][157][]158}}
\newlabel{p:prop:first-moment}{{6.2}{158}{First-moment exactness}{theorem.6.2}{}}
\newlabel{p:prop:first-moment@cref}{{[proposition][2][6]6.2}{[1][158][]158}}
\newlabel{p:sec:generality}{{6.4.3}{159}{One skeleton, many mechanisms}{subsection.6.4.3}{}}
\newlabel{p:sec:generality@cref}{{[subsection][3][6,4]6.4.3}{[1][159][]159}}
\newlabel{p:eq:skeleton}{{6.11}{159}{One skeleton, many mechanisms}{equation.6.4.11}{}}
\newlabel{p:eq:skeleton@cref}{{[equation][11][6]6.11}{[1][159][]159}}
\newlabel{p:tab:sg}{{6.4}{159}{One skeleton, many mechanisms. Each row is a different single-cell story; all induce \cref {p:eq:skeleton}, and differ only in the pair of kernels they supply}{table.6.4}{}}
\newlabel{p:tab:sg@cref}{{[table][4][6]6.4}{[1][159][]159}}
\newlabel{p:sec:relation}{{6.4.4}{159}{Relation to existing multiscale models}{subsection.6.4.4}{}}
\newlabel{p:sec:relation@cref}{{[subsection][4][6,4]6.4.4}{[1][159][]159}}
\newlabel{p:fig:gillespie}{{6.4}{160}{Stochastic ground truth. Ensemble means of an independent Gillespie simulation --- birth--death--catastrophe cells with impulse releases at burst times, exponential clearance, no reinfection --- against \crefrange {p:eq:Icell}{p:eq:Vfree}. Left column $\mu =0$, right column $\mu >0$; (a)--(b) infected cells, (c)--(d) free particles. The curves are indistinguishable at plotting resolution and the relative $L^2$ discrepancy of $1$--$2\%$ is sampling error: the renewal equations are not an approximation to the simulation, and this is the check that says so}{figure.6.4}{}}
\newlabel{p:fig:gillespie@cref}{{[figure][4][6]6.4}{[1][160][]160}}
\newlabel{p:sec:verification}{{6.4.5}{160}{Numerical verification}{subsection.6.4.5}{}}
\newlabel{p:sec:verification@cref}{{[subsection][5][6,4]6.4.5}{[1][160][]160}}
\newlabel{p:sec:projection}{{6.5}{161}{The standard model is a projection}{section.6.5}{}}
\newlabel{p:sec:projection@cref}{{[section][5][6]6.5}{[1][160][]161}}
\newlabel{p:eq:expphase}{{6.12}{161}{Exponential-phase equivalence}{equation.6.5.12}{}}
\newlabel{p:eq:expphase@cref}{{[equation][12][6]6.12}{[1][161][]161}}
\newlabel{p:eq:peff}{{6.13}{161}{Exponential-phase equivalence}{equation.6.5.13}{}}
\newlabel{p:eq:peff@cref}{{[equation][13][6]6.13}{[1][161][]161}}
\newlabel{p:thm:projection}{{6.3}{161}{Exponential-phase equivalence}{theorem.6.3}{}}
\newlabel{p:thm:projection@cref}{{[theorem][3][6]6.3}{[1][161][]161}}
\newlabel{p:eq:char}{{6.14}{161}{Exponential-phase equivalence}{equation.6.5.14}{}}
\newlabel{p:eq:char@cref}{{[equation][14][6]6.14}{[1][161][]161}}
\newlabel{p:eq:peff0}{{6.15}{161}{The ends of the map}{equation.6.5.15}{}}
\newlabel{p:eq:peff0@cref}{{[equation][15][6]6.15}{[1][161][]161}}
\newlabel{p:eq:peffinf}{{6.16}{161}{The ends of the map}{equation.6.5.16}{}}
\newlabel{p:eq:peffinf@cref}{{[equation][16][6]6.16}{[1][161][]161}}
\newlabel{p:prop:oldcell}{{6.4}{162}{The old-cell limit}{theorem.6.4}{}}
\newlabel{p:prop:oldcell@cref}{{[proposition][4][6]6.4}{[1][161][]162}}
\newlabel{p:eq:oldcell}{{6.17}{162}{The old-cell limit}{equation.6.5.17}{}}
\newlabel{p:eq:oldcell@cref}{{[equation][17][6]6.17}{[1][161][]162}}
\newlabel{p:sec:span}{{6.5.3}{162}{Monotonicity, and the size of the effect}{subsection.6.5.3}{}}
\newlabel{p:sec:span@cref}{{[subsection][3][6,5]6.5.3}{[1][162][]162}}
\newlabel{p:prop:peff-monotone}{{6.5}{162}{Monotonicity of the effective release rate}{theorem.6.5}{}}
\newlabel{p:prop:peff-monotone@cref}{{[proposition][5][6]6.5}{[1][162][]162}}
\newlabel{p:eq:peff-deriv}{{6.18}{162}{Monotonicity, and the size of the effect}{equation.6.5.18}{}}
\newlabel{p:eq:peff-deriv@cref}{{[equation][18][6]6.18}{[1][162][]162}}
\newlabel{p:eq:rho-closed}{{6.19}{162}{Monotonicity, and the size of the effect}{equation.6.5.19}{}}
\newlabel{p:eq:rho-closed@cref}{{[equation][19][6]6.19}{[1][162][]162}}
\newlabel{p:eq:peff-span}{{6.20}{162}{Monotonicity, and the size of the effect}{equation.6.5.20}{}}
\newlabel{p:eq:peff-span@cref}{{[equation][20][6]6.20}{[1][162][]162}}
\newlabel{p:fig:peff}{{6.5}{163}{The effective release rate (a) and removal rate (b) of \cref {p:eq:peff} against the epidemic growth rate $r$, for three parameter sets, over the whole domain $r>-\theta $ on which the map is defined. Open circles mark the $r=0$ values \cref {p:eq:peff0}; dotted segments mark the limits each curve runs toward. Note that $r=0$ is an interior point: the left-hand endpoint is the old-cell limit $\delta \mathbb {E}_{\rm QS}[X^2]$ and the right-hand one the young-cell limit $\delta $. Read the vertical extent of each curve as the size of the error made by treating $p$ as a constant of the pathogen rather than of the epidemic's phase}{figure.6.5}{}}
\newlabel{p:fig:peff@cref}{{[figure][5][6]6.5}{[1][162][]163}}
\newlabel{p:sec:generation-kernel}{{6.5.4}{163}{Threshold and growth rate as functionals of one kernel}{subsection.6.5.4}{}}
\newlabel{p:sec:generation-kernel@cref}{{[subsection][4][6,5]6.5.4}{[1][162][]163}}
\newlabel{p:eq:genkernel}{{6.21}{163}{Threshold and growth rate as functionals of one kernel}{equation.6.5.21}{}}
\newlabel{p:eq:genkernel@cref}{{[equation][21][6]6.21}{[1][162][]163}}
\newlabel{p:eq:R0gen}{{6.22}{163}{Threshold and growth rate as functionals of one kernel}{equation.6.5.22}{}}
\newlabel{p:eq:R0gen@cref}{{[equation][22][6]6.22}{[1][163][]163}}
\newlabel{p:thm:R0inv}{{6.6}{163}{$R_0$ is unchanged by bursting}{theorem.6.6}{}}
\newlabel{p:thm:R0inv@cref}{{[theorem][6][6]6.6}{[1][163][]163}}
\newlabel{p:fig:overlay-V}{{6.6}{164}{Both state variables, both sides of the threshold. Free particles $\Vfree (t)$ (top row) and infected cells $\Icell (t)$ (bottom) under the renewal system (solid) and under standard BMVR matched at $r=0$ (dashed), for a supercritical infection (left, $R_0=1.76$) and a subcritical one (right, $R_0=0.28$). The matched model gets the order of magnitude and not the shape; and the discrepancy is not confined to the release equation --- age structure moves the cell trajectory too, which is why matching $p$ alone cannot repair it. Parameters: $(\lambda ,\mu ,\delta )=(1,0,0.1)$ throughout, with $\gamma T=0.04$, $c=0.25$ on the left and $\gamma T=0.01$, $c=0.4$ on the right}{figure.6.6}{}}
\newlabel{p:fig:overlay-V@cref}{{[figure][6][6]6.6}{[1][164][]164}}
\newlabel{p:sec:overlays}{{6.5.5}{164}{What the projection loses}{subsection.6.5.5}{}}
\newlabel{p:sec:overlays@cref}{{[subsection][5][6,5]6.5.5}{[1][164][]164}}
\newlabel{p:fig:growth-phase}{{6.7}{165}{Why no single constant pair will do. One renewal trajectory (solid) against two standard parameter choices: the $r=0$ match (dashed) and the young-cell limit $p=\delta $, $d_{\Icell }=\mu +\delta $ (dash-dot). (a) The first six time units, where the $r=0$ match overshoots by more than an order of magnitude while the young-cell parameters are the closer of the two. (b) The full window, where the ordering has reversed and the young-cell model has collapsed. Each choice is correct at one growth rate and the trajectory passes through many, exactly as \cref {p:eq:peff} predicts. Both panels logarithmic in $\Vfree $. Parameters: $(\lambda ,\mu ,\delta )=(1,0,0.1)$, $\gamma T=0.025$, $c=0.2$, $R_0=1.38$}{figure.6.7}{}}
\newlabel{p:fig:growth-phase@cref}{{[figure][7][6]6.7}{[1][164][]165}}
\newlabel{p:sec:identifiability}{{6.5.6}{165}{Identifiability: a negative result}{subsection.6.5.6}{}}
\newlabel{p:sec:identifiability@cref}{{[subsection][6][6,5]6.5.6}{[1][165][]165}}
\newlabel{p:fig:identifiability}{{6.8}{166}{Population observables leave level sets rather than identifying the intracellular rates. With $c=1$ and $\gamma T=1$, the interior flooding point $(\lambda ,\mu ,\delta )=(1,0.5,1/3)$ has $R_0=2$ and $r_\star =0.294$. (a) One population observable leaves a level curve: in the slice $\lambda =1$ the contours $R_0=2$ and $r=r_\star $ are distinct; alternatives A and C share $R_0$ with the reference but not its growth rate. (b) With $\lambda $ free, a continuous family of triples shares both. Joint knowledge of threshold and exponential growth therefore does not identify $(\lambda ,\mu ,\delta )$}{figure.6.8}{}}
\newlabel{p:fig:identifiability@cref}{{[figure][8][6]6.8}{[1][165][]166}}
\newlabel{p:sec:flooding}{{6.6}{166}{Bursting versus budding: the flooding parameter~\texorpdfstring {$L$}{L}}{section.6.6}{}}
\newlabel{p:sec:flooding@cref}{{[section][6][6]6.6}{[1][166][]166}}
\newlabel{p:eq:q}{{6.23}{167}{The offspring of one infected cell}{equation.6.6.23}{}}
\newlabel{p:eq:q@cref}{{[equation][23][6]6.23}{[1][166][]167}}
\newlabel{p:eq:Goff}{{6.24}{167}{The offspring of one infected cell}{equation.6.6.24}{}}
\newlabel{p:eq:Goff@cref}{{[equation][24][6]6.24}{[1][167][]167}}
\newlabel{p:eq:offmoments}{{6.25}{167}{The offspring of one infected cell}{equation.6.6.25}{}}
\newlabel{p:eq:offmoments@cref}{{[equation][25][6]6.25}{[1][167][]167}}
\newlabel{p:eq:zext}{{6.26}{167}{Extinction, and the flooding advantage}{equation.6.6.26}{}}
\newlabel{p:eq:zext@cref}{{[equation][26][6]6.26}{[1][167][]167}}
\newlabel{p:def:L}{{6.7}{167}{The flooding parameter}{theorem.6.7}{}}
\newlabel{p:def:L@cref}{{[definition][7][6]6.7}{[1][167][]167}}
\newlabel{p:eq:Ldef}{{6.27}{167}{The flooding parameter}{equation.6.6.27}{}}
\newlabel{p:eq:Ldef@cref}{{[equation][27][6]6.27}{[1][167][]167}}
\newlabel{p:thm:flood}{{6.8}{167}{The flooding advantage}{theorem.6.8}{}}
\newlabel{p:thm:flood@cref}{{[theorem][8][6]6.8}{[1][167][]167}}
\newlabel{p:eq:flood}{{6.28}{167}{The flooding advantage}{equation.6.6.28}{}}
\newlabel{p:eq:flood@cref}{{[equation][28][6]6.28}{[1][167][]167}}
\newlabel{p:eq:floodcrit}{{6.29}{167}{The flooding advantage}{equation.6.6.29}{}}
\newlabel{p:eq:floodcrit@cref}{{[equation][29][6]6.29}{[1][167][]167}}
\newlabel{p:fig:L-landscape}{{6.9}{168}{Where a pathogen sits relative to the flooding boundary. The flooding parameter $L(\mu ,\delta )$ at $\lambda =1$; the exact curve $L=1$, equivalently $\delta _*=\lambda \mu /(\lambda +\mu )$, partitions the rate plane into a region where bursting lowers extinction risk and one where it raises it. The three marked points are the regimes of \cref {p:tab:trio}, $(\mu ,\delta )=(0,0.1)$, $(0.5,1/3)$ and $(0.9,0.1)$, with $L=1.10$, $1.00$ and $0.42$}{figure.6.9}{}}
\newlabel{p:fig:L-landscape@cref}{{[figure][9][6]6.9}{[1][168][]168}}
\newlabel{p:tab:flooding}{{6.5}{168}{Extinction probabilities from \cref {p:eq:zext} and the sign of their difference. The gap is $(L-1)(1/m-1)$ in every row}{table.6.5}{}}
\newlabel{p:tab:flooding@cref}{{[table][5][6]6.5}{[1][168][]168}}
\newlabel{p:fig:flooding-regimes}{{6.10}{169}{Extinction probabilities for bursting and matched-mean budding across the three regimes of \cref {p:tab:trio}. Each pair of curves begins at $m=1$; the signed gap follows \cref {p:eq:flood}. The $q$-windows differ because $m=1$ sits at different $q$ in each regime}{figure.6.10}{}}
\newlabel{p:fig:flooding-regimes@cref}{{[figure][10][6]6.10}{[1][168][]169}}
\newlabel{p:prop:L1}{{6.9}{169}{}{theorem.6.9}{}}
\newlabel{p:prop:L1@cref}{{[proposition][9][6]6.9}{[1][169][]169}}
\newlabel{p:eq:L1equiv}{{6.30}{169}{}{equation.6.6.30}{}}
\newlabel{p:eq:L1equiv@cref}{{[equation][30][6]6.30}{[1][169][]169}}
\newlabel{p:eq:varorder}{{6.31}{169}{}{equation.6.6.31}{}}
\newlabel{p:eq:varorder@cref}{{[equation][31][6]6.31}{[1][169][]169}}
\newlabel{p:sec:tradeoff}{{6.6.4}{170}{The price: invasion speed}{subsection.6.6.4}{}}
\newlabel{p:sec:tradeoff@cref}{{[subsection][4][6,6]6.6.4}{[1][170][]170}}
\newlabel{p:eq:rorder}{{6.32}{170}{The price: invasion speed}{equation.6.6.32}{}}
\newlabel{p:eq:rorder@cref}{{[equation][32][6]6.32}{[1][170][]170}}
\newlabel{p:tab:tradeoff}{{6.6}{170}{Deterministic growth rates at matched $R_0=2$, matched total yield and matched mean productive lifetime, with $c=1$. Budding invades faster on both sides of the flooding boundary and on it, so the ordering is not the same effect as \cref {p:thm:flood}}{table.6.6}{}}
\newlabel{p:tab:tradeoff@cref}{{[table][6][6]6.6}{[1][170][]170}}
\newlabel{p:prop:rorder}{{6.10}{170}{The ordering is a Laplace-transform comparison}{theorem.6.10}{}}
\newlabel{p:prop:rorder@cref}{{[proposition][10][6]6.10}{[1][170][]170}}
\newlabel{p:eq:laplace-order}{{6.33}{170}{The ordering is a Laplace-transform comparison}{equation.6.6.33}{}}
\newlabel{p:eq:laplace-order@cref}{{[equation][33][6]6.33}{[1][170][]170}}
\newlabel{p:fig:generation-times}{{6.11}{171}{The generation-time structure behind the first row of \cref {p:tab:tradeoff}. At $(\lambda ,\mu ,\delta )=(1,0,0.1)$, budding and bursting are matched at $R_0=2$, $c=1$, total yield and mean productive lifetime. The main curves are the normalised infection-generation laws $\mathcal A(\alpha )/R_0$, including the common extracellular waiting-time component; panel (b) isolates release age and shows why the budding law leads}{figure.6.11}{}}
\newlabel{p:fig:generation-times@cref}{{[figure][11][6]6.11}{[1][171][]171}}
\newlabel{p:fig:pareto}{{6.12}{172}{Tuning $\delta $ traces a strict extinction--growth trade-off, drawn parametrically as $\delta $ varies from $0.03$ to $0.8$ with $\lambda =1$, $\mu =0$, $c=1$ and $R_0=2$. The coupling is reset along the curve to hold $R_0=2$, and the budding comparator is matched in lifetime yield $V_\infty $ and in mean productive lifetime. There is no setting of $\delta $ that improves both coordinates: faster growth is bought with a higher extinction probability}{figure.6.12}{}}
\newlabel{p:fig:pareto@cref}{{[figure][12][6]6.12}{[1][171][]172}}
\newlabel{p:sec:reach}{{6.7}{172}{The reach of the construction}{section.6.7}{}}
\newlabel{p:sec:reach@cref}{{[section][7][6]6.7}{[1][172][]172}}
\newlabel{p:sec:reset}{{6.7.1}{172}{Recurrent release: the reset process}{subsection.6.7.1}{}}
\newlabel{p:sec:reset@cref}{{[subsection][1][6,7]6.7.1}{[1][172][]172}}
\newlabel{p:eq:reset-kernels}{{6.34}{172}{Recurrent release: the reset process}{equation.6.7.34}{}}
\newlabel{p:eq:reset-kernels@cref}{{[equation][34][6]6.34}{[1][172][]172}}
\newlabel{p:eq:reset-peff}{{6.35}{172}{Recurrent release: the reset process}{equation.6.7.35}{}}
\newlabel{p:eq:reset-peff@cref}{{[equation][35][6]6.35}{[1][172][]172}}
\newlabel{p:eq:reset-mature}{{6.36}{172}{Recurrent release: the reset process}{equation.6.7.36}{}}
\newlabel{p:eq:reset-mature@cref}{{[equation][36][6]6.36}{[1][172][]172}}
\newlabel{p:sec:partial}{{6.7.2}{173}{Partial release, and antiviral mode of action}{subsection.6.7.2}{}}
\newlabel{p:sec:partial@cref}{{[subsection][2][6,7]6.7.2}{[1][173][]173}}
\newlabel{p:eq:partial}{{6.37}{173}{Partial release, and antiviral mode of action}{equation.6.7.37}{}}
\newlabel{p:eq:partial@cref}{{[equation][37][6]6.37}{[1][173][]173}}
\newlabel{p:sec:hiv}{{6.7.3}{173}{Beyond birth--death--catastrophe: HIV}{subsection.6.7.3}{}}
\newlabel{p:sec:hiv@cref}{{[subsection][3][6,7]6.7.3}{[1][173][]173}}
\newlabel{p:eq:allee}{{6.38}{174}{Beyond birth--death--catastrophe: HIV}{equation.6.7.38}{}}
\newlabel{p:eq:allee@cref}{{[equation][38][6]6.38}{[1][174][]174}}
\newlabel{p:sec:discussion}{{6.8}{174}{Discussion}{section.6.8}{}}
\newlabel{p:sec:discussion@cref}{{[section][8][6]6.8}{[1][174][]174}}
\newlabel{p:sec:fitting-consequences}{{6.8.1}{174}{Consequences for fitting practice}{subsection.6.8.1}{}}
\newlabel{p:sec:fitting-consequences@cref}{{[subsection][1][6,8]6.8.1}{[1][174][]174}}
\newlabel{p:sec:limitations}{{6.8.2}{175}{Limitations}{subsection.6.8.2}{}}
\newlabel{p:sec:limitations@cref}{{[subsection][2][6,8]6.8.2}{[1][175][]175}}
\newlabel{p:sec:open-problems}{{6.8.3}{176}{Open problems}{subsection.6.8.3}{}}
\newlabel{p:sec:open-problems@cref}{{[subsection][3][6,8]6.8.3}{[1][176][]176}}
\newlabel{p:app:quoted}{{6.A}{177}{Single-cell results quoted from earlier chapters}{section.6.appendix.A}{}}
\newlabel{p:app:quoted@cref}{{[section][1][6]6.A}{[1][176][]177}}
\newlabel{p:eq:derived}{{6.39}{177}{Single-cell results quoted from earlier chapters}{equation.6.appendix.A.39}{}}
\newlabel{p:eq:derived@cref}{{[equation][39][6]6.39}{[1][176][]177}}
\newlabel{p:eq:AB}{{6.40}{177}{Single-cell results quoted from earlier chapters}{equation.6.appendix.A.40}{}}
\newlabel{p:eq:AB@cref}{{[equation][40][6]6.40}{[1][177][]177}}
\newlabel{p:eq:IDIhat}{{6.41}{177}{Single-cell results quoted from earlier chapters}{equation.6.appendix.A.41}{}}
\newlabel{p:eq:IDIhat@cref}{{[equation][41][6]6.41}{[1][177][]177}}
\newlabel{p:eq:J}{{6.42}{177}{Single-cell results quoted from earlier chapters}{equation.6.appendix.A.42}{}}
\newlabel{p:eq:J@cref}{{[equation][42][6]6.42}{[1][177][]177}}
\newlabel{p:eq:K}{{6.43}{177}{Single-cell results quoted from earlier chapters}{equation.6.appendix.A.43}{}}
\newlabel{p:eq:K@cref}{{[equation][43][6]6.43}{[1][177][]177}}
\newlabel{p:eq:V}{{6.44}{177}{Single-cell results quoted from earlier chapters}{equation.6.appendix.A.44}{}}
\newlabel{p:eq:V@cref}{{[equation][44][6]6.44}{[1][177][]177}}
\newlabel{p:eq:lifetime}{{6.45}{177}{Single-cell results quoted from earlier chapters}{equation.6.appendix.A.45}{}}
\newlabel{p:eq:lifetime@cref}{{[equation][45][6]6.45}{[1][177][]177}}
\newlabel{p:eq:burstlaw}{{6.46}{177}{Distribution-level facts}{equation.6.appendix.A.46}{}}
\newlabel{p:eq:burstlaw@cref}{{[equation][46][6]6.46}{[1][177][]177}}
\newlabel{p:eq:EW2}{{6.47}{177}{Distribution-level facts}{equation.6.appendix.A.47}{}}
\newlabel{p:eq:EW2@cref}{{[equation][47][6]6.47}{[1][177][]177}}
\newlabel{p:eq:gk}{{6.48}{178}{Multiplicity of infection}{equation.6.appendix.A.48}{}}
\newlabel{p:eq:gk@cref}{{[equation][48][6]6.48}{[1][177][]178}}
\newlabel{p:eq:Vk}{{6.49}{178}{Multiplicity of infection}{equation.6.appendix.A.49}{}}
\newlabel{p:eq:Vk@cref}{{[equation][49][6]6.49}{[1][177][]178}}
\newlabel{p:app:transforms}{{6.B}{178}{Laplace transforms of the kernels}{section.6.appendix.B}{}}
\newlabel{p:app:transforms@cref}{{[section][2][6]6.B}{[1][178][]178}}
\newlabel{p:eq:LIhat}{{6.50}{178}{Laplace transforms of the kernels}{equation.6.appendix.B.50}{}}
\newlabel{p:eq:LIhat@cref}{{[equation][50][6]6.50}{[1][178][]178}}
\newlabel{p:eq:Psi}{{6.51}{178}{Laplace transforms of the kernels}{equation.6.appendix.B.51}{}}
\newlabel{p:eq:Psi@cref}{{[equation][51][6]6.51}{[1][178][]178}}
\newlabel{p:eq:dLK}{{6.52}{178}{Laplace transforms of the kernels}{equation.6.appendix.B.52}{}}
\newlabel{p:eq:dLK@cref}{{[equation][52][6]6.52}{[1][178][]178}}
\newlabel{p:app:formulae}{{6.C}{178}{Master formula table}{section.6.appendix.C}{}}
\newlabel{p:app:formulae@cref}{{[section][3][6]6.C}{[1][178][]178}}
\newlabel{p:tab:formulae}{{6.7}{179}{Population-level quantities, in one place}{table.6.7}{}}
\newlabel{p:tab:formulae@cref}{{[table][7][6]6.7}{[1][178][]179}}
\newlabel{tt:ch:two-type}{{7}{180}{Type-specific catastrophe rates}{chapter.7}{}}
\newlabel{tt:ch:two-type@cref}{{[chapter][7][]7}{[1][180][]180}}
\newlabel{sec:introduction}{{7.1}{180}{Introduction}{section.7.1}{}}
\newlabel{sec:introduction@cref}{{[section][1][7]7.1}{[1][180][]180}}
\newlabel{sec:model}{{7.2}{182}{Model, assumptions, and backward equations}{section.7.2}{}}
\newlabel{sec:model@cref}{{[section][2][7]7.2}{[1][181][]182}}
\newlabel{eq:catastrophe-rate}{{7.1}{182}{The two-type process}{equation.7.2.1}{}}
\newlabel{eq:catastrophe-rate@cref}{{[equation][1][7]7.1}{[1][182][]182}}
\newlabel{fig:process-schematic}{{7.1}{183}{Two-type birth--death--conversion with a single global absorbing catastrophe. Type~1 (blue) gives birth at rate $\lambda _1$, dies at $\mu _1$, and converts irreversibly to type~2 (vermillion) at rate $\nu $; type~2 gives birth at $\lambda _2$ and dies at $\mu _2$. In state $(x,y)$ catastrophe fires at rate $\delta _1x+\delta _2y$ and absorbs the process in $(\R ,\R )$. $S(t)$ and $G(t)$ are the probabilities that this has not occurred by time $t$, from one type-1 or type-2 individual respectively. \explorerlink {https://markov.volvox.site/-/adam74aldridge/two-type-bdc}{Simulate}}{figure.7.1}{}}
\newlabel{fig:process-schematic@cref}{{[figure][1][7]7.1}{[1][182][]183}}
\newlabel{prop:feynman-kac}{{7.2}{183}{Occupation-time transform}{theorem.7.2}{}}
\newlabel{prop:feynman-kac@cref}{{[proposition][2][7]7.2}{[1][182][]183}}
\newlabel{eq:feynman-kac}{{7.2}{183}{Occupation-time transform}{equation.7.2.2}{}}
\newlabel{eq:feynman-kac@cref}{{[subequation][2][7]7.2}{[1][182][]183}}
\newlabel{fig:genealogy}{{7.2}{184}{One continuous-time realisation starting from a single type-1 individual. Filled nodes mark births; blue-ringed vermillion nodes mark type~$1\to 2$ conversion; crosses mark ordinary death; red discs mark lineages killed by the global catastrophe (dashed wavefront: simultaneous absorption in $(\R ,\R )$). Parameters: $\lambda _1=1.0$, $\mu _1=0.55$, $\nu =0.40$, $\lambda _2=0.90$, $\mu _2=0.50$, $\delta _1=0.05$, $\delta _2=0.08$ (Gillespie, seed 2869). \explorerlink {https://files.volvox.site/users/pages/6a6852d29e3764e44c80325e/files/1785221941917-TT_bdc_genealogy_visualiser.html}{Genealogy visualiser}}{figure.7.2}{}}
\newlabel{fig:genealogy@cref}{{[figure][2][7]7.2}{[1][182][]184}}
\newlabel{eq:backward-unscaled}{{7.3}{184}{First-event calculation}{equation.7.2.3}{}}
\newlabel{eq:backward-unscaled@cref}{{[subequation][3][7]7.3}{[1][184][]184}}
\newlabel{eq:S-unscaled}{{7.3a}{184}{First-event calculation}{equation.7.2.3a}{}}
\newlabel{eq:S-unscaled@cref}{{[subequation][1][7,3]7.3a}{[1][184][]184}}
\newlabel{eq:G-unscaled}{{7.3b}{184}{First-event calculation}{equation.7.2.3b}{}}
\newlabel{eq:G-unscaled@cref}{{[subequation][2][7,3]7.3b}{[1][184][]184}}
\newlabel{eq:initial-slopes}{{7.4}{184}{First-event calculation}{equation.7.2.4}{}}
\newlabel{eq:initial-slopes@cref}{{[equation][4][7]7.4}{[1][184][]184}}
\newlabel{fig:path-exposure}{{7.3}{185}{Why the no-catastrophe probability depends on occupation times, not only terminal counts. Two deterministic histories on $t\in [0,10]$ share $(X_0,Y_0)=(1,0)$ and $(X_{10},Y_{10})=(2,1)$: (a)~early peak then die-back; (b)~late bloom. (c)~Rates $\lambda =\delta _1X+\delta _2Y$ share the terminal hazard but differ in integrated exposure by about $2.4$. (d)~Pathwise weights $\exp \{-\int \lambda \,ds\}$ averaged in \cref {prop:feynman-kac} ($\delta _1=0.15$, $\delta _2=0.10$)}{figure.7.3}{}}
\newlabel{fig:path-exposure@cref}{{[figure][3][7]7.3}{[1][183][]185}}
\newlabel{fig:first-event}{{7.4}{185}{\relax }{figure.7.4}{}}
\newlabel{fig:first-event@cref}{{[figure][4][7]7.4}{[1][184][]185}}
\newlabel{eq:lambda1-positive}{{7.5}{186}{Dimensionless form}{equation.7.2.5}{}}
\newlabel{eq:lambda1-positive@cref}{{[equation][5][7]7.5}{[1][184][]186}}
\newlabel{eq:scaled-rates}{{7.6}{186}{Dimensionless form}{equation.7.2.6}{}}
\newlabel{eq:scaled-rates@cref}{{[equation][6][7]7.6}{[1][186][]186}}
\newlabel{eq:backward-scaled}{{7.7}{186}{Dimensionless form}{equation.7.2.7}{}}
\newlabel{eq:backward-scaled@cref}{{[subequation][7][7]7.7}{[1][186][]186}}
\newlabel{eq:S-scaled}{{7.7a}{186}{Dimensionless form}{equation.7.2.7a}{}}
\newlabel{eq:S-scaled@cref}{{[subequation][1][7,7]7.7a}{[1][186][]186}}
\newlabel{eq:G-scaled}{{7.7b}{186}{Dimensionless form}{equation.7.2.7b}{}}
\newlabel{eq:G-scaled@cref}{{[subequation][2][7,7]7.7b}{[1][186][]186}}
\newlabel{sec:autonomous}{{7.3}{186}{The autonomous type-2 probability}{section.7.3}{}}
\newlabel{sec:autonomous@cref}{{[section][3][7]7.3}{[1][186][]186}}
\newlabel{eq:Delta2}{{7.8}{186}{The autonomous type-2 probability}{equation.7.3.8}{}}
\newlabel{eq:Delta2@cref}{{[equation][8][7]7.8}{[1][186][]186}}
\newlabel{eq:type2-roots}{{7.9}{186}{The autonomous type-2 probability}{equation.7.3.9}{}}
\newlabel{eq:type2-roots@cref}{{[equation][9][7]7.9}{[1][186][]186}}
\newlabel{eq:G-factored}{{7.10}{186}{The autonomous type-2 probability}{equation.7.3.10}{}}
\newlabel{eq:G-factored@cref}{{[equation][10][7]7.10}{[1][186][]186}}
\newlabel{fig:type2-phase}{{7.5}{187}{Geometry of the autonomous type-2 probability. (a)~Right-hand side $f(G)=\lambda _2(G-r_{2,-})(G-r_{2,+})$, with $0\le r_{2,-}<1<r_{2,+}$ so only $r_{2,-}$ is an admissible long-time probability. (b)~Phase line: $r_{2,-}$ attracting, $r_{2,+}$ repelling; $G(0)=1$ is drawn into $r_{2,-}$. (c)~Relaxation of $G(t)$ from $1$ toward $r_{2,-}$. Parameters: $\lambda _2=0.9$, $\mu _2=0.35$, $\delta _2=0.20$}{figure.7.5}{}}
\newlabel{fig:type2-phase@cref}{{[figure][5][7]7.5}{[1][186][]187}}
\newlabel{eq:type2-alpha-q}{{7.11}{187}{The autonomous type-2 probability}{equation.7.3.11}{}}
\newlabel{eq:type2-alpha-q@cref}{{[equation][11][7]7.11}{[1][186][]187}}
\newlabel{prop:G-solution}{{7.3}{187}{Exact type-2 probability}{theorem.7.3}{}}
\newlabel{prop:G-solution@cref}{{[proposition][3][7]7.3}{[1][186][]187}}
\newlabel{eq:type2-z}{{7.12}{187}{Exact type-2 probability}{equation.7.3.12}{}}
\newlabel{eq:type2-z@cref}{{[equation][12][7]7.12}{[1][186][]187}}
\newlabel{eq:G-solution}{{7.13}{187}{Exact type-2 probability}{equation.7.3.13}{}}
\newlabel{eq:G-solution@cref}{{[equation][13][7]7.13}{[1][187][]187}}
\newlabel{fig:G-and-z2}{{7.6}{188}{(a)~$G(t)$ from $1$ toward $r_{2,-}$. (b)~Coordinate $z_2(\lambda _1t)$ increasing from $z_{2,0}<0$ toward $0$. (c)~Linear-fractional map $G=r_{2,-}-\alpha _2 z_2/(1-z_2)$. Parameters: $\lambda _1=1$, $\lambda _2=0.85$, $\mu _2=0.40$, $\delta _2=0.18$}{figure.7.6}{}}
\newlabel{fig:G-and-z2@cref}{{[figure][6][7]7.6}{[1][188][]188}}
\newlabel{sec:main_result}{{7.4}{188}{Hypergeometric reduction and main result}{section.7.4}{}}
\newlabel{sec:main_result@cref}{{[section][4][7]7.4}{[1][188][]188}}
\newlabel{fig:derivation-pipeline}{{7.7}{189}{\relax }{figure.7.7}{}}
\newlabel{fig:derivation-pipeline@cref}{{[figure][7][7]7.7}{[1][188][]189}}
\newlabel{eq:h-quadratic}{{7.14}{189}{Selecting the stable constant part}{equation.7.4.14}{}}
\newlabel{eq:h-quadratic@cref}{{[equation][14][7]7.14}{[1][189][]189}}
\newlabel{eq:h-lower-root}{{7.15}{189}{Selecting the stable constant part}{equation.7.4.15}{}}
\newlabel{eq:h-lower-root@cref}{{[equation][15][7]7.15}{[1][189][]189}}
\newlabel{eq:X-riccati}{{7.16}{189}{Selecting the stable constant part}{equation.7.4.16}{}}
\newlabel{eq:X-riccati@cref}{{[equation][16][7]7.16}{[1][189][]189}}
\newlabel{sec:linearising}{{7.4.2}{190}{Linearising the Riccati equation}{subsection.7.4.2}{}}
\newlabel{sec:linearising@cref}{{[subsection][2][7,4]7.4.2}{[1][189][]190}}
\newlabel{eq:riccati-substitution}{{7.17}{190}{Linearising the Riccati equation}{equation.7.4.17}{}}
\newlabel{eq:riccati-substitution@cref}{{[equation][17][7]7.17}{[1][189][]190}}
\newlabel{eq:Z-linear}{{7.18}{190}{Linearising the Riccati equation}{equation.7.4.18}{}}
\newlabel{eq:Z-linear@cref}{{[equation][18][7]7.18}{[1][190][]190}}
\newlabel{eq:chain-Z}{{7.19}{190}{Change of independent variable}{equation.7.4.19}{}}
\newlabel{eq:chain-Z@cref}{{[equation][19][7]7.19}{[1][190][]190}}
\newlabel{eq:chain-Z2}{{7.20}{190}{Change of independent variable}{equation.7.4.20}{}}
\newlabel{eq:chain-Z2@cref}{{[equation][20][7]7.20}{[1][190][]190}}
\newlabel{eq:hypergeometric}{{7.21}{190}{Change of independent variable}{equation.7.4.21}{}}
\newlabel{eq:hypergeometric@cref}{{[equation][21][7]7.21}{[1][190][]190}}
\newlabel{eq:C-generalized}{{7.22}{190}{Change of independent variable}{equation.7.4.22}{}}
\newlabel{eq:C-generalized@cref}{{[equation][22][7]7.22}{[1][190][]190}}
\newlabel{eq:gauss-canonical-solutions}{{7.23}{191}{Matching to the Gauss equation}{equation.7.4.23}{}}
\newlabel{eq:gauss-canonical-solutions@cref}{{[subequation][23][7]7.23}{[1][190][]191}}
\newlabel{eq:ABC-relations}{{7.24}{191}{Matching to the Gauss equation}{equation.7.4.24}{}}
\newlabel{eq:ABC-relations@cref}{{[equation][24][7]7.24}{[1][191][]191}}
\newlabel{eq:AB-explicit}{{7.25}{191}{Matching to the Gauss equation}{equation.7.4.25}{}}
\newlabel{eq:AB-explicit@cref}{{[equation][25][7]7.25}{[1][191][]191}}
\newlabel{sec:closed-formula}{{7.4.3}{191}{The closed formula}{subsection.7.4.3}{}}
\newlabel{sec:closed-formula@cref}{{[subsection][3][7,4]7.4.3}{[1][191][]191}}
\newlabel{eq:S-from-Phi}{{7.26}{191}{Rebuilding $\widehat S$ from $\Phi $}{equation.7.4.26}{}}
\newlabel{eq:S-from-Phi@cref}{{[equation][26][7]7.26}{[1][191][]191}}
\newlabel{eq:UV-definitions}{{7.27}{191}{Ansatz for $\Phi $: the two-dimensional solution space}{equation.7.4.27}{}}
\newlabel{eq:UV-definitions@cref}{{[subequation][27][7]7.27}{[1][191][]191}}
\newlabel{eq:L-definition}{{7.28}{192}{Initial condition and the constant $L$}{equation.7.4.28}{}}
\newlabel{eq:L-definition@cref}{{[equation][28][7]7.28}{[1][192][]192}}
\newlabel{eq:Phi-linear-combination}{{7.29}{192}{Choosing the mix: coefficients $K_U$ and $K_V$}{equation.7.4.29}{}}
\newlabel{eq:Phi-linear-combination@cref}{{[equation][29][7]7.29}{[1][192][]192}}
\newlabel{eq:K-coefficients}{{7.30}{192}{Choosing the mix: coefficients $K_U$ and $K_V$}{equation.7.4.30}{}}
\newlabel{eq:K-coefficients@cref}{{[subequation][30][7]7.30}{[1][192][]192}}
\newlabel{tt:thm:main}{{7.4}{193}{Exact finite-time no-catastrophe probability}{theorem.7.4}{}}
\newlabel{tt:thm:main@cref}{{[theorem][4][7]7.4}{[1][193][]193}}
\newlabel{eq:S-final}{{7.31}{193}{Exact finite-time no-catastrophe probability}{equation.7.4.31}{}}
\newlabel{eq:S-final@cref}{{[equation][31][7]7.31}{[1][193][]193}}
\newlabel{eq:unscaled-recovery}{{7.32}{193}{Exact finite-time no-catastrophe probability}{equation.7.4.32}{}}
\newlabel{eq:unscaled-recovery@cref}{{[equation][32][7]7.32}{[1][193][]193}}
\newlabel{fig:bdc-survival}{{7.8}{194}{No-catastrophe probabilities $S(t)$ (type~1 start) and $G(t)$ (type~2 start) at fixed rates $\lambda _1=1.08$, $\mu _1=0.2$, $\nu =0.56$, $\delta _1=0.013$, $\lambda _2=1.13$, $\mu _2=0.93$, $\delta _2=0.611$. Asymptotic levels $h=\lim S$ and $r_{2,-}=\lim G$ are indicated on the panel. Solid curves: closed form of \cref {tt:thm:main}; open circles: RK4 integration of the unscaled backward ODEs. Rates are labelled on the panel. \explorerlink {https://files.volvox.site/users/pages/6a6852d29e3764e44c80325e/files/1785222252388-TT_bdc_survival_visualiser.html}{Survival visualiser}}{figure.7.8}{}}
\newlabel{fig:bdc-survival@cref}{{[figure][8][7]7.8}{[1][193][]194}}
\newlabel{fig:bdc-density}{{7.9}{194}{Defective catastrophe-time densities $-S'(t)$ and $-G'(t)$ for the same rates as \cref {fig:bdc-survival}. Solid curves from the closed form; open circles from the same RK4 check of the backward ODEs}{figure.7.9}{}}
\newlabel{fig:bdc-density@cref}{{[figure][9][7]7.9}{[1][193][]194}}
\newlabel{note:parameter-map}{{7.6}{194}{Nested parameters}{theorem.7.6}{}}
\newlabel{note:parameter-map@cref}{{[remark][6][7]7.6}{[1][194][]194}}
\newlabel{sec:interpretation}{{7.4.4}{195}{Mathematical interpretation}{subsection.7.4.4}{}}
\newlabel{sec:interpretation@cref}{{[subsection][4][7,4]7.4.4}{[1][195][]195}}
\newlabel{sec:biology}{{7.5}{195}{Intracellular application: comparing catastrophe-rate regimes}{section.7.5}{}}
\newlabel{sec:biology@cref}{{[section][5][7]7.5}{[1][195][]195}}
\newlabel{tab:bio-mapping}{{7.1}{196}{Reading of mathematical components in an early intracellular infection setting}{table.7.1}{}}
\newlabel{tab:bio-mapping@cref}{{[table][1][7]7.1}{[1][195][]196}}
\newlabel{fig:bio-regimes}{{7.10}{198}{Competing catastrophe-rate regimes under an intracellular \textit {Y.\ pestis} reading. (A)~Survival $S(t)$ from one early founder; (B)~conditional intensity $q_S=-S'/S$. Shared growth rates with $\delta _1+\delta _2=c=0.30$ (coefficient normalisation, not matched realised exposure): EQ $(0.15,0.15)$, MAT $(0.05,0.25)$, GATE $(0,0.30)$, EARLY $(0.25,0.05)$ ($\lambda _1=1$, $\mu _1=0.45$, $\nu =0.35$, $\lambda _2=0.85$, $\mu _2=0.40$). Parameters are illustrative, not fitted. MAT is a working hypothesis; GATE is the type-2-only extreme ($S'(0)=0$). Exact closed form for $S$; $q_S$ from the ODE on the solution path}{figure.7.10}{}}
\newlabel{fig:bio-regimes@cref}{{[figure][10][7]7.10}{[1][198][]198}}
\newlabel{sec:conclusion}{{7.6}{198}{Conclusion}{section.7.6}{}}
\newlabel{sec:conclusion@cref}{{[section][6][7]7.6}{[1][198][]198}}
\newlabel{app:product-to-exponential}{{7.A}{199}{From product of survival factors to the pathwise exponential}{section.7.appendix.A}{}}
\newlabel{app:product-to-exponential@cref}{{[section][1][7]7.A}{[1][199][]199}}
\newlabel{app:hypergeometric-details}{{7.B}{199}{Hypergeometric derivatives and resonant parameters}{section.7.appendix.B}{}}
\newlabel{app:hypergeometric-details@cref}{{[section][2][7]7.B}{[1][199][]199}}
\newlabel{eq:U-derivative}{{7.33}{199}{Hypergeometric derivatives and resonant parameters}{equation.7.appendix.B.33}{}}
\newlabel{eq:U-derivative@cref}{{[equation][33][7]7.33}{[1][199][]199}}
\newlabel{eq:V-derivative}{{7.34}{199}{Hypergeometric derivatives and resonant parameters}{equation.7.appendix.B.34}{}}
\newlabel{eq:V-derivative@cref}{{[equation][34][7]7.34}{[1][199][]199}}
\newlabel{app:limits}{{7.C}{200}{Checks, limits, and the earlier branching model}{section.7.appendix.C}{}}
\newlabel{app:limits@cref}{{[section][3][7]7.C}{[1][199][]200}}
\newlabel{cor:long-time}{{7.7}{200}{Avoiding catastrophe forever}{theorem.7.7}{}}
\newlabel{cor:long-time@cref}{{[corollary][7][7]7.7}{[1][200][]200}}
\newlabel{eq:equal-catastrophe-diagonal}{{7.35}{200}{Equal catastrophe rates}{equation.7.appendix.C.35}{}}
\newlabel{eq:equal-catastrophe-diagonal@cref}{{[equation][35][7]7.35}{[1][200][]200}}
\newlabel{sec:delta2-zero}{{7.C.3}{200}{Type 2 cannot trigger catastrophe}{subsection.7.appendix.C.3}{}}
\newlabel{sec:delta2-zero@cref}{{[subsection][3][7,3]7.C.3}{[1][200][]200}}
\newlabel{eq:S-delta2-zero-ode}{{7.36}{200}{Type 2 cannot trigger catastrophe}{equation.7.appendix.C.36}{}}
\newlabel{eq:S-delta2-zero-ode@cref}{{[equation][36][7]7.36}{[1][200][]200}}
\newlabel{eq:R0}{{7.37}{201}{Type 2 cannot trigger catastrophe}{equation.7.appendix.C.37}{}}
\newlabel{eq:R0@cref}{{[equation][37][7]7.37}{[1][200][]201}}
\newlabel{prop:delta2-zero}{{7.8}{201}{Elementary type-2-zero boundary}{theorem.7.8}{}}
\newlabel{prop:delta2-zero@cref}{{[proposition][8][7]7.8}{[1][201][]201}}
\newlabel{eq:no-catastrophe-trivial}{{7.38}{201}{Type 2 cannot trigger catastrophe}{equation.7.appendix.C.38}{}}
\newlabel{eq:no-catastrophe-trivial@cref}{{[equation][38][7]7.38}{[1][201][]201}}
\newlabel{sec:lambda1-zero}{{7.C.4}{201}{No type-1 births}{subsection.7.appendix.C.4}{}}
\newlabel{sec:lambda1-zero@cref}{{[subsection][4][7,3]7.C.4}{[1][201][]201}}
\newlabel{eq:S-lambda1-zero}{{7.39}{201}{No type-1 births}{equation.7.appendix.C.39}{}}
\newlabel{eq:S-lambda1-zero@cref}{{[equation][39][7]7.39}{[1][201][]201}}
\newlabel{sec:lambda2-zero}{{7.C.5}{201}{No type-2 births}{subsection.7.appendix.C.5}{}}
\newlabel{sec:lambda2-zero@cref}{{[subsection][5][7,3]7.C.5}{[1][201][]201}}
\newlabel{eq:G-lambda2-zero}{{7.40}{201}{No type-2 births}{equation.7.appendix.C.40}{}}
\newlabel{eq:G-lambda2-zero@cref}{{[equation][40][7]7.40}{[1][201][]201}}
\newlabel{eq:bessel-initial-data}{{7.41}{202}{No type-2 births}{equation.7.appendix.C.41}{}}
\newlabel{eq:bessel-initial-data@cref}{{[equation][41][7]7.41}{[1][201][]202}}
\newlabel{prop:lambda2-zero}{{7.9}{202}{Zero type-2 birth boundary}{theorem.7.9}{}}
\newlabel{prop:lambda2-zero@cref}{{[proposition][9][7]7.9}{[1][202][]202}}
\newlabel{eq:S-lambda2-zero}{{7.42}{202}{Zero type-2 birth boundary}{equation.7.appendix.C.42}{}}
\newlabel{eq:S-lambda2-zero@cref}{{[equation][42][7]7.42}{[1][202][]202}}
\newlabel{fig:ak-contrast}{{7.11}{203}{Antal--Krapivsky terminal-population generating functions (left) versus the no-catastrophe probabilities studied here (right). Both use the same two-type one-way branching skeleton. Left: $\mathcal A(x,y,t)=\E [x^{X_t}y^{Y_t}]$ with initial data $\mathcal A(x,y,0)=x$, $\mathcal B(x,y,0)=y$, and no global absorbing catastrophe. Right: $S(t)=\Prob _{(1,0)}(\tau _{\!c}>t)$ and $G(t)=\Prob _{(0,1)}(\tau _{\!c}>t)$ with $S(0)=G(0)=1$ and catastrophe rate $\delta _1x+\delta _2y$}{figure.7.11}{}}
\newlabel{fig:ak-contrast@cref}{{[figure][11][7]7.11}{[1][202][]203}}
\newlabel{app:AK-recovery}{{7.D}{203}{Equation-level recovery of Antal--Krapivsky}{section.7.appendix.D}{}}
\newlabel{app:AK-recovery@cref}{{[section][4][7]7.D}{[1][202][]203}}
\newlabel{sec:AK-recovery}{{7.D}{203}{Equation-level recovery of Antal--Krapivsky}{section.7.appendix.D}{}}
\newlabel{sec:AK-recovery@cref}{{[section][4][7]7.D}{[1][202][]203}}
\newlabel{app:catastrophe-pgf}{{7.E}{203}{A catastrophe-aware population generating function}{section.7.appendix.E}{}}
\newlabel{app:catastrophe-pgf@cref}{{[section][5][7]7.E}{[1][203][]203}}
\newlabel{eq:killed-pgf-definitions}{{7.43}{203}{A catastrophe-aware population generating function}{equation.7.appendix.E.43}{}}
\newlabel{eq:killed-pgf-definitions@cref}{{[subequation][43][7]7.43}{[1][203][]203}}
\newlabel{prop:killed-pgf-backward}{{7.10}{204}{Backward system for the killed generating functions}{theorem.7.10}{}}
\newlabel{prop:killed-pgf-backward@cref}{{[proposition][10][7]7.10}{[1][204][]204}}
\newlabel{eq:killed-pgf-backward}{{7.44}{204}{Backward system for the killed generating functions}{equation.7.appendix.E.44}{}}
\newlabel{eq:killed-pgf-backward@cref}{{[subequation][44][7]7.44}{[1][204][]204}}
\newlabel{eq:killed-pgf-initial}{{7.45}{204}{Backward system for the killed generating functions}{equation.7.appendix.E.45}{}}
\newlabel{eq:killed-pgf-initial@cref}{{[equation][45][7]7.45}{[1][204][]204}}
\newlabel{eq:extended-pgf}{{7.46}{204}{The absorbing state and conditional populations}{equation.7.appendix.E.46}{}}
\newlabel{eq:extended-pgf@cref}{{[equation][46][7]7.46}{[1][204][]204}}
\newlabel{eq:conditional-pgf}{{7.47}{204}{The absorbing state and conditional populations}{equation.7.appendix.E.47}{}}
\newlabel{eq:conditional-pgf@cref}{{[equation][47][7]7.47}{[1][204][]204}}
\newlabel{eq:pgf-z-coordinate}{{7.48}{205}{Reuse of the hypergeometric solution}{equation.7.appendix.E.48}{}}
\newlabel{eq:pgf-z-coordinate@cref}{{[equation][48][7]7.48}{[1][205][]205}}
\newlabel{eq:pgf-F2-solution}{{7.49}{205}{Reuse of the hypergeometric solution}{equation.7.appendix.E.49}{}}
\newlabel{eq:pgf-F2-solution@cref}{{[equation][49][7]7.49}{[1][205][]205}}
\newlabel{eq:pgf-F1-solution}{{7.50}{205}{Reuse of the hypergeometric solution}{equation.7.appendix.E.50}{}}
\newlabel{eq:pgf-F1-solution@cref}{{[equation][50][7]7.50}{[1][205][]205}}
\newlabel{sec:three-pgf-distributions}{{7.E.3}{205}{Three distributions encoded by the killed PGF}{subsection.7.appendix.E.3}{}}
\newlabel{sec:three-pgf-distributions@cref}{{[subsection][3][7,5]7.E.3}{[1][205][]205}}
\newlabel{sec:ultimate-release}{{7.E.4}{206}{Ultimate release counts}{subsection.7.appendix.E.4}{}}
\newlabel{sec:ultimate-release@cref}{{[subsection][4][7,5]7.E.4}{[1][206][]206}}
\newlabel{eq:release-def}{{7.51}{206}{Ultimate release counts}{equation.7.appendix.E.51}{}}
\newlabel{eq:release-def@cref}{{[equation][51][7]7.51}{[1][206][]206}}
\newlabel{fig:dist-at-catastrophe}{{7.12a}{207}{At catastrophe}{subfigure.7.12.1}{}}
\newlabel{fig:dist-at-catastrophe@cref}{{[subfigure][1][7,12]7.12a}{[1][206][]207}}
\newlabel{sub@fig:dist-at-catastrophe}{{a}{207}{At catastrophe}{subfigure.7.12.1}{}}
\newlabel{sub@fig:dist-at-catastrophe@cref}{{[subfigure][1][7,12]7.12a}{[1][206][]207}}
\newlabel{fig:dist-no-catastrophe}{{7.12b}{207}{No catastrophe}{subfigure.7.12.2}{}}
\newlabel{fig:dist-no-catastrophe@cref}{{[subfigure][2][7,12]7.12b}{[1][206][]207}}
\newlabel{sub@fig:dist-no-catastrophe}{{b}{207}{No catastrophe}{subfigure.7.12.2}{}}
\newlabel{sub@fig:dist-no-catastrophe@cref}{{[subfigure][2][7,12]7.12b}{[1][206][]207}}
\newlabel{fig:dist-active}{{7.12c}{207}{Active}{subfigure.7.12.3}{}}
\newlabel{fig:dist-active@cref}{{[subfigure][3][7,12]7.12c}{[1][206][]207}}
\newlabel{sub@fig:dist-active}{{c}{207}{Active}{subfigure.7.12.3}{}}
\newlabel{sub@fig:dist-active@cref}{{[subfigure][3][7,12]7.12c}{[1][206][]207}}
\newlabel{fig:three-pgf-distributions}{{7.12}{207}{Three distributions encoded by the killed PGF at fixed time $t=3$, from one type-1 individual. (a)~Pre-catastrophe counts conditional on catastrophe at time $t$ (density sense). (b)~Counts conditional on no catastrophe up to $t$ (includes ordinary extinction at the origin). (c)~Active counts conditional on no catastrophe and non-extinction (log scale). Rates: $\lambda _1=0.75$, $\mu _1=0.3$, $\nu =0.7$, $\delta _1=0.001$, $\lambda _2=0.75$, $\mu _2=0.25$, $\delta _2=0.083$. Colour intensity scales are independent across panels. \explorerlink {https://files.volvox.site/users/pages/6a6852d29e3764e44c80325e/files/1785222034060-TT_bdc_distribution_visualiser.html}{Distribution visualiser}}{figure.7.12}{}}
\newlabel{fig:three-pgf-distributions@cref}{{[figure][12][7]7.12}{[1][206][]207}}
\newlabel{rem:release-origin}{{7.11}{207}{Generic regime versus boundaries}{theorem.7.11}{}}
\newlabel{rem:release-origin@cref}{{[remark][11][7]7.11}{[1][206][]207}}
\newlabel{eq:release-joint-pmf}{{7.52}{207}{Ultimate release counts}{equation.7.appendix.E.52}{}}
\newlabel{eq:release-joint-pmf@cref}{{[equation][52][7]7.52}{[1][207][]207}}
\newlabel{eq:release-pgf}{{7.53}{207}{Ultimate release counts}{equation.7.appendix.E.53}{}}
\newlabel{eq:release-pgf@cref}{{[equation][53][7]7.53}{[1][207][]207}}
\newlabel{eq:release-conditional}{{7.54}{207}{Ultimate release counts}{equation.7.appendix.E.54}{}}
\newlabel{eq:release-conditional@cref}{{[equation][54][7]7.54}{[1][207][]207}}
\newlabel{eq:release-mixture}{{7.55}{207}{Ultimate release counts}{equation.7.appendix.E.55}{}}
\newlabel{eq:release-mixture@cref}{{[equation][55][7]7.55}{[1][207][]207}}
\newlabel{fig:release-joint}{{7.13a}{208}{Joint law}{subfigure.7.13.1}{}}
\newlabel{fig:release-joint@cref}{{[subfigure][1][7,13]7.13a}{[1][208][]208}}
\newlabel{sub@fig:release-joint}{{a}{208}{Joint law}{subfigure.7.13.1}{}}
\newlabel{sub@fig:release-joint@cref}{{[subfigure][1][7,13]7.13a}{[1][208][]208}}
\newlabel{fig:release-marginals}{{7.13b}{208}{Marginals}{subfigure.7.13.2}{}}
\newlabel{fig:release-marginals@cref}{{[subfigure][2][7,13]7.13b}{[1][208][]208}}
\newlabel{sub@fig:release-marginals}{{b}{208}{Marginals}{subfigure.7.13.2}{}}
\newlabel{sub@fig:release-marginals@cref}{{[subfigure][2][7,13]7.13b}{[1][208][]208}}
\newlabel{fig:release-law}{{7.13}{208}{Unconditional ultimate release law $(X_{\mathrm {rel}},Y_{\mathrm {rel}})$ from one type-1 individual. (a)~Joint masses $\Prob _{(1,0)}\{X_{\mathrm {rel}}=m,\,Y_{\mathrm {rel}}=n\}$; the origin atom is forever-avoidance $\sigma _1(\infty )=h$ (clearance without release in the generic regime). (b)~Marginals $\Prob \{X_{\mathrm {rel}}=k\}$ (blue) and $\Prob \{Y_{\mathrm {rel}}=k\}$ (vermillion). Rates (on panel~(a)): $\lambda _1=0.56$, $\mu _1=0$, $\nu =0.12$, $\delta _1=0.012$, $\lambda _2=0.74$, $\mu _2=0.06$, $\delta _2=0.623$. Lattice truncation of the killed master equations; residual active mass at the numerical horizon is negligible. \explorerlink {https://files.volvox.site/users/pages/6a6852d29e3764e44c80325e/files/1785222143569-TT_bdc_release_visualiser.html}{Release visualiser}}{figure.7.13}{}}
\newlabel{fig:release-law@cref}{{[figure][13][7]7.13}{[1][208][]208}}
\newlabel{eq:release-marg-X}{{7.56}{208}{Ultimate release counts}{equation.7.appendix.E.56}{}}
\newlabel{eq:release-marg-X@cref}{{[equation][56][7]7.56}{[1][208][]208}}
\newlabel{eq:release-marg-Y}{{7.57}{208}{Ultimate release counts}{equation.7.appendix.E.57}{}}
\newlabel{eq:release-marg-Y@cref}{{[equation][57][7]7.57}{[1][208][]208}}
\newlabel{path:ch:pathogenesis}{{8}{209}{Bursting and budding pathogenesis}{chapter.8}{}}
\newlabel{path:ch:pathogenesis@cref}{{[chapter][8][]8}{[1][209][]209}}
\newlabel{path:sec:introduction}{{8.1}{209}{Introduction}{section.8.1}{}}
\newlabel{path:sec:introduction@cref}{{[section][1][8]8.1}{[1][209][]209}}
\newlabel{path:sec:advantages}{{8.1.1}{209}{Three candidate advantages}{subsection.8.1.1}{}}
\newlabel{path:sec:advantages@cref}{{[subsection][1][8,1]8.1.1}{[1][209][]209}}
\newlabel{path:tab:notation}{{8.1}{211}{Notation. The upper block is taken from \ChCore {} and \ChPop {}; the lower block carries meanings specific to this chapter, and those letters were chosen to avoid collisions with the ones those meanings would otherwise have taken}{table.8.1}{}}
\newlabel{path:tab:notation@cref}{{[table][1][8]8.1}{[1][210][]211}}
\newlabel{path:sec:schedule}{{8.2}{211}{Why the release schedule alone cannot matter}{section.8.2}{}}
\newlabel{path:sec:schedule@cref}{{[section][2][8]8.2}{[1][211][]211}}
\newlabel{path:sec:chainedruptures}{{8.2.1}{211}{Successive birth--death--catastrophe}{subsection.8.2.1}{}}
\newlabel{path:sec:chainedruptures@cref}{{[subsection][1][8,2]8.2.1}{[1][211][]211}}
\newlabel{path:fig:basic}{{8.1}{212}{Four frames of a particle simulation, showing what the chained model leaves out. (a) Intracellular replication before bursting: the load accumulates out of reach. (b) The cell bursts and the contents span the extracellular space before finding new hosts --- the interval the chained model collapses to a point. (c) Macrophages are motile; they wander the extracellular matrix phagocytosing debris and extracellular pathogens alike, so which cell a released agent reaches is not settled by diffusion alone. (d) A second round of bursting, here with a larger load. Under identical intracellular rates the growth rate is the same however the transfers between (a) and (d) are arranged, which is why the argument of this section forces a nonlinearity into the model. \explorerlink {https://files.volvox.site/users/69b423f5faa2bfbb5b696c61/pages/6a8495d018f4e12409b5ecf1/files/1789007285465-fau0qe-burst_suppression.html}{Simulate}}{figure.8.1}{}}
\newlabel{path:fig:basic@cref}{{[figure][1][8]8.1}{[1][212][]212}}
\newlabel{path:sec:intracellular}{{8.3}{213}{Intracellular depletion: replicator and suppressor}{section.8.3}{}}
\newlabel{path:sec:intracellular@cref}{{[section][3][8]8.3}{[1][213][]213}}
\newlabel{path:sec:granules}{{8.3.1}{213}{The granule store}{subsection.8.3.1}{}}
\newlabel{path:sec:granules@cref}{{[subsection][1][8,3]8.3.1}{[1][213][]213}}
\newlabel{path:sec:rsmodel}{{8.3.2}{213}{The replicator--suppressor process}{subsection.8.3.2}{}}
\newlabel{path:sec:rsmodel@cref}{{[subsection][2][8,3]8.3.2}{[1][213][]213}}
\newlabel{path:fig:rs}{{8.2}{214}{The two opposing forces, simulated. Intracellular granules are in red, pathogenic particles in green; when one of each makes contact, both are annihilated, and that mutual annihilation is the depletion the model turns on. (a) Early, with the store still nearly full. (b) Later in the same run, with most of the store spent: the pathogen count is little changed, but the suppression rate $\varsigma i j$ has fallen with $j$, so what remains of the defence can no longer keep pace with the birth rate. \explorerlink {https://files.volvox.site/users/69b423f5faa2bfbb5b696c61/pages/6a8495d018f4e12409b5ecf1/files/1789007285465-fau0qe-burst_suppression.html}{Simulate}}{figure.8.2}{}}
\newlabel{path:fig:rs@cref}{{[figure][2][8]8.2}{[1][213][]214}}
\newlabel{path:def:rs}{{8.1}{214}{Replicator--suppressor process}{theorem.8.1}{}}
\newlabel{path:def:rs@cref}{{[definition][1][8]8.1}{[1][214][]214}}
\newlabel{path:eq:statespace}{{8.1}{214}{Replicator--suppressor process}{equation.8.3.1}{}}
\newlabel{path:eq:statespace@cref}{{[equation][1][8]8.1}{[1][214][]214}}
\newlabel{path:eq:rsrates}{{8.2}{214}{Replicator--suppressor process}{equation.8.3.2}{}}
\newlabel{path:eq:rsrates@cref}{{[equation][2][8]8.2}{[1][214][]214}}
\newlabel{path:fig:statespace}{{8.3}{214}{The state space $\mathcal {S}$ of \cref {path:def:rs}. Births move right and spend no granule; suppressions move down and left, spending one of each. The column $X_t=0$ is extinction; the row $Q_t=0$ is a spent store, from which only births remain, so survival is certain. That is why the diagonal $i=j$ matters}{figure.8.3}{}}
\newlabel{path:fig:statespace@cref}{{[figure][3][8]8.3}{[1][214][]214}}
\newlabel{path:eq:survival}{{8.3}{215}{The replicator--suppressor process}{equation.8.3.3}{}}
\newlabel{path:eq:survival@cref}{{[equation][3][8]8.3}{[1][215][]215}}
\newlabel{path:rem:noreplenish}{{8.2}{215}{No replenishment}{theorem.8.2}{}}
\newlabel{path:rem:noreplenish@cref}{{[remark][2][8]8.2}{[1][215][]215}}
\newlabel{path:fig:fate}{{8.4}{215}{The three fates, read off \cref {path:fig:statespace}. In $i>j$ survival is certain: extinction would need $i$ suppressions and only $j$ are available (\cref {path:prop:certainty}). In $1\le i\le j$ both outcomes have positive probability, and $\pi _{i,j}$ is what the rest of this section computes. The boundary is the diagonal, and it is sharp}{figure.8.4}{}}
\newlabel{path:fig:fate@cref}{{[figure][4][8]8.4}{[1][215][]215}}
\newlabel{path:sec:rsfirststep}{{8.3.3}{215}{First-step analysis}{subsection.8.3.3}{}}
\newlabel{path:sec:rsfirststep@cref}{{[subsection][3][8,3]8.3.3}{[1][215][]215}}
\newlabel{path:eq:onestep}{{8.4}{216}{First-step analysis}{equation.8.3.4}{}}
\newlabel{path:eq:onestep@cref}{{[equation][4][8]8.4}{[1][215][]216}}
\newlabel{path:eq:firststep}{{8.5}{216}{First-step analysis}{equation.8.3.5}{}}
\newlabel{path:eq:firststep@cref}{{[equation][5][8]8.5}{[1][216][]216}}
\newlabel{path:eq:matrixsystem}{{8.6}{216}{First-step analysis}{equation.8.3.6}{}}
\newlabel{path:eq:matrixsystem@cref}{{[equation][6][8]8.6}{[1][216][]216}}
\newlabel{path:fig:levels}{{8.5}{216}{The level structure of the first-step system \cref {path:eq:firststep}. States are grouped by the level $j-i$. A suppression stays on its diagonal; a birth moves one level down. No transition raises the level, so $\mathbf {A}$ in \cref {path:eq:matrixsystem} is block lower-triangular in this ordering and the levels may be solved in turn, starting from level~$0$}{figure.8.5}{}}
\newlabel{path:fig:levels@cref}{{[figure][5][8]8.5}{[1][216][]216}}
\newlabel{path:prop:certainty}{{8.3}{216}{Certainty threshold}{theorem.8.3}{}}
\newlabel{path:prop:certainty@cref}{{[proposition][3][8]8.3}{[1][216][]216}}
\newlabel{path:tab:diagonal}{{8.2}{217}{The diagonal law \cref {path:eq:diagonal} at $\lambda =1$. Reading across a row shows the decay in cohort size at fixed rates; reading down a column shows how much slower that decay becomes as suppression strengthens relative to birth. The bottom row is the borderline case $\lambda =\varsigma $, where the failure probability is exactly $1/(j+1)$ --- a cohort of twenty-five matched to twenty-five granules still fails one time in twenty-six}{table.8.2}{}}
\newlabel{path:tab:diagonal@cref}{{[table][2][8]8.2}{[1][217][]217}}
\newlabel{path:prop:diagonal}{{8.4}{217}{The diagonal law}{theorem.8.4}{}}
\newlabel{path:prop:diagonal@cref}{{[proposition][4][8]8.4}{[1][217][]217}}
\newlabel{path:eq:diagonal}{{8.7}{217}{The diagonal law}{equation.8.3.7}{}}
\newlabel{path:eq:diagonal@cref}{{[equation][7][8]8.7}{[1][217][]217}}
\newlabel{path:sec:rsnumerics}{{8.3.4}{217}{The survival maps}{subsection.8.3.4}{}}
\newlabel{path:sec:rsnumerics@cref}{{[subsection][4][8,3]8.3.4}{[1][217][]217}}
\newlabel{path:fig:repsupp20}{{8.6}{218}{Survival probability $1-\pi _{i,j}$ at $\lambda =1$ and $N=20$, for $\varsigma =0.1$, $0.2$ and $0.3$ (left to right). Each lower panel is the full grid of initial conditions, with the dashed line marking the diagonal $i=j$; each upper panel is the marginal at a full store, $1-\pi _{i,N}$, which is the case a freshly engulfing macrophage presents. Everything below the diagonal is certain survival (\cref {path:prop:certainty}); the transition above it steepens as $\varsigma $ falls, so that a weakly armed cell is not merely less effective but less discriminating about dose}{figure.8.6}{}}
\newlabel{path:fig:repsupp20@cref}{{[figure][6][8]8.6}{[1][217][]218}}
\newlabel{path:fig:repsupp10}{{8.7}{218}{The same quantity at $N=10$, again with $\lambda =1$ and $\varsigma =0.1$, $0.2$, $0.3$. Halving the store does not halve the dose a pathogen needs: the certainty boundary moves with $N$, so the requirement is a matter of the ratio of cohort to store rather than of the cohort alone}{figure.8.7}{}}
\newlabel{path:fig:repsupp10@cref}{{[figure][7][8]8.7}{[1][217][]218}}
\newlabel{path:sec:removal}{{8.4}{219}{Extracellular depletion: fixed removal}{section.8.4}{}}
\newlabel{path:sec:removal@cref}{{[section][4][8]8.4}{[1][219][]219}}
\newlabel{path:sec:playingfield}{{8.4.1}{219}{The playing field}{subsection.8.4.1}{}}
\newlabel{path:sec:playingfield@cref}{{[subsection][1][8,4]8.4.1}{[1][219][]219}}
\newlabel{path:eq:logremoval}{{8.8}{219}{The playing field}{equation.8.4.8}{}}
\newlabel{path:eq:logremoval@cref}{{[equation][8][8]8.8}{[1][219][]219}}
\newlabel{path:eq:fixedremoval}{{8.9}{219}{The playing field}{equation.8.4.9}{}}
\newlabel{path:eq:fixedremoval@cref}{{[equation][9][8]8.9}{[1][219][]219}}
\newlabel{path:fig:loads}{{8.8}{220}{Two runs of the same simulation differing in released number alone; every rate is the same. Depletable extracellular immune factors are consumed by the agents they remove, so a large load creates a safety in numbers --- a ``flooding effect'' --- and the killers that act early are not available later. The comparison to draw is between~(c) and~(f). In~(c) the large load has cleared a patch of the killer field and survivors remain inside it, poised on the surrounding host cells; in~(f) the small load is exhausted before the field is, and nothing persists to the next round of absorption and replication. \explorerlink {https://files.volvox.site/users/69b423f5faa2bfbb5b696c61/pages/6a8495d018f4e12409b5ecf1/files/1789007285465-fau0qe-burst_suppression.html}{Simulate}}{figure.8.8}{}}
\newlabel{path:fig:loads@cref}{{[figure][8][8]8.8}{[1][219][]220}}
\newlabel{path:rem:reset}{{8.5}{220}{Killer prevalence resets}{theorem.8.5}{}}
\newlabel{path:rem:reset@cref}{{[remark][5][8]8.5}{[1][220][]220}}
\newlabel{path:rem:onetarget}{{8.6}{220}{One target cell, or many}{theorem.8.6}{}}
\newlabel{path:rem:onetarget@cref}{{[remark][6][8]8.6}{[1][220][]220}}
\newlabel{path:fig:logremoval}{{8.9a}{221}{the saturating law \cref {path:eq:logremoval}}{subfigure.8.9.1}{}}
\newlabel{path:fig:logremoval@cref}{{[subfigure][1][8,9]8.9a}{[1][219][]221}}
\newlabel{sub@path:fig:logremoval}{{a}{221}{the saturating law \cref {path:eq:logremoval}}{subfigure.8.9.1}{}}
\newlabel{sub@path:fig:logremoval@cref}{{[subfigure][1][8,9]8.9a}{[1][219][]221}}
\newlabel{path:fig:capped}{{8.9b}{221}{the step idealisation \cref {path:eq:fixedremoval}}{subfigure.8.9.2}{}}
\newlabel{path:fig:capped@cref}{{[subfigure][2][8,9]8.9b}{[1][219][]221}}
\newlabel{sub@path:fig:capped}{{b}{221}{the step idealisation \cref {path:eq:fixedremoval}}{subfigure.8.9.2}{}}
\newlabel{sub@path:fig:capped@cref}{{[subfigure][2][8,9]8.9b}{[1][219][]221}}
\newlabel{path:fig:removal}{{8.9}{221}{The casualty count against burst size: the law the biology suggests, and the caricature the calculation uses. In~(a) the vertical axis is the number of agents killed, not a rate --- what falls as the burst grows is the killed \emph {fraction}, since the curve flattens while the burst size does not. Doubling a load of ten adds well under one casualty, which is the whole of the flooding intuition in one line. In~(b) that flatness is pushed to its limit: below the budget every agent dies, and above it the plateau at $\vartheta $ means a burst pays the same toll however large it is, so all the marginal yield survives. The reason for adopting the caricature is that it makes the replication ratio computable in closed form}{figure.8.9}{}}
\newlabel{path:fig:removal@cref}{{[figure][9][8]8.9}{[1][219][]221}}
\newlabel{path:sec:ratio}{{8.4.2}{221}{The replication ratio}{subsection.8.4.2}{}}
\newlabel{path:sec:ratio@cref}{{[subsection][2][8,4]8.4.2}{[1][220][]221}}
\newlabel{path:eq:Er}{{8.10}{221}{The replication ratio}{equation.8.4.10}{}}
\newlabel{path:eq:Er@cref}{{[equation][10][8]8.10}{[1][221][]221}}
\newlabel{path:eq:Rdef}{{8.11}{221}{The replication ratio}{equation.8.4.11}{}}
\newlabel{path:eq:Rdef@cref}{{[equation][11][8]8.11}{[1][221][]221}}
\newlabel{path:eq:Rcond}{{8.12}{221}{The replication ratio}{equation.8.4.12}{}}
\newlabel{path:eq:Rcond@cref}{{[equation][12][8]8.12}{[1][221][]221}}
\newlabel{path:eq:burstlaw}{{8.13}{222}{The replication ratio}{equation.8.4.13}{}}
\newlabel{path:eq:burstlaw@cref}{{[equation][13][8]8.13}{[1][221][]222}}
\newlabel{path:eq:Rpartial}{{8.14}{222}{The replication ratio}{equation.8.4.14}{}}
\newlabel{path:eq:Rpartial@cref}{{[equation][14][8]8.14}{[1][221][]222}}
\newlabel{path:eq:FM}{{8.15}{222}{The replication ratio}{equation.8.4.15}{}}
\newlabel{path:eq:FM@cref}{{[equation][15][8]8.15}{[1][222][]222}}
\newlabel{path:sec:normalisation}{{8.4.2.1}{222}{The law is normalised, and the tail is free}{subsubsection.8.4.2.1}{}}
\newlabel{path:sec:normalisation@cref}{{[subsubsection][1][8,4,2]8.4.2.1}{[1][222][]222}}
\newlabel{path:eq:broot}{{8.16}{222}{The law is normalised, and the tail is free}{equation.8.4.16}{}}
\newlabel{path:eq:broot@cref}{{[equation][16][8]8.16}{[1][222][]222}}
\newlabel{path:eq:tail}{{8.17}{222}{The law is normalised, and the tail is free}{equation.8.4.17}{}}
\newlabel{path:eq:tail@cref}{{[equation][17][8]8.17}{[1][222][]222}}
\newlabel{path:sec:proliferability}{{8.4.3}{222}{Analysis of proliferability}{subsection.8.4.3}{}}
\newlabel{path:sec:proliferability@cref}{{[subsection][3][8,4]8.4.3}{[1][222][]222}}
\newlabel{path:prop:ratio}{{8.7}{223}{Fixed-removal replication ratio}{theorem.8.7}{}}
\newlabel{path:prop:ratio@cref}{{[proposition][7][8]8.7}{[1][222][]223}}
\newlabel{path:eq:R0}{{8.18}{223}{Fixed-removal replication ratio}{equation.8.4.18}{}}
\newlabel{path:eq:R0@cref}{{[equation][18][8]8.18}{[1][222][]223}}
\newlabel{path:cor:threshold}{{8.8}{223}{Threshold}{theorem.8.8}{}}
\newlabel{path:cor:threshold@cref}{{[corollary][8][8]8.8}{[1][223][]223}}
\newlabel{path:cor:whythreshold}{{8.9}{223}{Why removal creates a threshold at all}{theorem.8.9}{}}
\newlabel{path:cor:whythreshold@cref}{{[corollary][9][8]8.9}{[1][223][]223}}
\newlabel{path:rem:partial}{{8.10}{223}{Partial-sum expansion}{theorem.8.10}{}}
\newlabel{path:rem:partial@cref}{{[remark][10][8]8.10}{[1][223][]223}}
\newlabel{path:eq:Rintermediate}{{8.19}{223}{Partial-sum expansion}{equation.8.4.19}{}}
\newlabel{path:eq:Rintermediate@cref}{{[equation][19][8]8.19}{[1][223][]223}}
\newlabel{path:rem:integer}{{8.11}{223}{Scope: $\vartheta $ is an integer}{theorem.8.11}{}}
\newlabel{path:rem:integer@cref}{{[remark][11][8]8.11}{[1][223][]223}}
\newlabel{path:fig:ratio}{{8.10}{224}{The replication ratio \cref {path:eq:R0} against $\delta $ at $\lambda =0.8$, $\mu =0$, for budgets $\vartheta =0$, $1$ and $2$, with $V_\infty $ dashed for reference throughout. In panel~(a) the two curves coincide, since at zero removal the surviving release is the whole release; that panel has no crossing, which is \cref {path:cor:whythreshold} drawn. Elsewhere the marked point is the crossing $\mathcal {R}=1$ of \cref {path:cor:threshold}, at $\delta =0.800$ and $\delta =0.494$: to the left of it the infection proliferates, to the right the removal extinguishes it. Each further unit of budget moves the crossing left, so a pathogen facing a heavier toll needs a lower catastrophe rate --- a larger burst --- to stay supercritical}{figure.8.10}{}}
\newlabel{path:fig:ratio@cref}{{[figure][10][8]8.10}{[1][223][]224}}
\newlabel{path:sec:selection}{{8.5}{224}{Bursting and budding: where the advantage lives}{section.8.5}{}}
\newlabel{path:sec:selection@cref}{{[section][5][8]8.5}{[1][224][]224}}
\newlabel{path:sec:settled}{{8.5.1}{224}{What is settled, and what was left open}{subsection.8.5.1}{}}
\newlabel{path:sec:settled@cref}{{[subsection][1][8,5]8.5.1}{[1][224][]224}}
\newlabel{path:eq:L}{{8.20}{224}{What is settled, and what was left open}{equation.8.5.20}{}}
\newlabel{path:eq:L@cref}{{[equation][20][8]8.20}{[1][224][]224}}
\newlabel{path:tab:reversal}{{8.3}{225}{The reversal in six regimes. The last two columns name which strategy wins under linear clearance and which under threshold removal, and they disagree in every row except the tie. The middle row is the boundary $L=1$, where the bursting offspring law is exactly geometric (\cref {p:prop:L1}) and the two ratios agree to the digits shown. The fourth row is strongly supercritical on both sides, so the reversal is not an artefact of one of them being near extinction}{table.8.3}{}}
\newlabel{path:tab:reversal@cref}{{[table][3][8]8.3}{[1][225][]225}}
\newlabel{path:sec:reversal}{{8.5.2}{225}{The reversal}{subsection.8.5.2}{}}
\newlabel{path:sec:reversal@cref}{{[subsection][2][8,5]8.5.2}{[1][224][]225}}
\newlabel{path:prop:reversal}{{8.12}{225}{The flooding criterion reverses under threshold removal}{theorem.8.12}{}}
\newlabel{path:prop:reversal@cref}{{[proposition][12][8]8.12}{[1][225][]225}}
\newlabel{path:eq:reversal}{{8.21}{225}{The flooding criterion reverses under threshold removal}{equation.8.5.21}{}}
\newlabel{path:eq:reversal@cref}{{[equation][21][8]8.21}{[1][225][]225}}
\newlabel{path:fig:reversal}{{8.11}{226}{The sign change, drawn. The gap $\mathcal {R}_{\mathrm {burst}}-\mathcal {R}_{\mathrm {bud}}$ against $\delta $ at $\lambda =1$, $\mu =0.5$, for three removal budgets. All three curves cross zero at the same place, the flooding boundary $\delta _{*}=\lambda \mu /(\lambda +\mu )=1/3$ of \cref {p:eq:floodcrit}, and they cross it in the direction opposite to \cref {p:thm:flood}: bursting wins to the \emph {left}, where $L<1$. Raising the budget deepens the advantage on the left without moving the boundary, because the boundary is a property of the two laws and not of the toll}{figure.8.11}{}}
\newlabel{path:fig:reversal@cref}{{[figure][11][8]8.11}{[1][225][]226}}
\newlabel{path:sec:twoarguments}{{8.5.2}{226}{The reversal}{figure.8.11}{}}
\newlabel{path:sec:twoarguments@cref}{{[subsection][2][8,5]8.5.2}{[1][225][]226}}
\newlabel{path:rem:scope}{{8.13}{226}{Scope of the comparison}{theorem.8.13}{}}
\newlabel{path:rem:scope@cref}{{[remark][13][8]8.13}{[1][226][]226}}
\newlabel{path:sec:spectrum}{{8.5.3}{226}{Placement on the spectrum}{subsection.8.5.3}{}}
\newlabel{path:sec:spectrum@cref}{{[subsection][3][8,5]8.5.3}{[1][226][]226}}
\newlabel{path:fig:tollspectrum}{{8.12}{226}{Per-delivery toll along the budding--bursting spectrum of \cref {p:fig:spectrum}. Moving right, fewer deliveries are made and a smaller total is paid. The ordering is the batch count, not the distributional reversal}{figure.8.12}{}}
\newlabel{path:fig:tollspectrum@cref}{{[figure][12][8]8.12}{[1][226][]226}}
\newlabel{path:sec:ypestis}{{8.5.4}{227}{\texorpdfstring {\yp {}}{Y. pestis}}{subsection.8.5.4}{}}
\newlabel{path:sec:ypestis@cref}{{[subsection][4][8,5]8.5.4}{[1][227][]227}}
\newlabel{path:tab:ypestis}{{8.4}{227}{A reading of already-named objects under the \yp {} mapping. The formulae themselves do not depend on it}{table.8.4}{}}
\newlabel{path:tab:ypestis@cref}{{[table][4][8]8.4}{[1][227][]227}}
\newlabel{path:sec:discussion}{{8.6}{228}{Discussion}{section.8.6}{}}
\newlabel{path:sec:discussion@cref}{{[section][6][8]8.6}{[1][227][]228}}
\newlabel{path:sec:limitations}{{8.6.1}{228}{Limitations}{subsection.8.6.1}{}}
\newlabel{path:sec:limitations@cref}{{[subsection][1][8,6]8.6.1}{[1][227][]228}}
\newlabel{path:sec:open}{{8.6.2}{228}{Open problems}{subsection.8.6.2}{}}
\newlabel{path:sec:open@cref}{{[subsection][2][8,6]8.6.2}{[1][228][]228}}
\newlabel{path:sec:forward}{{8.6.3}{229}{Connections forward}{subsection.8.6.3}{}}
\newlabel{path:sec:forward@cref}{{[subsection][3][8,6]8.6.3}{[1][229][]229}}
\newlabel{path:app:verification}{{8.A}{229}{Numerical verification}{section.8.appendix.A}{}}
\newlabel{path:app:verification@cref}{{[section][1][8]8.A}{[1][229][]229}}
\newlabel{path:tab:verification}{{8.5}{230}{The verification record. The burst-law entry is the only one whose agreement is limited by anything other than floating-point arithmetic: its error is set by the truncation tail, not by the formula. Every other row agrees to within a few units in the last place, which is what one expects of two algebraically equivalent evaluations rather than of a formula that happens to fit}{table.8.5}{}}
\newlabel{path:tab:verification@cref}{{[table][5][8]8.5}{[1][229][]230}}
\newlabel{path:app:computational}{{8.B}{230}{Computational notes}{section.8.appendix.B}{}}
\newlabel{path:app:computational@cref}{{[section][2][8]8.B}{[1][230][]230}}
\newlabel{path:app:solver}{{8.B.1}{230}{The first-step solver}{subsection.8.appendix.B.1}{}}
\newlabel{path:app:solver@cref}{{[subsection][1][8,2]8.B.1}{[1][230][]230}}
\newlabel{path:app:quadtree}{{8.B.2}{231}{Quadtree sampling}{subsection.8.appendix.B.2}{}}
\newlabel{path:app:quadtree@cref}{{[subsection][2][8,2]8.B.2}{[1][231][]231}}
\newlabel{path:fig:quadtree}{{8.13}{231}{The quadtree index during a simulation run, with the cells local to the macrophage shaded pink; only the particles inside those cells are tested for collision. (a) A sparse field, where the index buys little. (b) A dense one, where it buys a great deal: the shaded region is a small fraction of the frame, so the number of collision tests grows far more slowly than the number of particles. This is why the large-load runs of \cref {path:fig:loads} are feasible at all}{figure.8.13}{}}
\newlabel{path:fig:quadtree@cref}{{[figure][13][8]8.13}{[1][231][]231}}
\newlabel{var:ch:rates}{{9}{232}{Models with non-constant rates}{chapter.9}{}}
\newlabel{var:ch:rates@cref}{{[chapter][9][]9}{[1][232][]232}}
\newlabel{var:sec:introduction}{{9.1}{232}{Introduction}{section.9.1}{}}
\newlabel{var:sec:introduction@cref}{{[section][1][9]9.1}{[1][232][]232}}
\newlabel{var:tab:models}{{9.1}{233}{The models, and the mechanism by which a rate moves in each. The last column is the shape of the chapter: \cref {var:sec:speciation} is its mathematical end, and the rows below it halt at a closed-form fragment or at a simulation. The two below the rule are carried as bare specifications in \cref {var:app:plague,var:app:publicgood}, the birth--plague model in two variants, so the count of distinct specifications is eight}{table.9.1}{}}
\newlabel{var:tab:models@cref}{{[table][1][9]9.1}{[1][233][]233}}
\newlabel{var:sec:notation}{{9.1.1}{233}{Notation}{subsection.9.1.1}{}}
\newlabel{var:sec:notation@cref}{{[subsection][1][9,1]9.1.1}{[1][233][]233}}
\newlabel{var:tab:notation}{{9.2}{234}{Notation, organised by model. The rates $\lambda $, $\mu $ and $\delta $, the roots $a$ and $b$ and the burst size $\Kb $ carry their thesis-wide meanings everywhere in this chapter and are never redefined. $V$ and $c$ are shared between \cref {var:sec:rrrr,var:sec:ca} with the \emph {same} meanings, a potential and the rate at which it is spent, which is the point of the pairing: what differs between the two models is that \cref {var:sec:ca} charges the inflow per cell and the outflow per type. Where one letter genuinely serves two models --- $\alpha $ as an exponent in \cref {var:sec:pimentel,var:sec:pimentelplus}, $N$ as an abundance in \cref {var:sec:speciation} and an integer capacity in \cref {var:sec:pimentel}, $X$ as a count of individuals in \cref {var:sec:rrrr} and of types in \cref {var:sec:ca} --- the two uses are separated by this table and never occur in the same section}{table.9.2}{}}
\newlabel{var:tab:notation@cref}{{[table][2][9]9.2}{[1][233][]234}}
\newlabel{var:sec:speciation}{{9.2}{235}{Logistic speciation}{section.9.2}{}}
\newlabel{var:sec:speciation@cref}{{[section][2][9]9.2}{[1][234][]235}}
\newlabel{var:sec:variationmotivation}{{9.2.1}{235}{Variation, and what suppresses it}{subsection.9.2.1}{}}
\newlabel{var:sec:variationmotivation@cref}{{[subsection][1][9,2]9.2.1}{[1][234][]235}}
\newlabel{var:sec:potp}{{9.2.2}{235}{The push of the past}{subsection.9.2.2}{}}
\newlabel{var:sec:potp@cref}{{[subsection][2][9,2]9.2.2}{[1][235][]235}}
\newlabel{var:sec:speciationmodel}{{9.2.3}{235}{The model}{subsection.9.2.3}{}}
\newlabel{var:sec:speciationmodel@cref}{{[subsection][3][9,2]9.2.3}{[1][235][]235}}
\newlabel{var:eq:logistic}{{9.1}{235}{The model}{equation.9.2.1}{}}
\newlabel{var:eq:logistic@cref}{{[equation][1][9]9.1}{[1][235][]235}}
\newlabel{var:eq:logisticsol}{{9.2}{236}{The model}{equation.9.2.2}{}}
\newlabel{var:eq:logisticsol@cref}{{[equation][2][9]9.2}{[1][235][]236}}
\newlabel{var:eq:transitions}{{9.3}{236}{The model}{equation.9.2.3}{}}
\newlabel{var:eq:transitions@cref}{{[equation][3][9]9.3}{[1][236][]236}}
\newlabel{var:eq:rate1}{{9.4}{236}{The model}{equation.9.2.4}{}}
\newlabel{var:eq:rate1@cref}{{[equation][4][9]9.4}{[1][236][]236}}
\newlabel{var:eq:ratesolved}{{9.5}{236}{The model}{equation.9.2.5}{}}
\newlabel{var:eq:ratesolved@cref}{{[equation][5][9]9.5}{[1][236][]236}}
\newlabel{var:rem:scope}{{9.1}{236}{What the rate law is, and is not, claiming}{theorem.9.1}{}}
\newlabel{var:rem:scope@cref}{{[remark][1][9]9.1}{[1][236][]236}}
\newlabel{var:sec:marginal}{{9.2.4}{236}{The marginal law}{subsection.9.2.4}{}}
\newlabel{var:sec:marginal@cref}{{[subsection][4][9,2]9.2.4}{[1][236][]236}}
\newlabel{var:eq:Bdef}{{9.6}{236}{The marginal law}{equation.9.2.6}{}}
\newlabel{var:eq:Bdef@cref}{{[equation][6][9]9.6}{[1][236][]236}}
\newlabel{var:prop:yule-geom}{{9.2}{236}{Marginal law of an inhomogeneous Yule process}{theorem.9.2}{}}
\newlabel{var:prop:yule-geom@cref}{{[proposition][2][9]9.2}{[1][236][]236}}
\newlabel{var:eq:geomlaw}{{9.7}{236}{Marginal law of an inhomogeneous Yule process}{equation.9.2.7}{}}
\newlabel{var:eq:geomlaw@cref}{{[equation][7][9]9.7}{[1][236][]236}}
\newlabel{var:sec:mean}{{9.2.5}{237}{The mean, directly}{subsection.9.2.5}{}}
\newlabel{var:sec:mean@cref}{{[subsection][5][9,2]9.2.5}{[1][237][]237}}
\newlabel{var:eq:meanode}{{9.8}{237}{The mean, directly}{equation.9.2.8}{}}
\newlabel{var:eq:meanode@cref}{{[equation][8][9]9.8}{[1][237][]237}}
\newlabel{var:eq:intbeta}{{9.9}{237}{The mean, directly}{equation.9.2.9}{}}
\newlabel{var:eq:intbeta@cref}{{[equation][9][9]9.9}{[1][237][]237}}
\newlabel{var:eq:Bt}{{9.10}{237}{The mean, directly}{equation.9.2.10}{}}
\newlabel{var:eq:Bt@cref}{{[equation][10][9]9.10}{[1][237][]237}}
\newlabel{var:sec:limit}{{9.2.6}{237}{The limit law}{subsection.9.2.6}{}}
\newlabel{var:sec:limit@cref}{{[subsection][6][9,2]9.2.6}{[1][237][]237}}
\newlabel{var:cor:logistic-limit}{{9.3}{237}{Limit law under the logistic rate}{theorem.9.3}{}}
\newlabel{var:cor:logistic-limit@cref}{{[corollary][3][9]9.3}{[1][237][]237}}
\newlabel{var:eq:Binf}{{9.11}{237}{Limit law under the logistic rate}{equation.9.2.11}{}}
\newlabel{var:eq:Binf@cref}{{[equation][11][9]9.11}{[1][237][]237}}
\newlabel{var:eq:limitmean}{{9.12}{237}{Limit law under the logistic rate}{equation.9.2.12}{}}
\newlabel{var:eq:limitmean@cref}{{[equation][12][9]9.12}{[1][237][]237}}
\newlabel{var:fig:logspec-limit}{{9.1}{238}{What \cref {var:cor:logistic-limit} says, drawn from the closed form. Panel~(a): the law of $S_\infty $ is geometric, so most of its mass sits at small counts however large the mean, and its coefficient of variation is close to one --- fitting this model to a single clade is fitting a mean to one draw from that. Panel~(b): the expected final count $\ex {S_\infty }=(K/N_0)^{\sigma K/\varrho }$ depends on the five constants $\sigma ,K,\varrho ,N_0,\xi $ only through the two dimensionless groups on the axes, so every point on a contour is a genus with the same expected diversity reached from different parameters}{figure.9.1}{}}
\newlabel{var:fig:logspec-limit@cref}{{[figure][1][9]9.1}{[1][238][]238}}
\newlabel{var:sec:altrates}{{9.2.7}{238}{The other rate laws}{subsection.9.2.7}{}}
\newlabel{var:sec:altrates@cref}{{[subsection][7][9,2]9.2.7}{[1][238][]238}}
\newlabel{var:eq:altrates}{{9.13}{238}{The other rate laws}{equation.9.2.13}{}}
\newlabel{var:eq:altrates@cref}{{[equation][13][9]9.13}{[1][238][]238}}
\newlabel{var:eq:beta3}{{9.14}{238}{The other rate laws}{equation.9.2.14}{}}
\newlabel{var:eq:beta3@cref}{{[equation][14][9]9.14}{[1][238][]238}}
\newlabel{var:fig:lsmean}{{9.2}{239}{Mean species count \cref {var:eq:Bt} for $N_0=1,2,5$, with $K=10$, $\varrho =1$ and $\sigma =0.08$; dashed lines are the limits $(K/N_0)^{\sigma K/\varrho }$. A smaller founding abundance leaves more integrated unused capacity, so the limiting mean is higher. This is the pure-birth model of \cref {var:eq:transitions}, not the birth--death numerical variant recorded at \cref {m:fig:lsmean}; \cref {var:sec:withextinction} is what a constant extinction rate does to this one. \explorerlink {https://markov.volvox.site/-/adam74aldridge/logistic-environment-type-count-process}{Simulate}}{figure.9.2}{}}
\newlabel{var:fig:lsmean@cref}{{[figure][2][9]9.2}{[1][238][]239}}
\newlabel{var:eq:mixture}{{9.15}{239}{The other rate laws}{equation.9.2.15}{}}
\newlabel{var:eq:mixture@cref}{{[equation][15][9]9.15}{[1][239][]239}}
\newlabel{var:sec:withextinction}{{9.2.8}{239}{What extinction does}{subsection.9.2.8}{}}
\newlabel{var:sec:withextinction@cref}{{[subsection][8][9,2]9.2.8}{[1][239][]239}}
\newlabel{var:eq:meanwithdeath}{{9.16}{240}{What extinction does}{equation.9.2.16}{}}
\newlabel{var:eq:meanwithdeath@cref}{{[equation][16][9]9.16}{[1][239][]240}}
\newlabel{var:rem:open-speciation}{{9.2.8}{240}{Open questions}{section*.66}{}}
\newlabel{var:rem:open-speciation@cref}{{[subsection][8][9,2]9.2.8}{[1][240][]240}}
\newlabel{var:fig:extinction}{{9.3}{241}{Logistic speciation with a constant extinction rate, at $(K,N_0,\varrho ,\sigma ,\mu )=(10,5,0.8,0.8,0.4)$ from $S_0=1$: $160$ exact realisations in pale blue, their mean over $400$ paths in bold. Origination follows $\beta _3(t)=\sigma N'(t)$, drawn below against the constant $\mu $. $N_0=K/2$ starts the abundance at the inflection of \cref {var:eq:logistic}, so origination is at its largest at $t=0$ and decays from there. The mean peaks at $10.0$ species at $t=3.24$, and the rates cross at $t=3.29$: the two coincide, as \cref {var:eq:meanwithdeath} requires. Extinction of a genus's last species is suppressed, which is why the paths have a floor at one and the mean settles at $1.02$ rather than decaying to zero. Set against the same origination law without extinction, where \cref {var:eq:beta3} gives a limiting mean of $\ee ^{\sigma (K-N_0)}\approx 55$ approached from below and never left, the cost of a constant $\mu $ is the whole limit law}{figure.9.3}{}}
\newlabel{var:fig:extinction@cref}{{[figure][3][9]9.3}{[1][240][]241}}
\newlabel{var:sec:pimentel}{{9.3}{241}{Competitive reversal: the Pimentel model}{section.9.3}{}}
\newlabel{var:sec:pimentel@cref}{{[section][3][9]9.3}{[1][241][]241}}
\newlabel{var:sec:pimentelexperiment}{{9.3.1}{241}{The experiment}{subsection.9.3.1}{}}
\newlabel{var:sec:pimentelexperiment@cref}{{[subsection][1][9,3]9.3.1}{[1][241][]241}}
\newlabel{var:sec:pimentelmodel}{{9.3.2}{242}{The model}{subsection.9.3.2}{}}
\newlabel{var:sec:pimentelmodel@cref}{{[subsection][2][9,3]9.3.2}{[1][242][]242}}
\newlabel{var:eq:kernel}{{9.17}{242}{The model}{equation.9.3.17}{}}
\newlabel{var:eq:kernel@cref}{{[equation][17][9]9.17}{[1][242][]242}}
\newlabel{var:eq:blockprobs}{{9.18}{242}{The model}{equation.9.3.18}{}}
\newlabel{var:eq:blockprobs@cref}{{[equation][18][9]9.18}{[1][242][]242}}
\newlabel{var:eq:resistupdate}{{9.19}{243}{The model}{equation.9.3.19}{}}
\newlabel{var:eq:resistupdate@cref}{{[equation][19][9]9.19}{[1][243][]243}}
\newlabel{var:eq:dR}{{9.20}{243}{The model}{equation.9.3.20}{}}
\newlabel{var:eq:dR@cref}{{[equation][20][9]9.20}{[1][243][]243}}
\newlabel{var:sec:pimentelresults}{{9.3.3}{243}{Results}{subsection.9.3.3}{}}
\newlabel{var:sec:pimentelresults@cref}{{[subsection][3][9,3]9.3.3}{[1][243][]243}}
\newlabel{var:fig:pim-two-runs}{{9.4}{244}{Two realisations of the Pimentel chain at $\alpha =1$, $N=100$, $\varepsilon =10$, the physical case. Each reverses: an early advantage is not kept, and the two histories, started from the same symmetric state, absorb at opposite boundaries}{figure.9.4}{}}
\newlabel{var:fig:pim-two-runs@cref}{{[figure][4][9]9.4}{[1][243][]244}}
\newlabel{var:fig:pim-runs}{{9.5}{244}{Two independent realisations from each of two initial conditions, at $\alpha =10$, $N=10^{5}$, $\varepsilon =1$, over $2\times 10^{6}$ steps. Panel~(a): started evenly matched, and so near the critical state, the runs separate early and stay apart. Panel~(b): started with one in the majority, and so away from it, they track one another. Proximity to the critical state, not the size of the noise, decides how far apart two histories of the same model end up}{figure.9.5}{}}
\newlabel{var:fig:pim-runs@cref}{{[figure][5][9]9.5}{[1][243][]244}}
\newlabel{var:sec:critical}{{9.3.4}{244}{The moving critical point}{subsection.9.3.4}{}}
\newlabel{var:sec:critical@cref}{{[subsection][4][9,3]9.3.4}{[1][243][]244}}
\newlabel{var:eq:icrit}{{9.21}{244}{The moving critical point}{equation.9.3.21}{}}
\newlabel{var:eq:icrit@cref}{{[equation][21][9]9.21}{[1][244][]244}}
\newlabel{var:fig:pim-critical}{{9.6}{245}{A single realisation at $\alpha =3$, $N=10^{4}$, $\varepsilon =1$, with the critical state $i_{c,t}$ of \cref {var:eq:icrit} overlaid. The critical state does not stay put. It is dragged towards whichever boundary the leading population is heading for, and saturates near $0$ or near $N$ while that population is well ahead; each reversal of $A_t$ follows $i_{c,t}$ crossing it}{figure.9.6}{}}
\newlabel{var:fig:pim-critical@cref}{{[figure][6][9]9.6}{[1][245][]245}}
\newlabel{var:sec:automaton}{{9.3.5}{245}{A spatial version}{subsection.9.3.5}{}}
\newlabel{var:sec:automaton@cref}{{[subsection][5][9,3]9.3.5}{[1][245][]245}}
\newlabel{var:fig:automaton}{{9.7}{246}{Six frames from the spatial version, left to right along the top row then the bottom: advance, squeeze, fringe, crossing, recovery, fixation. Unoccupied cells are green in the first frame. Red advances until blue is confined to one corner, where a pale fringe of elevated resistance appears; blue then crosses back and takes the domain. The reversal is spatial as well as numerical. \explorerlink {https://files.volvox.site/users/69b423f5faa2bfbb5b696c61/pages/6a84961018f4e12409b5ed0a/files/1789007399583-rl4747-pimentel.html}{Simulate}}{figure.9.7}{}}
\newlabel{var:fig:automaton@cref}{{[figure][7][9]9.7}{[1][245][]246}}
\newlabel{var:sec:pimentelplus}{{9.3.6}{246}{Coupling speciation to competition}{subsection.9.3.6}{}}
\newlabel{var:sec:pimentelplus@cref}{{[subsection][6][9,3]9.3.6}{[1][245][]246}}
\newlabel{var:eq:fitkernel}{{9.22}{246}{Coupling speciation to competition}{equation.9.3.22}{}}
\newlabel{var:eq:fitkernel@cref}{{[equation][22][9]9.22}{[1][246][]246}}
\newlabel{var:eq:fitupdate}{{9.23}{246}{Coupling speciation to competition}{equation.9.3.23}{}}
\newlabel{var:eq:fitupdate@cref}{{[equation][23][9]9.23}{[1][246][]246}}
\newlabel{var:eq:subspec}{{9.24}{247}{Coupling speciation to competition}{equation.9.3.24}{}}
\newlabel{var:eq:subspec@cref}{{[equation][24][9]9.24}{[1][247][]247}}
\newlabel{var:eq:fitpenalty}{{9.25}{247}{Coupling speciation to competition}{equation.9.3.25}{}}
\newlabel{var:eq:fitpenalty@cref}{{[equation][25][9]9.25}{[1][247][]247}}
\newlabel{var:rem:open-plus}{{9.3.6}{247}{Open questions}{section*.67}{}}
\newlabel{var:rem:open-plus@cref}{{[subsection][6][9,3]9.3.6}{[1][247][]247}}
\newlabel{var:rem:open-pimentel}{{9.3.6}{247}{Open questions}{section*.68}{}}
\newlabel{var:rem:open-pimentel@cref}{{[subsection][6][9,3]9.3.6}{[1][247][]247}}
\newlabel{var:sec:rrrr}{{9.4}{248}{Rise, run, ruin, rejuvenation}{section.9.4}{}}
\newlabel{var:sec:rrrr@cref}{{[section][4][9]9.4}{[1][248][]248}}
\newlabel{var:sec:rrrrmotivation}{{9.4.1}{248}{The cycle the model is for}{subsection.9.4.1}{}}
\newlabel{var:sec:rrrrmotivation@cref}{{[subsection][1][9,4]9.4.1}{[1][248][]248}}
\newlabel{var:sec:rrrrmodel}{{9.4.2}{248}{The model}{subsection.9.4.2}{}}
\newlabel{var:sec:rrrrmodel@cref}{{[subsection][2][9,4]9.4.2}{[1][248][]248}}
\newlabel{var:eq:rrrrV}{{9.26}{248}{The model}{equation.9.4.26}{}}
\newlabel{var:eq:rrrrV@cref}{{[equation][26][9]9.26}{[1][248][]248}}
\newlabel{var:eq:rrrr}{{9.27}{248}{The model}{equation.9.4.27}{}}
\newlabel{var:eq:rrrr@cref}{{[equation][27][9]9.27}{[1][248][]248}}
\newlabel{var:sec:rrrrsims}{{9.4.3}{249}{What the simulations show}{subsection.9.4.3}{}}
\newlabel{var:sec:rrrrsims@cref}{{[subsection][3][9,4]9.4.3}{[1][249][]249}}
\newlabel{var:fig:vbc-two}{{9.8}{249}{Two realisations of \cref {var:eq:rrrrV,var:eq:rrrr}: potential $V(t)$ in vermillion, agency count $X_t$ in blue ($\phi =4$, $c=0.8$, $\lambda =0.5$, $\delta =0.25$, $\omega =0.15$). Every cycle has the same four parts in the same order --- potential builds while $X_t$ is small, $X_t$ grows into the surplus, the surplus is consumed, the collapse follows --- and no two are copies of one another}{figure.9.8}{}}
\newlabel{var:fig:vbc-two@cref}{{[figure][8][9]9.8}{[1][249][]249}}
\newlabel{var:fig:vbc-many}{{9.9}{249}{Several realisations from identical parameters and identical initial conditions. The first cycle is nearly common to all of them; later cycles are not, each catastrophe having reset the phase by a random amount. The parameters fix the amplitude, not the timing}{figure.9.9}{}}
\newlabel{var:fig:vbc-many@cref}{{[figure][9][9]9.9}{[1][249][]249}}
\newlabel{var:rem:open-rrrr}{{9.4.3}{250}{Open questions}{section*.69}{}}
\newlabel{var:rem:open-rrrr@cref}{{[subsection][3][9,4]9.4.3}{[1][250][]250}}
\newlabel{var:sec:ca}{{9.5}{250}{Speciation as catastrophe}{section.9.5}{}}
\newlabel{var:sec:ca@cref}{{[section][5][9]9.5}{[1][250][]250}}
\newlabel{var:sec:camotivation}{{9.5.1}{250}{The same hazard, with the inflow made endogenous}{subsection.9.5.1}{}}
\newlabel{var:sec:camotivation@cref}{{[subsection][1][9,5]9.5.1}{[1][250][]250}}
\newlabel{var:sec:castate}{{9.5.2}{251}{State}{subsection.9.5.2}{}}
\newlabel{var:sec:castate@cref}{{[subsection][2][9,5]9.5.2}{[1][251][]251}}
\newlabel{var:sec:cacompetition}{{9.5.3}{251}{Competition}{subsection.9.5.3}{}}
\newlabel{var:sec:cacompetition@cref}{{[subsection][3][9,5]9.5.3}{[1][251][]251}}
\newlabel{var:eq:castrength}{{9.28}{251}{Competition}{equation.9.5.28}{}}
\newlabel{var:eq:castrength@cref}{{[equation][28][9]9.28}{[1][251][]251}}
\newlabel{var:eq:catypeIII}{{9.29}{252}{Competition}{equation.9.5.29}{}}
\newlabel{var:eq:catypeIII@cref}{{[equation][29][9]9.29}{[1][251][]252}}
\newlabel{var:sec:camutation}{{9.5.4}{252}{Mutation and the potential}{subsection.9.5.4}{}}
\newlabel{var:sec:camutation@cref}{{[subsection][4][9,5]9.5.4}{[1][252][]252}}
\newlabel{var:eq:camut}{{9.30}{252}{Mutation and the potential}{equation.9.5.30}{}}
\newlabel{var:eq:camut@cref}{{[equation][30][9]9.30}{[1][252][]252}}
\newlabel{var:eq:caode}{{9.31}{252}{Mutation and the potential}{equation.9.5.31}{}}
\newlabel{var:eq:caode@cref}{{[equation][31][9]9.31}{[1][252][]252}}
\newlabel{var:fig:ca-support}{{9.10}{253}{The same seed at the same time ($t=2000$), with neighbour support on and off. With it, domains are compact and their boundaries short: seven per cent of occupied cells lie on a species boundary. Without it, boundaries wander, domains interpenetrate, and the boundary fraction rises to thirty-six per cent. The ragged state also carries far more species at once --- $1335$ against $26$ --- because nothing gathers a domain together and small holdings are never squeezed out. Hue labels the taxonomy and nothing else: it is carried along by species the competition rules have already made, and no cell ever changes hands because of a colour. \explorerlink {https://files.volvox.site/users/69b423f5faa2bfbb5b696c61/pages/6a84961018f4e12409b5ed0a/files/1789008065523-qk5hsl-multiplicative_dissipation.html}{Simulate}}{figure.9.10}{}}
\newlabel{var:fig:ca-support@cref}{{[figure][10][9]9.10}{[1][252][]253}}
\newlabel{var:sec:cafragment}{{9.5.5}{253}{Fragmentation}{subsection.9.5.5}{}}
\newlabel{var:sec:cafragment@cref}{{[subsection][5][9,5]9.5.5}{[1][253][]253}}
\newlabel{var:eq:cafrag}{{9.32}{254}{Fragmentation}{equation.9.5.32}{}}
\newlabel{var:eq:cafrag@cref}{{[equation][32][9]9.32}{[1][254][]254}}
\newlabel{var:fig:ca-burst}{{9.11}{254}{A burst, across the single frame in which it happens, at $t\approx 1239$. The species held $12{,}826$ cells --- a fifth of the lattice --- and still had $V_i=755$ in hand when \cref {var:eq:cafrag} fired; all $1444$ of its sub-species became species at once, and it was the only species to fragment in that frame. Nothing has moved: the outline of the parent is exactly preserved and its neighbours are untouched. What has changed is the law inside that outline. Each fragment is drawn in a fresh hue here so that the division is unmistakable. \explorerlink {https://files.volvox.site/users/69b423f5faa2bfbb5b696c61/pages/6a84961018f4e12409b5ed0a/files/1789008065523-qk5hsl-multiplicative_dissipation.html}{Simulate}}{figure.9.11}{}}
\newlabel{var:fig:ca-burst@cref}{{[figure][11][9]9.11}{[1][254][]254}}
\newlabel{var:sec:caconsequences}{{9.5.6}{254}{What the simulations show}{subsection.9.5.6}{}}
\newlabel{var:sec:caconsequences@cref}{{[subsection][6][9,5]9.5.6}{[1][254][]254}}
\newlabel{var:eq:cabalance}{{9.33}{255}{The population sits at the balance point}{equation.9.5.33}{}}
\newlabel{var:eq:cabalance@cref}{{[equation][33][9]9.33}{[1][255][]255}}
\newlabel{var:fig:ca-sequence}{{9.12}{255}{One run on a $256^2$ lattice. At $t=250$ a single species holds six per cent of the ground and already carries $443$ sub-species, visible as mottling within one narrow band of hue. By $t=600$ it has burst six times and the fragments are still spreading into open ground. The lattice closes at about $t=1100$; by $t=3000$ the run holds some fifty living species out of $125{,}000$ ever created. \explorerlink {https://files.volvox.site/users/69b423f5faa2bfbb5b696c61/pages/6a84961018f4e12409b5ed0a/files/1789008065523-qk5hsl-multiplicative_dissipation.html}{Simulate}}{figure.9.12}{}}
\newlabel{var:fig:ca-sequence@cref}{{[figure][12][9]9.12}{[1][255][]255}}
\newlabel{var:rem:open-ca}{{9.5.6}{255}{Open questions}{section*.74}{}}
\newlabel{var:rem:open-ca@cref}{{[subsection][6][9,5]9.5.6}{[1][255][]255}}
\newlabel{var:fig:ca-founder}{{9.13}{256}{The founder species, followed from its first cell to the burst that ends it at $t=354$. Potential $V_i$ above, sub-species count $X_i$ below against the balance point $\varphi N_i/c$ of \cref {var:eq:cabalance}; the dotted line marks the burst. The lower panel is the argument of \cref {var:sec:caconsequences} in one picture --- the dashed curve is area, rescaled, and diversity tracks it. The species is not bankrupt when it bursts: it still holds $V_i\approx 730$ on $8689$ cells, but by then it carries $964$ sub-species and the hazard \cref {var:eq:cafrag} has grown large enough to fire}{figure.9.13}{}}
\newlabel{var:fig:ca-founder@cref}{{[figure][13][9]9.13}{[1][255][]256}}
\newlabel{var:app:plague}{{9.A}{257}{The birth--plague model}{section.9.appendix.A}{}}
\newlabel{var:app:plague@cref}{{[section][1][9]9.A}{[1][257][]257}}
\newlabel{var:eq:plague}{{9.34}{257}{The birth--plague model}{equation.9.appendix.A.34}{}}
\newlabel{var:eq:plague@cref}{{[equation][34][9]9.34}{[1][257][]257}}
\newlabel{var:eq:vitalityV}{{9.35}{257}{The vitality-coupled variant}{equation.9.appendix.A.35}{}}
\newlabel{var:eq:vitalityV@cref}{{[equation][35][9]9.35}{[1][257][]257}}
\newlabel{var:eq:vitality}{{9.36}{257}{The vitality-coupled variant}{equation.9.appendix.A.36}{}}
\newlabel{var:eq:vitality@cref}{{[equation][36][9]9.36}{[1][257][]257}}
\newlabel{var:fig:bpplots}{{9.14}{258}{Two realisations of \cref {var:eq:plague} at $(\lambda ,\delta ,\chi ,\mu )=(1,\,0.1,\,0.02,\,10)$: uninfected count $X_t$ in blue, infected count $Y_t$ in vermillion. The clearance rate is large relative to the contact rate, so infected individuals never accumulate and the epidemic acts as a brake on $X_t$ rather than as a collapse of it. Set beside \cref {var:fig:vbc-many}, the fluctuations never reset a phase. \explorerlink {https://markov.volvox.site/-/adam74aldridge/birth-plague-model}{Simulate}}{figure.9.14}{}}
\newlabel{var:fig:bpplots@cref}{{[figure][14][9]9.14}{[1][257][]258}}
\newlabel{var:app:publicgood}{{9.B}{258}{Public good accumulation}{section.9.appendix.B}{}}
\newlabel{var:app:publicgood@cref}{{[section][2][9]9.B}{[1][258][]258}}
\newlabel{var:eq:publicgood}{{9.37}{258}{Public good accumulation}{equation.9.appendix.B.37}{}}
\newlabel{var:eq:publicgood@cref}{{[equation][37][9]9.37}{[1][258][]258}}
\newlabel{var:eq:kstar}{{9.38}{258}{The minimum founding cohort}{equation.9.appendix.B.38}{}}
\newlabel{var:eq:kstar@cref}{{[equation][38][9]9.38}{[1][258][]258}}
\newlabel{var:fig:kstar}{{9.15}{259}{$K^\ast $ computed at $(\lambda _0,\alpha ,g,c)=(0.6,0.25,0.15,0.45)$, so that $P$ falls whenever $X_t<c/g=3$, from $600$ paths per point. The survival probability as a function of $K$ is a steep sigmoid through one half, and $K^\ast $ itself rises with the excess death rate $\mu -\lambda _0$. The dashed line in panel~(b) is $c/g$: $K^\ast $ sits above it, as it must, because a founding cohort smaller than the clearance threshold cannot start the loop. At $\mu -\lambda _0=1.2$ no $K\le 28$ crossed one half, which is a statement about the parameter point rather than a failure of the definition. This is one parameter slice and a numerical estimate, not a closed form: unlike \cref {path:cor:threshold}, which is a statement about a budget, $K^\ast $ is defined by a limit of a survival probability and there is no reason to expect it in closed form}{figure.9.15}{}}
\newlabel{var:fig:kstar@cref}{{[figure][15][9]9.15}{[1][259][]259}}
\newlabel{var:app:pgf}{{9.C}{259}{The inhomogeneous Yule generating function}{section.9.appendix.C}{}}
\newlabel{var:app:pgf@cref}{{[section][3][9]9.C}{[1][259][]259}}
\newlabel{var:sec:pgf}{{9.C.1}{259}{The generating function}{subsection.9.appendix.C.1}{}}
\newlabel{var:sec:pgf@cref}{{[subsection][1][9,3]9.C.1}{[1][259][]259}}
\newlabel{var:eq:fke}{{9.39}{259}{The generating function}{equation.9.appendix.C.39}{}}
\newlabel{var:eq:fke@cref}{{[equation][39][9]9.39}{[1][259][]259}}
\newlabel{var:eq:pgfpde}{{9.40}{259}{The generating function}{equation.9.appendix.C.40}{}}
\newlabel{var:eq:pgfpde@cref}{{[equation][40][9]9.40}{[1][259][]259}}
\newlabel{var:eq:charics}{{9.41}{259}{The generating function}{equation.9.appendix.C.41}{}}
\newlabel{var:eq:charics@cref}{{[equation][41][9]9.41}{[1][259][]259}}
\newlabel{var:eq:chareqs}{{9.42}{260}{The generating function}{equation.9.appendix.C.42}{}}
\newlabel{var:eq:chareqs@cref}{{[equation][42][9]9.42}{[1][259][]260}}
\newlabel{var:eq:charsep}{{9.43}{260}{The generating function}{equation.9.appendix.C.43}{}}
\newlabel{var:eq:charsep@cref}{{[equation][43][9]9.43}{[1][260][]260}}
\newlabel{var:eq:pgf}{{9.44}{260}{The generating function}{equation.9.appendix.C.44}{}}
\newlabel{var:eq:pgf@cref}{{[equation][44][9]9.44}{[1][260][]260}}
\newlabel{var:sec:distribution}{{9.C.2}{260}{The distribution}{subsection.9.appendix.C.2}{}}
\newlabel{var:sec:distribution@cref}{{[subsection][2][9,3]9.C.2}{[1][260][]260}}
\newlabel{var:eq:geompgf}{{9.45}{260}{The distribution}{equation.9.appendix.C.45}{}}
\newlabel{var:eq:geompgf@cref}{{[equation][45][9]9.45}{[1][260][]260}}
\newlabel{ch:conclusion}{{10}{261}{Conclusion}{chapter.10}{}}
\newlabel{ch:conclusion@cref}{{[chapter][10][]10}{[1][261][]261}}
\newlabel{con:sec:contributions}{{10.1}{261}{Contributions, chapter by chapter}{section.10.1}{}}
\newlabel{con:sec:contributions@cref}{{[section][1][10]10.1}{[1][261][]261}}
\newlabel{con:tab:contributions}{{10.1}{263}{What each chapter settles, and what it hands on. The right-hand column is drawn from the chapters' own statements of scope, not added here}{table.10.1}{}}
\newlabel{con:tab:contributions@cref}{{[table][1][10]10.1}{[1][263][]263}}
\newlabel{con:sec:questions}{{10.2}{264}{The four questions, answered}{section.10.2}{}}
\newlabel{con:sec:questions@cref}{{[section][2][10]10.2}{[1][263][]264}}
\newlabel{con:sec:q1}{{10.2.1}{264}{Why \texorpdfstring {\yp {}}{Y. pestis} switches phase}{subsection.10.2.1}{}}
\newlabel{con:sec:q1@cref}{{[subsection][1][10,2]10.2.1}{[1][263][]264}}
\newlabel{con:sec:q2}{{10.2.2}{264}{Whether bursting pays}{subsection.10.2.2}{}}
\newlabel{con:sec:q2@cref}{{[subsection][2][10,2]10.2.2}{[1][264][]264}}
\newlabel{con:sec:q3}{{10.2.3}{265}{An intracellular replacement for the standard model}{subsection.10.2.3}{}}
\newlabel{con:sec:q3@cref}{{[subsection][3][10,2]10.2.3}{[1][265][]265}}
\newlabel{con:sec:q4}{{10.2.4}{265}{Rates that move}{subsection.10.2.4}{}}
\newlabel{con:sec:q4@cref}{{[subsection][4][10,2]10.2.4}{[1][265][]265}}
\newlabel{con:sec:quadrant}{{10.3}{265}{The quadrant argument}{section.10.3}{}}
\newlabel{con:sec:quadrant@cref}{{[section][3][10]10.3}{[1][265][]265}}
\newlabel{con:sec:quadrant-ordering}{{10.3.1}{265}{Three environments, one ordering}{subsection.10.3.1}{}}
\newlabel{con:sec:quadrant-ordering@cref}{{[subsection][1][10,3]10.3.1}{[1][265][]265}}
\newlabel{con:eq:ordering}{{10.1}{266}{Three environments, one ordering}{equation.10.3.1}{}}
\newlabel{con:eq:ordering@cref}{{[equation][1][10]10.1}{[1][265][]266}}
\newlabel{con:eq:births}{{10.2}{266}{Three environments, one ordering}{equation.10.3.2}{}}
\newlabel{con:eq:births@cref}{{[equation][2][10]10.2}{[1][266][]266}}
\newlabel{con:eq:netordering}{{10.3}{266}{Three environments, one ordering}{equation.10.3.3}{}}
\newlabel{con:eq:netordering@cref}{{[equation][3][10]10.3}{[1][266][]266}}
\newlabel{con:rem:alternatives}{{10.1}{266}{What the argument does not consider}{theorem.10.1}{}}
\newlabel{con:rem:alternatives@cref}{{[remark][1][10]10.1}{[1][266][]266}}
\newlabel{con:sec:quadrant-threshold}{{10.3.2}{266}{Where the two strategies cross}{subsection.10.3.2}{}}
\newlabel{con:sec:quadrant-threshold@cref}{{[subsection][2][10,3]10.3.2}{[1][266][]266}}
\newlabel{con:eq:fractions}{{10.4}{266}{Where the two strategies cross}{equation.10.3.4}{}}
\newlabel{con:eq:fractions@cref}{{[equation][4][10]10.4}{[1][266][]266}}
\newlabel{con:eq:Rpro}{{10.5}{266}{Where the two strategies cross}{equation.10.3.5}{}}
\newlabel{con:eq:Rpro@cref}{{[equation][5][10]10.5}{[1][266][]266}}
\newlabel{con:eq:Ranti}{{10.6}{266}{Where the two strategies cross}{equation.10.3.6}{}}
\newlabel{con:eq:Ranti@cref}{{[equation][6][10]10.6}{[1][266][]266}}
\newlabel{con:eq:Rpro-linear}{{10.7}{266}{Where the two strategies cross}{equation.10.3.7}{}}
\newlabel{con:eq:Rpro-linear@cref}{{[equation][7][10]10.7}{[1][266][]266}}
\newlabel{con:eq:threshold-general}{{10.8}{267}{Where the two strategies cross}{equation.10.3.8}{}}
\newlabel{con:eq:threshold-general@cref}{{[equation][8][10]10.8}{[1][267][]267}}
\newlabel{con:eq:threshold}{{10.9}{267}{Where the two strategies cross}{equation.10.3.9}{}}
\newlabel{con:eq:threshold@cref}{{[equation][9][10]10.9}{[1][267][]267}}
\newlabel{con:fig:quadrant}{{10.1}{267}{The crossing. The expected net replication rate under a pro-phagocytic strategy falls linearly in the neutrophil fraction $p_N$, because a larger fraction of uptake events delivers the bacterium into the more hostile of the two cell types; the anti-phagocytic rate does not depend on $p_N$ at all. The single crossing is $p_N^{\ast }$ of \cref {con:eq:threshold}, and it lies strictly between $0$ and $1$ exactly when the ordering \cref {con:eq:ordering} holds}{figure.10.1}{}}
\newlabel{con:fig:quadrant@cref}{{[figure][1][10]10.1}{[1][267][]267}}
\newlabel{con:sec:quadrant-reading}{{10.3.3}{267}{What the threshold says}{subsection.10.3.3}{}}
\newlabel{con:sec:quadrant-reading@cref}{{[subsection][3][10,3]10.3.3}{[1][267][]267}}
\newlabel{con:sec:limits}{{10.4}{268}{Limitations and open problems}{section.10.4}{}}
\newlabel{con:sec:limits@cref}{{[section][4][10]10.4}{[1][268][]268}}
\newlabel{con:sec:closing}{{10.5}{269}{Closing}{section.10.5}{}}
\newlabel{con:sec:closing@cref}{{[section][5][10]10.5}{[1][269][]269}}
\newlabel{oi:details}{{5.5}{270}{Originality of the results}{section*.78}{}}
\newlabel{oi:details@cref}{{[table][5][5]5.5}{[1][270][]270}}
\newlabel{oi:statements}{{5.5}{270}{Originality of the results}{section*.79}{}}
\newlabel{oi:statements@cref}{{[table][5][5]5.5}{[1][270][]270}}
\newlabel{sup:ai-use}{{5.5}{273}{Statement on the use of generative AI}{section*.81}{}}
\newlabel{sup:ai-use@cref}{{[table][5][5]5.5}{[1][273][]273}}
\newlabel{sup:ai-account}{{5.5}{273}{Statement on the use of generative AI}{section*.82}{}}
\newlabel{sup:ai-account@cref}{{[table][5][5]5.5}{[1][273][]273}}
\makeatother
